



\chapter{Conically smooth stratified spaces}

We now turn to the theory of \emph{conically smooth stratified spaces}. These were first introduced in the work of Ayala--Francis--Tanaka and further elaborated in the subsequent works of Ayala--Francis--Rozenblyum \cite{AFT1, AFR1, AFR2}. They will form the geometric foundations for our stratified higher bordism categories, and their adoption for that role is undoubtedly one of the more novel constructions contained in this thesis. Part of this novelty lies in the recency of their introduction, and in the bespoke categorical toolkit required to master them well enough so as to use them effectively. 

It is a stated goal of this chapter to not only introduce the reader to the main definitions and results on conically smooth stratified spaces, but also expand, elaborate and even refine the exposition of these geometric objects. The overarching goal being to encourage their broader adoption in contexts where stratified spaces naturally arise.

\section{Stratified spaces}

When we introduce smooth manifolds formally, it is customary to first introduce the notion of topological space followed by the definition of topological manifold, after which smooth manifolds are defined as topological manifolds with an additional notion of smooth atlas. This is the trajectory we will emulate for conically smooth stratified spaces; the first section will introduce the equivalent to topological manifolds for stratified spaces.       

\begin{definition}\label{ch4:: def: upward closed}
    Let $P$ be a \textit{partially ordered set}, henceforth referred to by its customary shorthand poset. We say that a subset $U$ of $P$ is \textbf{upward closed} if for every $x \in U$, its upward closure, i.e., the set of elements $x\leq y$, is contained in $U$. 
\end{definition} 

\begin{observation}
Every poset is naturally endowed with the structure of category, where for every $p,q \in P$ there is a unique morphism $q \longra p$ if and only if $q\leq p$. Composition derives from the transitivity of $\leq$ while associativity is trivial.
\end{observation}

\begin{observation}
Given a poset $P$, there is an upward closed subset $P_{\geq p}$ associated to every $p \in P$ which is defined as 
\begin{equation}
    P_{\geq p} = \{ q \in P\,:\, p \leq q \}.
\end{equation}
Interpreting $P$ as a category, the subset $P_{\geq p}$, which is itself a sub-poset of $P$, can also be interpreted as the overcategory $p/P$.  
\end{observation}

\begin{definition}\label{ch4:: def: strat_topology}
    Let $P$ be a poset. We define a topology for $P$ for which an open subset $U \subset P$ is \emph{open} if it is upward closed. We will call this topology the Alexandroff topology. 
\end{definition}

\begin{example}
    The finite totally ordered set $[n]$ is a poset for every $n\geq -1$. The open subsets of $[n]$ are given by the image of the maps 
    \begin{equation}
    \langle n-k, n-k +1, \ldots, n \rangle \maps [k] \hookra [n]
    \end{equation} for every $0\leq k\leq n$. Alternatively, we can write these as $[n]_{\geq k}$.
\end{example}

\begin{observation}
We check that this is in fact a topology. Let $\{U_{i}\}_{i=1}^{n}$ be a finite list of open subsets of $P$, if $x \in U_1 \cap \cdots \cap U_n$ then if $y\geq x$ we have $y\in U_1,\cdots, U_n$ it thus belongs to the finite intersection.       
\end{observation}

A map $P \xra{f} Q$ of posets is continuous if and only if it is order-preserving. Indeed, the preimage of an upward closed subset under an order-preserving map is again upward closed, and conversely continuity forces $x\leq y$ to imply $f(x)\leq f(y)$. 

\begin{observation}\label{ch4: obs:: posets as tops}
    This means that there exists a fully faithful functor 
    \begin{equation}
        \mathsf{Posets} \hookra \Top
    \end{equation}
    from the category of posets with the Alexandroff topology and continuous maps thereamong, to the category of topological spaces.
\end{observation}

As noted by Ayala--Francis--Tanaka, the next definition first appears in Lurie's \emph{Higher Algebra} \cite[Appendix~A.5]{LurieHA}. 

\begin{definition}\label{ch4:: def: p-stratified space}
    Let $P$ be a poset, we define a $P$-stratified space to be a pair 
    \begin{equation}
        (X,S\maps X \longra P)
    \end{equation} where $X$ is a topological space and $S\maps X \longra P$ is a continuous map of topological spaces. 
\end{definition}

For $p\in P$, we write $X_p=S^{-1}(p)$ and call it the $p$-stratum of $X$.

\begin{observation}
There are many different definitions of stratified space in the literature, but all coalesce around the notion of $P$-stratified space. 
\end{observation}


\begin{observation}
    Ordinary topological spaces, or unstratified topological spaces, are encoded as stratified spaces with a trivial stratification, or a stratification by the trivial poset $P\simeq *$.
\end{observation}

As with any good definition of a new mathematical object, we must give a definition of stratified continuous map, which we will henceforth just call \emph{stratified map}. 
\begin{definition}\label{ch4:: def: stratified map}
    Let $(X, S\maps X \longra P)$ and $(Y, S'\maps Y \longra Q)$ be $P$-stratified and $Q$-stratified spaces, respectively. A \textbf{continuous stratified map} $f \maps (X,S\maps X \longra P ) \longra (Y,S'\maps Y \longra Q)$ is a pair $(f \maps X \longra Y, \alpha\maps P \longra Q)$ where $f$ and $\alpha$ are continuous maps such that: 
    \begin{equation}
        \begin{tikzcd}[ampersand replacement=\&]
            X \&\& Y \\
            \\
            P \&\& Q
            \arrow["f"{description}, from=1-1, to=1-3]
            \arrow["S"', from=1-1, to=3-1]
            \arrow["{S'}", from=1-3, to=3-3]
            \arrow["\alpha"{description}, from=3-1, to=3-3]
        \end{tikzcd}
    \end{equation}
    A pair of morphisms $(f,\alpha)$ and $(g,\beta)$ is composed by horizontally stacking the respective commutative diagrams. The composition is well-defined on the nose.
\end{definition}


\begin{remark}
    One compact and categorical way of thinking about how the theory of stratified topological spaces might be defined is as the following pullback category:
    \begin{equation}
        \mathsf{Ar}(\Top) \underset{\Top}{\times} \mathsf{Posets},
    \end{equation} 
    where the arrow $\mathsf{Ar}(\Top) \longra \Top$ sends a continuous map $X \longra Y$ to its target $Y$.
\end{remark}

\nid The next example indicates how $P$-stratifications connect to the typically adopted notion of stratification as a decomposition into locally closed subspaces. 

\begin{example}
    Let $X$ be a topological space, and let $P = \{a < b\}$ be a two-element poset. The Alexandroff topology is such that the only open subsets are $P$, $\{b\} \subset P$, and the $\emptyset$. A continuous map $S\maps X \longra P$ is specified entirely by the choice of a closed subset $X_0 \subseteq X$, corresponding to the pre-image of $a\in P$. Therefore, a $P$-stratified space in this case is equivalent to a topological space $X$ and a choice of a closed subset $X_0 \subseteq X$ — the continuous map $S$ sends points of $X_0$ to $a$ and points of $X\setminus X_0$ to $b$.
\end{example}

\begin{example}
    Posets do not need to be linearly ordered. As an example, let $P$ denote the \emph{span-shaped} poset
    \begin{equation}
        P = \big(l\, \la * \ra r \, \big) 
    \end{equation}  
    A $P$-stratified space $X$ consists of two open strata $X_l$ and $X_r$, together with a closed stratum $X_*$ contained in the complement of $X_l \cup X_r$. 
    Let us illustrate with a concrete example how a $P$-stratified torus looks like:
    \begin{equation}
    \tikzfig{figures/Strat/STR1}{\begin{tikzpicture}
	\begin{pgfonlayer}{nodelayer}
		\node [style=none] (0) at (-5, 0) {};
		\node [style=none] (1) at (5, 0) {};
		\node [style=none] (2) at (5, 0) {};
		\node [style=none] (3) at (-5, 0) {};
		\node [style=none] (4) at (-1.75, 0) {};
		\node [style=none] (5) at (1.75, 0) {};
		\node [style=none] (6) at (-1.5, -0.4) {};
		\node [style=none] (7) at (1.5, -0.4) {};
		\node [style=none] (8) at (0, -2.925) {};
		\node [style=none] (9) at (0, -0.75) {};
		\node [style=none] (10) at (0, 0.5) {};
		\node [style=none] (11) at (0, 2.925) {};
		\node [style=none] (12) at (0, 3.525) {$X_*$};
		\node [style=none] (13) at (0, -3.6) {$X_*$};
		\node [style=none] (14) at (-3.25, 0) {$X_{l}$};
		\node [style=none] (15) at (3.25, 0) {$X_{r}$};
	\end{pgfonlayer}
	\begin{pgfonlayer}{edgelayer}
		\draw [style={open_int_blue_2}, in=180, out=180, looseness=0.75] (10.center) to (11.center);
		\draw [style={open_int_blue_2_dash}, in=0, out=0, looseness=0.75] (11.center) to (10.center);
		\draw [style={open_int_blue_2}, in=180, out=180, looseness=0.75] (8.center) to (9.center);
		\draw [style={open_int_blue_2_dash}, in=0, out=0, looseness=0.75] (9.center) to (8.center);
		\draw [style=geometric, in=270, out=-90] (0.center) to (1.center);
		\draw [style=geometric, in=90, out=90] (2.center) to (3.center);
		\draw [style=geometric, in=270, out=-90, looseness=0.75] (4.center) to (5.center);
		\draw [style=geometric, in=90, out=90] (6.center) to (7.center);
	\end{pgfonlayer}
\end{tikzpicture}
}
    \end{equation}
    The two strata $X_{l}$ and $X_{r}$ are open cylinders while $X_*$ is the closed subspace which is contained in the complement of $X_{l} \cup X_{r} \subsetneq X$.
\end{example}

\begin{example}
To hone in on the intuition for what a stratified map is, we note the following pair of examples 
\begin{equation}
    \tikzfig{figures/Strat/STR10}{\begin{tikzpicture}
	\begin{pgfonlayer}{nodelayer}
		\node [style=0 strat] (4) at (-5, 0) {};
		\node [style=0 strat] (5) at (0, 0) {};
		\node [style=0 strat] (6) at (4.75, 0) {};
		\node [style=0 strat] (7) at (-5, 6.5) {};
		\node [style=0 strat] (8) at (0, 8.5) {};
		\node [style=0 strat] (9) at (0, 4.5) {};
		\node [style=0 strat] (10) at (5, 6.5) {};
		\node [style=none] (11) at (0, 3) {};
		\node [style=none] (12) at (0, 1.25) {};
		\node [style=none] (13) at (0, -1.5) {stratified map};
	\end{pgfonlayer}
	\begin{pgfonlayer}{edgelayer}
		\draw [style=geometric] (4) to (5);
		\draw [style=geometric] (5) to (6);
		\draw [style=geometric, in=180, out=-90, looseness=0.75] (7) to (9);
		\draw [style=geometric, in=-90, out=0, looseness=0.75] (9) to (10);
		\draw [style=geometric, in=90, out=0, looseness=0.75] (8) to (10);
		\draw [style=geometric, in=-180, out=90, looseness=0.75] (7) to (8);
		\draw [style=dir] (11.center) to (12.center);
	\end{pgfonlayer}
\end{tikzpicture}
} \qquad \qquad \tikzfig{figures/Strat/STR11}{\begin{tikzpicture}
	\begin{pgfonlayer}{nodelayer}
		\node [style=0 strat] (4) at (-5, 0) {};
		\node [style=0 strat] (5) at (0, 0) {};
		\node [style=0 strat] (6) at (4.75, 0) {};
		\node [style=0 strat] (7) at (-5, 6.5) {};
		\node [style=0 strat] (8) at (0, 8.5) {};
		\node [style=0 strat] (10) at (5, 6.5) {};
		\node [style=none] (11) at (0, 3) {};
		\node [style=none] (12) at (0, 1.25) {};
		\node [style=none] (13) at (0, 4.5) {};
		\node [style=none] (14) at (0, -1.5) {non-stratified map};
	\end{pgfonlayer}
	\begin{pgfonlayer}{edgelayer}
		\draw [style=geometric] (4) to (5);
		\draw [style=geometric] (5) to (6);
		\draw [style=geometric, in=90, out=0, looseness=0.75] (8) to (10);
		\draw [style=geometric, in=-180, out=90, looseness=0.75] (7) to (8);
		\draw [style=dir] (11.center) to (12.center);
		\draw [style=geometric, in=180, out=-90, looseness=0.75] (7) to (13.center);
		\draw [style=geometric, in=0, out=-90, looseness=0.75] (10) to (13.center);
	\end{pgfonlayer}
\end{tikzpicture}
}
\end{equation}
\smallskip

\noindent where on the left we depict as a source of the map a stratified circle which has eight different strata: it has four 0-dimensional strata (the black dots) and four 1-dimensional open arcs; the target is a stratified closed interval with five different strata: it has three 0-dimensional strata and two 1-dimensional strata given by the open arcs. The map is the vertical projection induced by the framing of the page. On the right, the spaces are the same but for the absence of a single 0-dimensional stratum. The rightmost map fails to be stratified because the lower arc of the source stratified circle gets smeared across different strata.     
\end{example}

\begin{example}
    Consider the two stratified circles of the example before. The underlying spaces are the same, so we can ask whether the identity homeomorphism is stratified or not. The answer is yes, but only in one direction — the maps 
    \begin{equation}
        \tikzfig{figures/Strat/STR12}{\begin{tikzpicture}
	\begin{pgfonlayer}{nodelayer}
		\node [style=0 strat] (7) at (-12.5, 6.5) {};
		\node [style=0 strat] (8) at (-7.5, 8.5) {};
		\node [style=0 strat] (9) at (-7.5, 4.5) {};
		\node [style=0 strat] (10) at (-2.5, 6.5) {};
		\node [style=0 strat] (14) at (2.5, 6.5) {};
		\node [style=0 strat] (15) at (7.5, 8.5) {};
		\node [style=0 strat] (16) at (12.5, 6.5) {};
		\node [style=none] (17) at (7.5, 4.5) {};
		\node [style=none] (18) at (-2, 7.5) {};
		\node [style=none] (19) at (-2, 5.5) {};
		\node [style=none] (20) at (2, 7.5) {};
		\node [style=none] (21) at (2, 5.5) {};
		\node [style=none] (22) at (0, 9) {stratified map};
		\node [style=none] (23) at (0, 4) {non-stratified map};
	\end{pgfonlayer}
	\begin{pgfonlayer}{edgelayer}
		\draw [style=geometric, in=180, out=-90, looseness=0.75] (7) to (9);
		\draw [style=geometric, in=-90, out=0, looseness=0.75] (9) to (10);
		\draw [style=geometric, in=90, out=0, looseness=0.75] (8) to (10);
		\draw [style=geometric, in=-180, out=90, looseness=0.75] (7) to (8);
		\draw [style=geometric, in=90, out=0, looseness=0.75] (15) to (16);
		\draw [style=geometric, in=-180, out=90, looseness=0.75] (14) to (15);
		\draw [style=geometric, in=180, out=-90, looseness=0.75] (14) to (17.center);
		\draw [style=geometric, in=0, out=-90, looseness=0.75] (16) to (17.center);
		\draw [style=dir, bend left] (18.center) to (20.center);
		\draw [style=dir, bend left] (21.center) to (19.center);
	\end{pgfonlayer}
\end{tikzpicture}
}
    \end{equation}
    are not both stratified. Only the top map is, and we call it a \textbf{refinement} precisely because it is witnessed by a homeomorphism of the underlying topological spaces.
\end{example}

\medskip
We have a working definition of \emph{topological} stratified space and continuous stratified map, now we would like to articulate what it means for a stratified space to be openly embedded into another stratified space. This is a fundamental touchstone in the path towards defining an appropriate site of topological stratified manifolds — which, as we have previously mentioned, will be the scaffolding for a theory of \emph{smooth} stratified spaces. 

Before introducing the definition of stratified open embedding, let us give the definition for the unstratified case first as a warm-up and a refresher.

\begin{definition}
    Let $X$ and $Y$ be topological spaces and $f\maps X \longra Y$ a continuous map. We say that $f$ is an \textbf{open embedding} if $f$ is \emph{open}, i.e., sends open subsets of $X$ to open subsets of $Y$, and an embedding i.e., the image $f(X) \subseteq Y$ is homeomorphic to $X$. Alternatively, we say that $f$ is a \textbf{closed embedding} if it is \emph{closed}, i.e., if it sends closed subsets of $X$ to closed subsets of $Y$ and also an embedding.
\end{definition}

\begin{definition}\label{ch4:: def: strat open emb}
    A continuous stratified map $(f,\alpha)\maps (X\longra P) \longra (Y \longra Q)$ is a \textbf{stratified open embedding} if:
    \begin{itemize}
        \item the underlying map $f\maps X \longra Y$ is an open embedding of topological spaces;
        \item for every $p \in P$, the restriction
        \begin{equation}
            f_{|X_{p}} \maps X_{p} \longra Y_{\alpha(p)}
        \end{equation}
        is an open embedding of topological manifolds.
    \end{itemize}
\end{definition}

Now we turn to two ways of producing stratified spaces from old ones, which will be especially important for our theory — \emph{cones} and \emph{products}. Together they will generate local models for our singular spaces.

\begin{definition}
    Let $(S \maps X \longra P)$ and $(S' \maps Y \longra Q)$ be a pair of stratified topological spaces. We define the following continuous map
    \begin{equation}
    \big(S\times S' \maps (X \times Y) \longra (P \times Q) \big),
    \end{equation}
    where $S\times S'$ is the product of the stratification maps. We call this the \textbf{product stratification}.
\end{definition}


\begin{definition}\label{ch4:: def: cones}
    Let $(X\longra P)$ be a stratified topological space. We define its associated \textbf{open cone} $\sC(X \longra P)$ to be the stratified topological space obtained as follows: 
    \begin{itemize}
        \item $\sC(X)$: the open cone on $X$ is the pushout of topological spaces
        \begin{equation}
            \sC(X) \, \overset{\sf def}{=} \,  \big(X\times \RR_{\geq 0}\big) \coprod_{X\times \{0\}} *.
        \end{equation}
        \item $\sC(P)$: the cone on $P$ is obtained by formally adjoining a minimal element $*$. Equivalently as the pushout of posets
        \begin{equation}
            \sC(P) \, \overset{\sf def}{=} \, (P \times [1]) \coprod_{P \times \{0\}} *.
        \end{equation}
    \end{itemize}
    The stratification $\sC(X) \longra \sC(P)$ sends $x \in X_{p} \times \RR_{> 0}$ to $p$ and sends the cone point $* \in \sC(X)$ to $*\in \sC(P)$. The cone point is thus sent to the cone point.
\end{definition}

\begin{remark}
    The poset $P \times [1]$ has two canonically isomorphic copies of $P$ embedded in it, call them $P_0$ and $P_1$, together with an arrow connecting every element of $P_0$ to its image in $P_1$. By taking the cone we are essentially collapsing $P_0$ and thus getting a new minimal element called $*$.  
\end{remark}

Since we are attempting to define a notion of topological stratified manifold, we must first isolate the correct topological class of stratified spaces.
\begin{definition}\cite[Definition~2.1.15]{AFT1}
    A \textbf{$C^0$ stratified space} is a paracompact stratified topological space belonging to the smallest full subcategory of stratified topological spaces such that:
    \begin{enumerate}
        \item it contains the empty stratified space;
        \item if $(X \ra P)$ is compact and belongs to it, then so does its open cone $\sC(X \ra P)$;
        \item if $(X \ra P)$ belongs to it, then so does $\RR \times (X \ra P)$, where $\RR$ carries the trivial stratification;
        \item it is closed under stratified open embeddings;
        \item if a stratified topological space admits an open cover by objects in the subcategory, then it also belongs to the subcategory.
    \end{enumerate}
\end{definition}

Since spaces of the form $\RR^i\times \sC(Z)$ are the local models singled out by this definition, we now isolate them.
\begin{definition}\label{ch4:: def: C0 basics}
    A \textbf{$C^0$ basic} is a $C^0$ stratified space of the form 
    \begin{equation}
        \RR^{i} \times \sC(Z)
    \end{equation}
    where $Z$ is a compact $C^0$ stratified space, $i\geq 0$ and $\RR^{i}$ is taken as the usual topological manifold with the trivial stratification. 
\end{definition}

\begin{lemma}\cite[Lemma~2.2.2]{AFT1}
    Let $X \longra P$ be a stratified topological space whose underlying space is second countable and Hausdorff. Consider the collection 
    \begin{equation}
        \{f_\alpha\maps U_\alpha \hookra X\mkern1mu | \, f_\alpha \text{ an open embedding} \}
    \end{equation}
    where $U_\alpha$ ranges over all $C^0$ basics. This collection forms a basis for the topology of $X$ if and only if $X$ is a $C^0$ stratified space.
\end{lemma}

\subsection{Dimension and depth}

When dealing with topological manifolds, there is a fundamental local invariant --- \emph{dimension}. For $C^0$ stratified spaces, this notion demands revisiting, since not every point admits a neighborhood homeomorphic to an open $\RR^{n}$ for some $n \geq 0$. 

\begin{definition}
    Let $(X \longra P)$ be a $C^0$ non-empty stratified space. For $x\in X$, we define its \textbf{local dimension} $\dim_x(X)$ to be the \emph{covering dimension} of $X$ at $x$, and we define the dimension of $(X\longra P)$ to be 
    \begin{equation}
        \mathsf{dim}(X) \mkern1mu \overset{\sf def}{=} \mkern1mu \underset{x \in X}{\mathsf{sup}} \mkern1mu\mkern1mu \mathsf{dim}_{x}(X).
    \end{equation}
    We say that $X$ has \textbf{pure dimension $n$} if it has constant local dimension, i.e., if $\mathsf{dim}_{x}(X) = n$ for all $x\in X$. We adopt the convention $\mathsf{dim}(\emptyset) = -1$.
\end{definition}

\begin{definition}
    Let $(X \longra P)$ be a non-empty $C^0$ stratified space and let $x\in X_{p}$ belong to the stratum indexed by $p\in P$. We define the \textbf{local depth} of $X$ at $x$ by
    \begin{equation}
        \depth_{x}(X) \mkern1mu \overset{\sf def}{=} \mkern1mu \mathsf{dim}_{x}(X) - \mathsf{dim}(X_{p}),
    \end{equation}
    and we define the \textbf{depth} of $X$ by
    \begin{equation}
        \depth(X) \mkern1mu \overset{\sf def}{=} \mkern1mu \underset{x \in X}{\mathsf{sup}} \mkern1mu \depth_x(X).
    \end{equation}
    We adopt the convention $\depth(\emptyset)=-1$.
\end{definition}

\begin{observation}
Before wrapping up the fundamentals of the theory of $C^0$-stratified spaces, let us end with the following constructions and results.
\end{observation}

\begin{construction}
    Let $\big(X \ra P \big)$ be a $C^0$ stratified space. There is a stratification $\big(X \ra P^{\sf con}\big)$ of $X$, where the $P^{\sf con}$ has as a set of objects the set of path components 
    \begin{equation}
        \pi_0\Big(\coprod_{p \in P} X_p\Big).
    \end{equation}
    Let $q$ and $q'$ denote a pair of connected strata in $X$, this receives a partial order from the relation 
    \begin{equation}
        q \leq q' \quad \Leftrightarrow \quad X_q \subset \overline{X}_{q'}.  
    \end{equation}
    The poset $P$ receives a natural surjective poset map from $P^{\sf con}$, defined by sending a connected component $q$ to its respective stratum. This map of posets also has the special property of being \emph{discrete} — the fiber contains solely identity arrows. There is a sense in which this factorization of $X \ra P$ 
    \begin{equation}
        X \longra P^{\sf con} \longra P
    \end{equation}
    is initial.
\end{construction}

\begin{remark}
    If a stratification $X \ra Q$ is such that all of its strata are connected, we call it a \define{characteristic stratification} or the continuous map itself \define{characteristic}.
\end{remark}

\begin{lemma}
    The category of factorizations of $\big(X \ra P\big)$ as a characteristic map followed by a discrete map of posets, is contractible.
\end{lemma}
\begin{proof}
  This is proved in \cite[Lemma B.37]{DD21}.   
\end{proof}
\medskip

\nid This in particular means that whether we work with a $P$-stratification or its associated stratification by connected components, is essentially the same — we will often use this fact to stipulate stratifications merely by highlighting their individual connected components, without a direct mention of the continuous map to a poset $P$, or even the poset $P$.

\section{Conical differentiability and smoothness}

Unburdened by the desire or the need to be succinct, we will now introduce a thorough discussion on conical smoothness. The reason for doing this is twofold: the first is premised on our appreciation for the subtle, yet dizzying, inductive nature of its definition; the second on our perception of their usefulness as a toolkit to the future studies of topological field theories — potentially decorated with any assortment of other adjectives e.g., with \emph{topological defects}.  

\subsection{Conically smooth along the cone locus}
One might argue that the single most important contribution of Ayala--Francis--Tanaka in their foundational paper on stratified spaces (see \cite{AFT1}), was their introduction of the notion of conical smoothness along cone loci — a mild regularity condition that calls for an adaptation of the usual derivative for smooth manifolds — defined by the existence of a limit --, without appealing to the secant $k$-planes that come with working within an ambient manifold, as is the case for Whitney stratifications. This adaptation chooses instead to efficiently work with what is left of a vector space around every point in a $C^{0}$ stratified space. 

To initiate this discussion, let $(Z\ra P)$ be a compact topological stratified space and let $U = \RR^{i} \times \sC(Z)$ be a topological basic. So as to easily distinguish points in $U$, let us fix coordinates
\begin{equation}
    \big(u,[r,z]\big) \in U
\end{equation}
where $\big( u,r,z \big)\in \RR^{i} \times \RR_{\geq 0} \times Z$, and $[s,z]$ is the image of $(s,z)$ under the quotient map $\pi_Z \maps \RR_{\geq 0} \times Z \longra \sC(Z)$ collapsing the zero fiber onto the cone point. 

\begin{definition}
    Let $U = \RR^{i} \times \sC(Z)$ be a topological basic. We define its \textbf{tangent cone} to be the space
    \begin{equation}
        \sT U \,\overset{\sf def}{=}\, \sT\RR^{i} \times \sC(Z).
    \end{equation}
\end{definition}

Since $\RR^{i}$ is framed, we can re-write the factors in the tangent cone as follows:
\begin{equation}
    \sT \RR^{i} \times \sC(Z) \simeq (\RR^{i}_{\vec{v}} \times \RR^{i}) \times \sC(Z) \simeq \RR^{i}_{\vec{v}} \times \big(\RR^{i} \times \sC(Z)\big) = \RR^{i}_{\vec{v}} \times U.
\end{equation} 
\begin{observation}
This identification suggests we can use coordinates 
\begin{equation}
    \big(\vec{v},u,[s,z]\big),\quad \text{ with } \vec{v} \in \RR^{i}_{\vec{v}},\big(u,[s,z]\big)\in U.
\end{equation}
Now we explicate the sense in which the tangent cone captures some of what is left of a vector space in a basic neighborhood. In a genuine real vector space — the case, if $Z=\emptyset$ — one is allowed to translate vectors along any direction and to scale vectors by any scalar amount $\lambda \in \RR$. We can re-package this observation formally into a statement about how the group
\begin{equation}
    \mathsf{Aff}^{+}(\RR^{i}) = \big(\RR^{i} \rtimes \RR_{>0}\big)
\end{equation} 
freely acts on $\RR^{i}$, with the action determined by 
\begin{align}
    \alpha \maps (\RR^{i} \rtimes \RR_{>0}) \times \RR^{i} &\longra \RR^{i} \\
    (\vec{v},\lambda,u) &\longmapsto (\vec{v} + \lambda u).
\end{align}
Given an arbitrary continuous map $f\maps \RR^{i} \longra \RR^{j}$ and a choice of some point $u\in \RR^{i}$, this action is instrumental in defining the derivative of $f$ at $u$:
\begin{align}
    \label{ch4:: eq: quotient derivative}
    \sD f\maps \sT\RR^{i} \simeq \RR^{i}_{\vec{v}} \times \RR^{i} &\longra \sT\RR^{j} \simeq \RR^{j}_{\vec{v}} \times \RR^{j} \\
    \label{ch4:: eq: quotient derivative 2}
    (\vec{v},u) &\longmapsto \underset{\lambda \ra 0^+}{\mathsf{lim}} \Bigg(\frac{~~f\big(\alpha_{\lambda,\vec{v}}(u)\big)-f(u)~~}{\lambda},\,\, f(u)\Bigg)
\end{align}
If the limit does not exist for $(\vec{v},u)$, we say the derivative of $f$ at $u$ does not exist. We can do better — the action map $\alpha_{\lambda,\vec{v}}$ admits an inverse induced by the inverse in the group
\begin{equation}
    \alpha_{\lambda,\vec{v}}^{-1}(u) = \alpha_{\frac{1}{\lambda},-\frac{\vec{v}}{\lambda}}(u) = \frac{1}{\lambda}(u-\vec{v})
\end{equation}
and the limit in \ref{ch4:: eq: quotient derivative 2} can be re-written as 
\begin{equation}
     \sD f(\vec{v},u) \, = \, \bigg(\underset{\lambda \ra 0^+}{\mathsf{lim}}\,\, \alpha_{\lambda,f(u)}^{-1}\circ f \circ \alpha_{\lambda,\vec{v}},\,\,f(u)\bigg).
\end{equation}
Thus the derivative may be viewed as the continuous extension, along the inclusion $\RR^{i}\times \RR_{>0} \hookra \RR^{i}\times \RR_{\geq 0}$, of the \emph{conjugated action} on $f$.

The basic neighborhood $U=\RR^{i}\times \sC(Z)$, assuming that $Z$ is not-empty retains an action by $\mathsf{Aff}^{+}(\RR^{i})$. This is the key geometric observation permitting an immediate bootstrap of the conventional definition of differentiability into a context where maps require preservation of strata. The action we are indicating here is as follows:
\begin{align}
    \gamma \maps  \RR_{>0} \times \sT \RR^{i} \times \sC(Z) &\longra \RR_{>0} \times \sT \RR^{i} \times \sC(Z) \\
    (\lambda,u,\vec{v},[s,z]) &\longmapsto (\lambda, \lambda\vec{v} + u, u, [\lambda s,z]). 
\end{align}
This action captures the infinitesimal displacements which are permissible in the vicinity of a point sitting at the origin of a basic, in partial analogy with how the affine group acts on the smooth tangent bundle. 
\end{observation}

We illustrate the action $\gamma$ with two examples, with the intent to consolidate the exposition.\footnote{In both cases, color is used for mere visual aid, and is not in any sense fundamental or intrinsic to the data of a stratified space. Apologies in advance for the readers who might otherwise not have the ability to distinguish all colors.}

\begin{example}\label{ch4:: ex: scaling and translation}
    In the first example, we consider the basic $U=\RR \times \sC(S^{0}) \simeq \RR \times \sC(*_+ \sqcup *_{-})$ which we choose to represent as follows
    \begin{equation}
        \tikzfig{figures/Strat/STR2}{\begin{tikzpicture}
	\begin{pgfonlayer}{nodelayer}
		\node [style=none] (0) at (-2, 0) {};
		\node [style=none] (1) at (2, 0) {};
		\node [style=none] (2) at (-2, 0) {};
		\node [style=none] (3) at (2, 0) {};
		\node [style=none] (4) at (0, 2.05) {};
		\node [style=none] (5) at (0, -2.05) {};
		\node [style=none] (6) at (-2, 0) {};
		\node [style=none] (7) at (0, -2.05) {};
		\node [style=none] (8) at (0, 2.05) {};
		\node [style=none] (9) at (2, 0) {};
	\end{pgfonlayer}
	\begin{pgfonlayer}{edgelayer}
		\draw [style=ts-dashed-grey, in=90, out=90, looseness=1.75] (0.center) to (1.center);
		\draw [style=ts-dashed-grey, in=-90, out=-90, looseness=1.75] (2.center) to (3.center);
		\draw [style=defect patch 1] (6.center)
			 to [in=-180, out=90] (4.center)
			 to (5.center)
			 to [in=-90, out=180] cycle;
		\draw [style=defect patch 2] (9.center)
			 to [in=0, out=-90] (7.center)
			 to (8.center)
			 to [in=90, out=0] cycle;
		\draw [style={open_int_green}] (8.center) to (7.center);
	\end{pgfonlayer}
\end{tikzpicture}
} \quad \simeq \quad \RR \,\, \times \,\, \tikzfig{figures/Strat/STR3}{\begin{tikzpicture}
	\begin{pgfonlayer}{nodelayer}
		\node [style=none] (0) at (-3, 0) {};
		\node [style=none] (1) at (3, 0) {};
		\node [style=none] (2) at (0, 0) {};
		\node [style={endpoint_yellow_small}] (3) at (0, 0) {};
	\end{pgfonlayer}
	\begin{pgfonlayer}{edgelayer}
		\draw [style=defect line 1] (0.center) to (3.center);
		\draw [style=defect line 2] (3.center) to (1.center);
	\end{pgfonlayer}
\end{tikzpicture}
}
    \end{equation}
    with a choice of identification $*_+$ with orange (the left patch) and $*_{-}$ with pink (the right patch). The \emph{cone loci} is identified with the green strand. Note of course that we implicitly \emph{rounded} the corners to obtain a disk. We can attempt to visualize the action $\gamma$ of $\RR^{i} \rtimes \RR_{>0}$: 
    \begin{enumerate}
        \item The action of scaling $\gamma_{\lambda, 0}$ along the radial direction can be depicted as follows
        \begin{equation}
         \tikzfig{figures/Strat/STR2}{\begin{tikzpicture}
	\begin{pgfonlayer}{nodelayer}
		\node [style=none] (0) at (-2, 0) {};
		\node [style=none] (1) at (2, 0) {};
		\node [style=none] (2) at (-2, 0) {};
		\node [style=none] (3) at (2, 0) {};
		\node [style=none] (4) at (0, 2.05) {};
		\node [style=none] (5) at (0, -2.05) {};
		\node [style=none] (6) at (-2, 0) {};
		\node [style=none] (7) at (0, -2.05) {};
		\node [style=none] (8) at (0, 2.05) {};
		\node [style=none] (9) at (2, 0) {};
	\end{pgfonlayer}
	\begin{pgfonlayer}{edgelayer}
		\draw [style=ts-dashed-grey, in=90, out=90, looseness=1.75] (0.center) to (1.center);
		\draw [style=ts-dashed-grey, in=-90, out=-90, looseness=1.75] (2.center) to (3.center);
		\draw [style=defect patch 1] (6.center)
			 to [in=-180, out=90] (4.center)
			 to (5.center)
			 to [in=-90, out=180] cycle;
		\draw [style=defect patch 2] (9.center)
			 to [in=0, out=-90] (7.center)
			 to (8.center)
			 to [in=90, out=0] cycle;
		\draw [style={open_int_green}] (8.center) to (7.center);
	\end{pgfonlayer}
\end{tikzpicture}
}\,\,\,\, \xra{\,\,\,\,\gamma_{\lambda,0}\,\,\,\,}\,\,\,\, \tikzfig{figures/Strat/STR4}{\begin{tikzpicture}
	\begin{pgfonlayer}{nodelayer}
		\node [style=none] (7) at (0, -2.05) {};
		\node [style=none] (8) at (0, 2.05) {};
		\node [style=none] (9) at (2, 0) {};
		\node [style=none] (0) at (-2, 0) {};
		\node [style=none] (1) at (2, 0) {};
		\node [style=none] (2) at (-2, 0) {};
		\node [style=none] (3) at (2, 0) {};
		\node [style=none] (4) at (0, 2.05) {};
		\node [style=none] (5) at (0, -2.05) {};
		\node [style=none] (6) at (-2, 0) {};
		\node [style=none] (10) at (-1.25, 0) {};
		\node [style=none] (11) at (1.25, 0) {};
		\node [style=none] (12) at (-1.25, 0) {};
		\node [style=none] (13) at (1.25, 0) {};
		\node [style=none] (14) at (0, 1.3) {};
		\node [style=none] (15) at (0, -1.3) {};
		\node [style=none] (16) at (-1.25, 0) {};
		\node [style=none] (17) at (0, -1.3) {};
		\node [style=none] (18) at (0, 1.3) {};
		\node [style=none] (19) at (1.25, 0) {};
	\end{pgfonlayer}
	\begin{pgfonlayer}{edgelayer}
		\draw [style=defect patch 2] (9.center)
			 to [in=0, out=-90] (7.center)
			 to (8.center)
			 to [in=90, out=0] cycle;
		\draw [style=ts-dashed-grey, in=90, out=90, looseness=1.75] (0.center) to (1.center);
		\draw [style=ts-dashed-grey, in=-90, out=-90, looseness=1.75] (2.center) to (3.center);
		\draw [style=defect patch 1] (6.center)
			 to [in=-180, out=90] (4.center)
			 to (5.center)
			 to [in=-90, out=180] cycle;
		\draw [style=ts-dashed, in=90, out=90, looseness=1.75] (10.center) to (11.center);
		\draw [style=ts-dashed, in=-90, out=-90, looseness=1.75] (12.center) to (13.center);
		\draw [style={open_int_green}] (18.center) to (17.center);
		\draw [style=defect patch 1 dark] (16.center)
			 to [in=-180, out=90] (14.center)
			 to (15.center)
			 to [in=-90, out=180] cycle;
		\draw [style=defect patch 2 dark] (19.center)
			 to [in=0, out=-90] (17.center)
			 to (18.center)
			 to [in=90, out=0] cycle;
		\draw [style={open_int_green}] (8.center) to (7.center);
	\end{pgfonlayer}
\end{tikzpicture}
} 
        \end{equation}
        where we have chosen $0< \lambda < 1$. Note also that this map is just the identity on the $\RR$ i.e., the Euclidean coordinate. 
        \item The action of translation $\gamma_{1,\vec{v}}$ along the cone locus can be depicted as follows
        \begin{equation}
            \tikzfig{figures/Strat/STR6}{\begin{tikzpicture}
	\begin{pgfonlayer}{nodelayer}
		\node [style=none] (0) at (-2, 2) {};
		\node [style=none] (1) at (2, 2) {};
		\node [style=none] (2) at (-2, -2) {};
		\node [style=none] (3) at (2, -2) {};
		\node [style=none] (4) at (0, 2) {};
		\node [style=none] (5) at (0, -2) {};
		\node [style=none] (6) at (0, 2) {};
		\node [style=none] (7) at (0, -2) {};
		\node [style=none] (8) at (2, 2) {};
		\node [style=none] (9) at (2, -2) {};
		\node [style=none] (10) at (0, 0) {};
		\node [style=none] (11) at (0, 1) {};
		\node [style=none] (15) at (1.5, -1) {};
		\node [style=none] (17) at (0, -1) {};
		\node [style=none] (23) at (0, 2) {};
		\node [style=none] (24) at (0.5, 0.5) {$\scriptstyle{\vec{v}}$};
	\end{pgfonlayer}
	\begin{pgfonlayer}{edgelayer}
		\draw [style=defect patch 1] (0.center)
			 to (2.center)
			 to cycle
			 to (4.center)
			 to (5.center)
			 to (2.center);
		\draw [style=ts-dashed-grey] (2.center) to (3.center);
		\draw [style=ts-dashed-grey] (3.center) to (1.center);
		\draw [style=ts-dashed-grey] (1.center) to (0.center);
		\draw [style=defect patch 2] (6.center)
			 to (7.center)
			 to cycle
			 to (8.center)
			 to (9.center)
			 to (7.center);
		\draw [style={open_int_green}, in=90, out=-90] (6.center) to (7.center);
		\draw [style=dir thick, in=270, out=90] (10.center) to (11.center);
	\end{pgfonlayer}
\end{tikzpicture}
}\,\,\,\, \xra{\,\,\,\,\gamma_{1,\vec{v}}\,\,\,\,} \,\,\,\, \tikzfig{figures/Strat/STR5}{\begin{tikzpicture}
	\begin{pgfonlayer}{nodelayer}
		\node [style=none] (0) at (-2, 2) {};
		\node [style=none] (1) at (2, 2) {};
		\node [style=none] (2) at (-2, -2) {};
		\node [style=none] (3) at (2, -2) {};
		\node [style=none] (4) at (0, 2) {};
		\node [style=none] (5) at (0, -2) {};
		\node [style=none] (6) at (0, 2) {};
		\node [style=none] (7) at (0, -2) {};
		\node [style=none] (8) at (2, 2) {};
		\node [style=none] (9) at (2, -2) {};
		\node [style=none] (10) at (0, 0) {};
		\node [style=none] (11) at (0, 1) {};
		\node [style=none] (12) at (-1.5, 3) {};
		\node [style=none] (13) at (1.5, 3) {};
		\node [style=none] (14) at (-1.5, -1) {};
		\node [style=none] (15) at (1.5, -1) {};
		\node [style=none] (16) at (0, 3) {};
		\node [style=none] (17) at (0, -1) {};
		\node [style=none] (23) at (0, 2) {};
		\node [style=none] (18) at (0, 3) {};
		\node [style=none] (19) at (0, -1) {};
		\node [style=none] (20) at (1.5, 3) {};
		\node [style=none] (21) at (1.5, -1) {};
	\end{pgfonlayer}
	\begin{pgfonlayer}{edgelayer}
		\draw [style=defect patch 1] (0.center)
			 to (2.center)
			 to cycle
			 to (4.center)
			 to (5.center)
			 to (2.center);
		\draw [style=ts-dashed-grey] (2.center) to (3.center);
		\draw [style=ts-dashed-grey] (3.center) to (1.center);
		\draw [style=ts-dashed-grey] (1.center) to (0.center);
		\draw [style=defect patch 2] (6.center)
			 to (7.center)
			 to cycle
			 to (8.center)
			 to (9.center)
			 to (7.center);
		\draw [style=defect patch 1 dark] (12.center)
			 to (14.center)
			 to cycle
			 to (16.center)
			 to (17.center)
			 to (14.center);
		\draw [style=ts-dashed-grey] (14.center) to (15.center);
		\draw [style=ts-dashed-grey] (15.center) to (13.center);
		\draw [style=ts-dashed-grey] (13.center) to (12.center);
		\draw [style={open_int_green}] (18.center) to (19.center);
		\draw [style=ts-dashed-grey] (12.center) to (14.center);
		\draw [style=defect patch 2 dark] (18.center)
			 to (19.center)
			 to cycle
			 to (20.center)
			 to [in=90, out=-90] (21.center)
			 to (19.center);
		\draw [style={open_int_green}, in=90, out=-90] (6.center) to (7.center);
		\draw [style=dir thick, in=270, out=90] (10.center) to (11.center);
	\end{pgfonlayer}
\end{tikzpicture}
} 
        \end{equation}
        where we have made a light rescaling to better indicate the image of the mapping.
    \end{enumerate}
    Noticeably, both of these maps are stratified and homeomorphisms. Note also that, despite the fact that the underlying space is homeomorphic to $\RR^2$, we dispense of any translation along a direction not in the cone locus — as that action would budge the strata \emph{off} the cone locus, and circumvent any interpretation as a conically smooth infinitesimal action. Finally, observe that neither action interacts with the $z$ coordinate. That is because we cannot scale or translate points of an arbitrary compact topological stratified space. 
\end{example}


\begin{example}
    As a second example, let us consider the open cone on a circle stratified by its north and south poles $U = \sC(\SS^{1})$ where $\SS^1 = (q \maps S^{1} \ra [1])$ and $q$ sends ${N,S}$ to $\{0\}$ and the remaining two open arcs to $\{1\}$, which we depict as follows
    \begin{equation}
        \tikzfig{figures/Strat/STR7}{\begin{tikzpicture}
	\begin{pgfonlayer}{nodelayer}
		\node [style=none] (0) at (-2, 0) {};
		\node [style=none] (1) at (2, 0) {};
		\node [style=none] (2) at (-2, 0) {};
		\node [style=none] (3) at (2, 0) {};
		\node [style=none] (4) at (0, 2.05) {};
		\node [style=none] (5) at (0, -2.05) {};
		\node [style=none] (6) at (-2, 0) {};
		\node [style=none] (7) at (0, -2.05) {};
		\node [style=none] (8) at (0, 2.05) {};
		\node [style=none] (9) at (2, 0) {};
		\node [style=endpoint] (10) at (0, 0) {};
	\end{pgfonlayer}
	\begin{pgfonlayer}{edgelayer}
		\draw [style=ts-dashed-grey, in=90, out=90, looseness=1.75] (0.center) to (1.center);
		\draw [style=ts-dashed-grey, in=-90, out=-90, looseness=1.75] (2.center) to (3.center);
		\draw [style=defect patch 1] (6.center)
			 to [in=-180, out=90] (4.center)
			 to (5.center)
			 to [in=-90, out=180] cycle;
		\draw [style=defect patch 1] (9.center)
			 to [in=0, out=-90] (7.center)
			 to (8.center)
			 to [in=90, out=0] cycle;
		\draw [style={open_int_green}] (8.center) to (7.center);
	\end{pgfonlayer}
\end{tikzpicture}
} \,\,\,\, \simeq \,\,\,\, \text{open cone on } \,\,\tikzfig{figures/Strat/STR8}{\begin{tikzpicture}
	\begin{pgfonlayer}{nodelayer}
		\node [style=none] (11) at (-2, 0) {};
		\node [style=none] (14) at (2, 0) {};
		\node [style=none] (15) at (-2, 0) {};
		\node [style=none] (16) at (2, 0) {};
		\node [style={endpoint_yellow_small}] (17) at (0, 2.05) {};
		\node [style={endpoint_yellow_small}] (18) at (0, -2.05) {};
	\end{pgfonlayer}
	\begin{pgfonlayer}{edgelayer}
		\draw [style=defect line 1, bend left=90, looseness=1.75] (11.center) to (14.center);
		\draw [style=defect line 1, bend right=90, looseness=1.75] (15.center) to (16.center);
	\end{pgfonlayer}
\end{tikzpicture}
}.
    \end{equation}
    Despite diverging from \Cref{ch4:: ex: scaling and translation} only in the fact that it contains a point stratum at the origin, the implication is that it no longer admits any non-trivial action by translation. The only action is by scaling along the radial coordinate.
    \begin{equation}
        \tikzfig{figures/Strat/STR7}{\begin{tikzpicture}
	\begin{pgfonlayer}{nodelayer}
		\node [style=none] (0) at (-2, 0) {};
		\node [style=none] (1) at (2, 0) {};
		\node [style=none] (2) at (-2, 0) {};
		\node [style=none] (3) at (2, 0) {};
		\node [style=none] (4) at (0, 2.05) {};
		\node [style=none] (5) at (0, -2.05) {};
		\node [style=none] (6) at (-2, 0) {};
		\node [style=none] (7) at (0, -2.05) {};
		\node [style=none] (8) at (0, 2.05) {};
		\node [style=none] (9) at (2, 0) {};
		\node [style=endpoint] (10) at (0, 0) {};
	\end{pgfonlayer}
	\begin{pgfonlayer}{edgelayer}
		\draw [style=ts-dashed-grey, in=90, out=90, looseness=1.75] (0.center) to (1.center);
		\draw [style=ts-dashed-grey, in=-90, out=-90, looseness=1.75] (2.center) to (3.center);
		\draw [style=defect patch 1] (6.center)
			 to [in=-180, out=90] (4.center)
			 to (5.center)
			 to [in=-90, out=180] cycle;
		\draw [style=defect patch 1] (9.center)
			 to [in=0, out=-90] (7.center)
			 to (8.center)
			 to [in=90, out=0] cycle;
		\draw [style={open_int_green}] (8.center) to (7.center);
	\end{pgfonlayer}
\end{tikzpicture}
}\,\,\,\, \xra{\,\,\,\,\gamma_{\lambda,0}\,\,\,\,}\,\,\,\, \tikzfig{figures/Strat/STR9}{\begin{tikzpicture}
	\begin{pgfonlayer}{nodelayer}
		\node [style=none] (0) at (-2, 0) {};
		\node [style=none] (1) at (2, 0) {};
		\node [style=none] (2) at (-2, 0) {};
		\node [style=none] (3) at (2, 0) {};
		\node [style=none] (4) at (0, 2.05) {};
		\node [style=none] (5) at (0, -2.05) {};
		\node [style=none] (6) at (-2, 0) {};
		\node [style=none] (7) at (0, -2.05) {};
		\node [style=none] (8) at (0, 2.05) {};
		\node [style=none] (9) at (2, 0) {};
		\node [style=none] (10) at (-1.25, 0) {};
		\node [style=none] (11) at (1.25, 0) {};
		\node [style=none] (12) at (-1.25, 0) {};
		\node [style=none] (13) at (1.25, 0) {};
		\node [style=none] (14) at (0, 1.3) {};
		\node [style=none] (15) at (0, -1.3) {};
		\node [style=none] (16) at (-1.25, 0) {};
		\node [style=none] (17) at (0, -1.3) {};
		\node [style=none] (18) at (0, 1.3) {};
		\node [style=none] (19) at (1.25, 0) {};
		\node [style=endpoint] (20) at (0, 0) {};
	\end{pgfonlayer}
	\begin{pgfonlayer}{edgelayer}
		\draw [style=ts-dashed-grey, in=90, out=90, looseness=1.75] (0.center) to (1.center);
		\draw [style=ts-dashed-grey, in=-90, out=-90, looseness=1.75] (2.center) to (3.center);
		\draw [style=defect patch 1] (6.center)
			 to [in=-180, out=90] (4.center)
			 to (5.center)
			 to [in=-90, out=180] cycle;
		\draw [style=defect patch 1] (9.center)
			 to [in=0, out=-90] (7.center)
			 to (8.center)
			 to [in=90, out=0] cycle;
		\draw [style={open_int_green}] (8.center) to (7.center);
		\draw [style=ts-dashed-grey, in=90, out=90, looseness=1.75] (10.center) to (11.center);
		\draw [style=ts-dashed-grey, in=-90, out=-90, looseness=1.75] (12.center) to (13.center);
		\draw [style=defect patch 1 dark] (16.center)
			 to [in=-180, out=90] (14.center)
			 to (15.center)
			 to [in=-90, out=180] cycle;
		\draw [style=defect patch 1 dark] (19.center)
			 to [in=0, out=-90] (17.center)
			 to (18.center)
			 to [in=90, out=0] cycle;
		\draw [style={open_int_green}] (18.center) to (17.center);
	\end{pgfonlayer}
\end{tikzpicture}
} 
    \end{equation}
    The cone locus has changed from being 1- to 0-dimensional. This severely restricts the kind of infinitesimal automorphisms allowed, and as described in \Cref{ch4:: obs: con smooth circles} it has strong implications for the comparison between the ordinary notion of $C^{1}$ and the weaker $C^{1}$ along the cone locus.
\end{example}

Immediately, we are led to consider a notion of differentiability for basics modeled on the continuous extension of the conjugated action. Let us define this explicitly.


\begin{definition}\cite[Definition 3.1.4]{AFT1}
    \label{ch4:: def: conically smooth along}
    Given a pair of compact stratified topological spaces $\big(Z\ra P\big)$ and $\big(W\ra Q\big)$ and a continuous stratified map 
    \begin{equation}
        f \maps \RR^{i} \times \sC(Z) \longra \RR^{j} \times \sC(W),
    \end{equation} 
    we say that \textbf{$f$ is $C^1$ along $\RR^{i}$} if\footnote{Often this is instead stated as \emph{along the cone locus}, since $\RR^{i}\times \{*\}$ contains all the points in a basic which are in the deepest stratum.}:
    \begin{itemize}
        \item \textbf{(Cone):} The underlying map of posets $\sC(P) \longra \sC(Q)$ maps the cone point $*_P$ to the cone point $*_Q$.
        \item \textbf{($C^1$ along):} If the solid diagram 
        \begin{equation}
            \begin{tikzcd}
                {\RR_{>0}\times \sT\RR^{i}\times \sC(Z)} &&& {\RR_{>0}\times \sT\RR^{j}\times \sC(W)} \\
                \\
                {\RR_{\geq 0}\times \sT\RR^{i}\times \sC(Z)} &&& {\RR_{\geq 0}\times \sT\RR^{j}\times \sC(W)}
                \arrow["{\gamma^{-1}\circ f_{\Delta}\circ \gamma}", from=1-1, to=1-4]
                \arrow[hook, from=1-1, to=3-1]
                \arrow[hook, from=1-4, to=3-4]
                \arrow["{\sD f}"', dashed, from=3-1, to=3-4]
            \end{tikzcd}
        \end{equation}
        admits a continuous extension $\sD f$.  
    \end{itemize}
    For $r>1$, we say that $f$ is \textbf{$C^{r}$ along $\RR^{i}$} if it is $C^1$ along $\RR^{i}$ and, under the identification
    \begin{equation}
        \sT\RR^{i}\times \sC(Z) \simeq \RR^{i}\times \RR^{i}\times \sC(Z),
    \end{equation}
    the derivative
    \begin{equation}
        \sD f \maps \RR^{i}\times \RR^{i}\times \sC(Z) \longra \RR^{j}\times \RR^{j}\times \sC(W)
    \end{equation}
    is $C^{r-1}$ along $\RR^{i}\times \RR^{i}$. We say that $f$ is \textbf{conically smooth along $\RR^i$} if it is $C^r$ along $\RR^i$ for every $r\geq 1$.
\end{definition}

\begin{observation}\label{ch4:: obs: con smooth circles}
Conical smoothness along $\RR^{i}$ for a given basic, is a weaker notion than that of smoothness, even in cases where the cone point admits a neighborhood homeomorphic to an open disk. To see this, we observe the example outlined in \cite[Example 3.1.9]{AFT1}: choosing polar coordinates $\sC(S^{n-1})\cong \RR^{n}$ we note that any arbitrary homeomorphism $f\maps S^{n-1} \ra S^{n-1}$ will induce a conically smooth map (but generally not smooth) $\sC(f)\maps \sC(S^{n-1}) \ra \sC(S^{n-1})$ at the cone point $*$. 

\medskip

\nid This is sensible, since all we are demanding is that the \emph{radial directional derivatives} exist and are continuous — we are taking the limit of the conjugation of $f$ exclusively by the radial action, since no other is available. Yet, this is not enough to ensure $\sC(f)$ is $C^1$ at the origin — this example is an indicator of the type of regularity we will work with as we venture into conically smooth stratified geometry.   
\end{observation}



\subsection{Stratified maps}

The full inductive construction of the category of basics and the sheaf of maximal atlases in \cite[Section~3.2]{AFT1} is somewhat elaborate. Since our aim is to proceed to the geometric applications, we record only the streamlined formulation that we shall use: first the recursive definition on basics, and then the chartwise definitions for atlases and arbitrary maps.


\begin{definition}[Conically smooth maps between basics]\cite[Definition~3.3.1]{AFT1}
    \label{ch4:: def: conically smooth map}
    We declare the unique map $\emptyset \to \emptyset$ to be conically smooth. Assume inductively that conically smooth maps have already been defined for basics of strictly smaller depth. Let
    \begin{equation}
        f \maps \RR^{i}\times \sC(Z) \longra \RR^{j}\times \sC(Y)
    \end{equation}
    be a continuous stratified map between basics. Then $f$ is \define{conically smooth} if:
    \begin{enumerate}
        \item \textbf{Hits the cone-point:} if, on the level of posets, $f$ sends the cone-point to the cone-point, then
        \begin{itemize}
            \item $f$ is conically smooth along $\RR^{i}$ in the sense of \Cref{ch4:: def: conically smooth along};
            \item the restriction
            \begin{equation}
                f_{|f^{-1}(\RR^{j}\times Y \times \RR_{>0})} \maps f^{-1}(\RR^{j}\times Y \times \RR_{>0}) \longra \RR^{j}\times Y \times \RR_{>0}
            \end{equation}
            is conically smooth.
        \end{itemize}
        \item \textbf{Misses the cone-point:} if, on the level of posets, $f$ misses the cone-point, then it factors as
        \begin{equation}
            \begin{tikzcd}[ampersand replacement=\&]
                {\RR^{i}\times \sC(Z)} \&\& {\RR^{j}\times \sC(Y)} \\
                \\
                \& {\RR^{j}\times Y \times \RR_{>0}} \& {.}
                \arrow["f", from=1-1, to=1-3]
                \arrow["{f_0}"', dashed, from=1-1, to=3-2]
                \arrow[hook, from=3-2, to=1-3]
            \end{tikzcd}
        \end{equation}
        and we say that $f$ is conically smooth if $f_0$ is conically smooth.
    \end{enumerate}
\end{definition}

\begin{definition}\cite[Definition~3.2.10]{AFT1}
    A \textbf{conically smooth atlas} on a $C^0$ stratified space $M$ is a collection of charts
    \begin{equation}
        \big\{\varphi_\alpha \maps U_\alpha \hookra M\big\}_\alpha
    \end{equation}
    where each $U_\alpha$ is a basic, the images $\{\varphi_\alpha(U_\alpha)\}_\alpha$ cover $M$, and whenever
    \begin{equation}
        x \in \varphi_\alpha(U_\alpha)\cap \varphi_\beta(U_\beta),
    \end{equation}
    there exists a basic $W$ and conically smooth open embeddings
    \begin{equation}
        W \hookra U_\alpha, \qquad W \hookra U_\beta
    \end{equation}
    whose composites with $\varphi_\alpha$ and $\varphi_\beta$ agree on their images, and this common image contains $x$. Here, a \textbf{conically smooth open embedding between basics} means a conically smooth map between basics whose underlying stratified map is a stratified open embedding. A \textbf{conically smooth stratified space} is a $C^0$ stratified space equipped with a maximal conically smooth atlas.
\end{definition}

\begin{definition}
    Let $f \maps \big(M,\mathcal{A}\big) \longra \big(N,\mathcal{B}\big)$ be a continuous stratified map between conically smooth stratified spaces. We say that $f$ is \define{conically smooth} if for every pair of charts
    \begin{equation}
        (\varphi\maps U \hookra M) \in \mathcal{A}, \qquad (\psi \maps V \hookra N) \in \mathcal{B}
    \end{equation}
    such that
    \begin{equation}
       f\big(\varphi(U)\big) \subset \psi(V),
    \end{equation}
    the induced map between basics
    \begin{equation}
        \psi^{-1} \circ f \circ \varphi \maps U \longra V
    \end{equation}
    is conically smooth in the sense of \Cref{ch4:: def: conically smooth map}.
\end{definition}

\begin{proposition}
    If $M\xra{\,\, f \,\,} N \xra{\,\, g \,\,} Q$ is a pair of composable conically smooth maps, then the composite is also a conically smooth map. 
\end{proposition}
\begin{proof}
    It suffices to check this result again at the level of basics.
\end{proof}


\begin{observation}
    Let $f \maps M \longra N$ be a conically smooth map, and let $N$ be trivially stratified, i.e., a smooth manifold. Per \Cref{ch4:: def: conically smooth map} we have that since, locally, $f$ (as a map of posets) preserves the cone-point $ U = \RR^{i} \times \sC(Z) \longra \RR^{j} = V$ for some basics $U$ of $M$ and $V$ of $N$, it suffices to check that $f$ is conically smooth along $\RR^{i}$.
\end{observation}

\begin{observation}
    A map of smooth manifolds, equivalently of depth $0$ conically smooth stratified spaces regarded as trivially stratified, is conically smooth. 
\end{observation}

\begin{definition}
    Let $\strat$ denote the ordinary category whose objects are conically smooth stratified spaces and whose morphisms are conically smooth maps. Let $\snglr \subset \strat$ denote the wide subcategory whose morphisms are the conically smooth maps whose underlying stratified maps are stratified open embeddings. Let $\bsc \subset \snglr$ denote the full subcategory spanned by basics.
\end{definition}

\begin{proposition}\cite[Corollary~3.4.10]{AFT1}
    The category $\strat$ admits products.
\end{proposition}

Finally, we note the following useful result.

\begin{lemma}\cite[Lemma 3.4.5]{AFT1}
    Let $\big(M \ra P, \mathcal{A}\big)$ be a conically smooth stratified space, and let $Q \hookra P$ be a consecutive inclusion of posets. The pullback $C^0$ stratified space 
    \begin{equation}
        \label{eq 2}
        \big(M_{|Q} \ra Q\big) \longra (M \ra P)
    \end{equation}
    inherits a conically smooth structure such that the map \Cref{eq 2} is also conically smooth.
\end{lemma}

\begin{cor}
    Let $\big(M \ra P,\mathcal{A}\big)$ be a conically smooth stratified space. Then we have that for every $p\in P$:
    \begin{enumerate}
        \item $M_{\leq p}$ is a conically smooth stratified subspace.
        \item $M_{p}$ is a smooth manifold.
        \item $M_{\nless p}$ is a conically smooth stratified subspace.
    \end{enumerate}
\end{cor}
\begin{proof}
    Note that for $(1)$ we consider $P^{\leq p} \hookra P$; which is consecutive because: it is a full subcategory of $P$ (spanned by all elements $x\leq p$.), and because given a pair $(x\leq p), (y\leq p) \in P_{\leq p}$, we have that every $z\in P$ such that $(x \leq z \leq y)$ also satisfies $ z\leq p$. For $(2)$ we have that $\{p\} \hookra P$ is trivially consecutive. Lastly, $(3)$ is consecutive because it is a full subcategory — it is the full subcategory spanned by all the objects $x \in P$ such that $x \geq p$ or $x$ is not related to $p$. Moreover, given a pair $x,y \in P_{\nless}$, if $(x \leq z \leq y)$, surely $z\in P_{\nless p}$.      
\end{proof}

\subsubsection{Classes of maps}

Conically smooth stratified maps, by virtue of the potential for differing behaviour along distinct strata, accommodate a strictly larger wealth of sub-classes when compared to the smooth unstratified ones. 


\begin{definition}[Conically smooth submersion]
    \label{ch4:: def: con sm subm}
    Let $f\maps X \longra Y$ be a conically smooth map. We call it a \define{conically smooth submersion} if there is a collection of images $\{\varphi_\alpha(O_\alpha)\subseteq Y\}_\alpha$, indexed by commutative diagrams 
    \begin{equation}\label{eq:def:con sm submersion}
    \begin{tikzcd}[ampersand replacement=\&]
	{U_\alpha\times O_\alpha } \&\& X \\
	\\
	{O_\alpha} \&\& Y
	\arrow["{\psi_{\alpha}}", hook, from=1-1, to=1-3]
	\arrow["{\sf pr}"', from=1-1, to=3-1]
	\arrow["\lrcorner"{anchor=center, pos=0.125}, draw=none, from=1-1, to=3-3]
	\arrow["f", from=1-3, to=3-3]
	\arrow["{\varphi_\alpha}"', hook, from=3-1, to=3-3]
    \end{tikzcd}
    \end{equation}
    where the horizontal arrows are conically smooth open embeddings, and this collection constitutes a basis for the topology of $Y$.
\end{definition}


\begin{remark}\label{ch4:: rmk: con sm submersions}
    Note that \Cref{ch4:: def: con sm subm} looks different from the usual definition for submersive smooth maps. Namely, for categories of smooth manifolds and smooth maps therein, we record submersivity as a condition on the rank of the derivative map. It is a standard fact from differential topology that submersions are locally isomorphic to the standard projection
    \begin{equation}
    \RR^{n+k} \simeq \RR^n \times \RR^k\fib \RR^n.
    \end{equation}
    \Cref{ch4:: def: con sm subm} takes this characterization of submersions as their defining feature.
\end{remark}
\begin{observation}
    An implication of this definition is that, restricting to any stratum $X_p \subset X$, the induced map of smooth manifolds
    \begin{equation}
        f_{|X_p}\maps X_p \longra Y_{f(p)}
    \end{equation}
    is a smooth submersion. In particular, a conically smooth submersion among depth $0$ stratified spaces is an ordinary smooth submersion.
\end{observation}


\begin{definition}[Conically smooth fiber bundles]\label{ch4:: def: con sm fiber bundles}
    Let $f\maps X \longra Y$ be a conically smooth map. We call it a \define{conically smooth fiber bundle} if the collection of images $\varphi_\alpha(O_\alpha)\subseteq Y$, indexed by pullback diagrams 
    \begin{equation}\label{eq:def:fiber bundle}
    \begin{tikzcd}[ampersand replacement=\&]
	{F_\alpha\times O_\alpha } \&\& X \\
	\\
	{O_\alpha} \&\& Y
	\arrow["{\psi_{\alpha}}", hook, from=1-1, to=1-3]
	\arrow["{\sf pr}"', from=1-1, to=3-1]
	\arrow["\lrcorner"{anchor=center, pos=0.125}, draw=none, from=1-1, to=3-3]
	\arrow["f", from=1-3, to=3-3]
	\arrow["{\varphi_\alpha}"', hook, from=3-1, to=3-3]
    \end{tikzcd}
    \end{equation}
    where the horizontal arrows are conically smooth open embeddings, and this collection constitutes a basis for the topology of $Y$.
\end{definition}

\begin{example}
    In the figure below, we depict an example of a conically smooth fiber bundle $\pi \maps X \longra Y$ 
    \begin{equation}
        \tikzfig{figures/Strat/STR15}{\begin{tikzpicture}
	\begin{pgfonlayer}{nodelayer}
		\node [style=none] (0) at (-3, 1) {};
		\node [style=none] (1) at (-3, -0.5) {};
		\node [style=none] (2) at (3, 1) {};
		\node [style=none] (3) at (3, -0.5) {};
		\node [style=none] (4) at (0, 0) {};
		\node [style=none] (5) at (0.35, -0.125) {};
		\node [style=none] (6) at (0, -0.25) {};
		\node [style=none] (7) at (2.35, 0.3) {};
		\node [style=none] (8) at (2.2, 0.425) {};
		\node [style=none] (9) at (-1.75, 0.6) {};
		\node [style=none] (10) at (-2.15, 0.375) {};
		\node [style=0 strat] (11) at (3, -0.5) {};
		\node [style=0 strat] (12) at (3, 1) {};
		\node [style=0 strat] (13) at (-3, 1) {};
		\node [style=0 strat] (14) at (-3, -0.5) {};
		\node [style=none] (15) at (0, -1.5) {};
		\node [style=none] (16) at (0, -4) {};
		\node [style=none] (17) at (-3, -5.5) {};
		\node [style=none] (18) at (3, -5.5) {};
		\node [style=none] (19) at (3, -5.5) {};
		\node [style=none] (20) at (-3, -5.5) {};
		\node [style=0 strat] (21) at (-3, -5.5) {};
		\node [style=0 strat] (22) at (3, -5.5) {};
		\node [style=none] (23) at (-5.25, 0.5) {$X$};
		\node [style=none] (24) at (-5, -5.5) {$Y$};
	\end{pgfonlayer}
	\begin{pgfonlayer}{edgelayer}
		\draw [in=120, out=-90, looseness=0.50] (0.center) to (4.center);
		\draw [in=90, out=90, looseness=0.50] (0.center) to (2.center);
		\draw [in=-90, out=-60, looseness=0.75] (4.center) to (3.center);
		\draw [in=-90, out=30, looseness=0.75] (5.center) to (2.center);
		\draw [in=-150, out=-90, looseness=0.50] (1.center) to (6.center);
		\draw [in=-150, out=90, looseness=0.75] (1.center) to (10.center);
		\draw [in=150, out=15, looseness=0.75] (9.center) to (8.center);
		\draw [in=-30, out=90] (3.center) to (7.center);
		\draw [style=cd-ar] (15.center) to (16.center);
		\draw [in=90, out=90, looseness=0.50] (17.center) to (19.center);
		\draw [in=270, out=-90, looseness=0.50] (17.center) to (19.center);
	\end{pgfonlayer}
\end{tikzpicture}
}
    \end{equation}
    \medskip

    \nid For every point $y \in Y$ of the base, there is a neighborhood $y \in U \subset Y$ such that there is a conically smooth isomorphism
    \begin{equation}
        F \times U \, \simeq  \, \pi^{-1}(U) \simeq \,  U \times_{Y} X,
    \end{equation} where $F$ is the discrete space with two points.
\end{example}

\begin{example}
    Perhaps just as useful to the introduction of conically smooth fiber bundles, is the indication of a good counter-example. The projection map $\pi \maps X \longra Y$, depicted in the diagram below
    \begin{equation}
        \tikzfig{figures/Strat/STR16}{\begin{tikzpicture}
	\begin{pgfonlayer}{nodelayer}
		\node [style=none] (23) at (-5, 0) {$X$};
		\node [style=none] (24) at (-5, -5.5) {$Y$};
		\node [style=none] (25) at (-3, -5.5) {};
		\node [style=none] (26) at (7, -5.5) {};
		\node [style=none] (27) at (-3, -2) {};
		\node [style=none] (28) at (-3, 2) {};
		\node [style=none] (29) at (2, -2) {};
		\node [style=0 strat] (30) at (2, -2) {};
		\node [style=0 strat] (31) at (2, -1) {};
		\node [style=0 strat] (32) at (2, 1) {};
		\node [style=0 strat] (33) at (2, 2) {};
		\node [style=none] (34) at (-3, 2) {};
		\node [style=none] (35) at (7, 2) {};
		\node [style=none] (36) at (-3, -2) {};
		\node [style=none] (37) at (7, -2) {};
		\node [style=none] (38) at (-3, 0.25) {};
		\node [style=none] (39) at (7, 0) {};
		\node [style=none] (40) at (2, -2.75) {};
		\node [style=none] (41) at (2, -4.75) {};
		\node [style=0 strat] (42) at (2, -5.5) {};
	\end{pgfonlayer}
	\begin{pgfonlayer}{edgelayer}
		\draw (25.center) to (26.center);
		\draw [style=defect patch b] (36.center)
			 to (34.center)
			 to (35.center)
			 to (37.center)
			 to cycle;
		\draw (34.center) to (33);
		\draw (33) to (35.center);
		\draw (36.center) to (30);
		\draw (30) to (37.center);
		\draw (33) to (32);
		\draw (32) to (31);
		\draw (31) to (30);
		\draw [in=-180, out=0] (38.center) to (32);
		\draw [in=180, out=0] (31) to (39.center);
		\draw [style=cd-ar] (40.center) to (41.center);
	\end{pgfonlayer}
\end{tikzpicture}
}
    \end{equation}
    \medskip 

    \nid is not a conically smooth fiber bundle. Any conically smooth open embedding $U \hookra Y$ which contains the 0-stratum in the base, will not be in the basis for a topology of $Y$, since restricting to $U$ does not yield a stratified projection of a fixed stratified fiber $F$:
    \begin{equation}
        \tikzfig{figures/Strat/STR17}{\begin{tikzpicture}
	\begin{pgfonlayer}{nodelayer}
		\node [style=none] (25) at (0.75, -5.5) {};
		\node [style=none] (26) at (3.25, -5.5) {};
		\node [style=none] (27) at (0.75, -2) {};
		\node [style=none] (28) at (0.75, 2) {};
		\node [style=none] (29) at (2, -2) {};
		\node [style=0 strat] (30) at (2, -2) {};
		\node [style=0 strat] (31) at (2, -1) {};
		\node [style=0 strat] (32) at (2, 1) {};
		\node [style=0 strat] (33) at (2, 2) {};
		\node [style=none] (34) at (0.75, 2) {};
		\node [style=none] (35) at (3.25, 2) {};
		\node [style=none] (36) at (0.75, -2) {};
		\node [style=none] (37) at (3.25, -2) {};
		\node [style=none] (38) at (0.75, 0.75) {};
		\node [style=none] (39) at (3.25, -0.75) {};
		\node [style=none] (40) at (2, -2.75) {};
		\node [style=none] (41) at (2, -4.75) {};
		\node [style=0 strat] (42) at (2, -5.5) {};
		\node [style=none] (43) at (-1.5, -5.5) {$U$};
		\node [style=none] (44) at (-2.25, 0) {$U \times_{Y} X$ };
	\end{pgfonlayer}
	\begin{pgfonlayer}{edgelayer}
		\draw (25.center) to (26.center);
		\draw [style=defect patch b] (36.center)
			 to (34.center)
			 to (35.center)
			 to (37.center)
			 to cycle;
		\draw (34.center) to (33);
		\draw (33) to (35.center);
		\draw (36.center) to (30);
		\draw (30) to (37.center);
		\draw (33) to (32);
		\draw (32) to (31);
		\draw (31) to (30);
		\draw [in=-180, out=0] (38.center) to (32);
		\draw [in=180, out=0] (31) to (39.center);
		\draw [style=cd-ar] (40.center) to (41.center);
		\draw [style=ts-dashed] (34.center) to (36.center);
		\draw [style=ts-dashed] (35.center) to (37.center);
	\end{pgfonlayer}
\end{tikzpicture}
}
    \end{equation}
    \medskip
    
    \nid Modifying the total space   
    \begin{equation}
        \tikzfig{figures/Strat/STR17}{\begin{tikzpicture}
	\begin{pgfonlayer}{nodelayer}
		\node [style=none] (25) at (0.75, -5.5) {};
		\node [style=none] (26) at (3.25, -5.5) {};
		\node [style=none] (27) at (0.75, -2) {};
		\node [style=none] (28) at (0.75, 2) {};
		\node [style=none] (29) at (2, -2) {};
		\node [style=0 strat] (30) at (2, -2) {};
		\node [style=0 strat] (31) at (2, -1) {};
		\node [style=0 strat] (32) at (2, 1) {};
		\node [style=0 strat] (33) at (2, 2) {};
		\node [style=none] (34) at (0.75, 2) {};
		\node [style=none] (35) at (3.25, 2) {};
		\node [style=none] (36) at (0.75, -2) {};
		\node [style=none] (37) at (3.25, -2) {};
		\node [style=none] (38) at (0.75, 0.75) {};
		\node [style=none] (39) at (3.25, -0.75) {};
		\node [style=none] (40) at (2, -2.75) {};
		\node [style=none] (41) at (2, -4.75) {};
		\node [style=0 strat] (42) at (2, -5.5) {};
		\node [style=none] (43) at (-1.5, -5.5) {$U$};
		\node [style=none] (44) at (-2.25, 0) {$U \times_{Y} X$ };
	\end{pgfonlayer}
	\begin{pgfonlayer}{edgelayer}
		\draw (25.center) to (26.center);
		\draw [style=defect patch b] (36.center)
			 to (34.center)
			 to (35.center)
			 to (37.center)
			 to cycle;
		\draw (34.center) to (33);
		\draw (33) to (35.center);
		\draw (36.center) to (30);
		\draw (30) to (37.center);
		\draw (33) to (32);
		\draw (32) to (31);
		\draw (31) to (30);
		\draw [in=-180, out=0] (38.center) to (32);
		\draw [in=180, out=0] (31) to (39.center);
		\draw [style=cd-ar] (40.center) to (41.center);
		\draw [style=ts-dashed] (34.center) to (36.center);
		\draw [style=ts-dashed] (35.center) to (37.center);
	\end{pgfonlayer}
\end{tikzpicture}
} \quad \squigra \quad \tikzfig{figures/Strat/STR18}{\begin{tikzpicture}
	\begin{pgfonlayer}{nodelayer}
		\node [style=none] (25) at (0.5, -5.5) {};
		\node [style=none] (26) at (3.5, -5.5) {};
		\node [style=none] (27) at (0.5, -2) {};
		\node [style=none] (28) at (0.5, 2) {};
		\node [style=none] (29) at (2, -2) {};
		\node [style=0 strat] (30) at (2, -2) {};
		\node [style=0 strat] (31) at (2, 0) {};
		\node [style=0 strat] (32) at (2, 0) {};
		\node [style=0 strat] (33) at (2, 2) {};
		\node [style=none] (34) at (0.5, 2) {};
		\node [style=none] (35) at (3.5, 2) {};
		\node [style=none] (36) at (0.5, -2) {};
		\node [style=none] (37) at (3.5, -2) {};
		\node [style=none] (38) at (0.5, 0.25) {};
		\node [style=none] (39) at (3.5, 0.25) {};
		\node [style=none] (40) at (2, -2.75) {};
		\node [style=none] (41) at (2, -4.75) {};
		\node [style=0 strat] (42) at (2, -5.5) {};
	\end{pgfonlayer}
	\begin{pgfonlayer}{edgelayer}
		\draw (25.center) to (26.center);
		\draw [style=defect patch b] (36.center)
			 to (34.center)
			 to (35.center)
			 to (37.center)
			 to cycle;
		\draw (34.center) to (33);
		\draw (33) to (35.center);
		\draw (36.center) to (30);
		\draw (30) to (37.center);
		\draw (33) to (32);
		\draw (32) to (31);
		\draw (31) to (30);
		\draw [in=-180, out=0] (38.center) to (32);
		\draw [in=180, out=0] (31) to (39.center);
		\draw [style=cd-ar] (40.center) to (41.center);
		\draw [style=ts-dashed] (34.center) to (36.center);
		\draw [style=ts-dashed] (35.center) to (37.center);
	\end{pgfonlayer}
\end{tikzpicture}
} 
    \end{equation}
    \medskip

    \nid turns it into a conically smooth fiber bundle with fiber 
    \begin{equation}
     \tikzfig{figures/Strat/STR19}{\begin{tikzpicture}
	\begin{pgfonlayer}{nodelayer}
		\node [style=none] (29) at (2, -2) {};
		\node [style=0 strat] (30) at (2, -2) {};
		\node [style=0 strat] (31) at (2, 0) {};
		\node [style=0 strat] (32) at (2, 0) {};
		\node [style=0 strat] (33) at (2, 2) {};
		\node [style=none] (34) at (0.25, 0) {$F$  };
		\node [style=none] (35) at (1, -0.075) {$=$};
	\end{pgfonlayer}
	\begin{pgfonlayer}{edgelayer}
		\draw (33) to (32);
		\draw (32) to (31);
		\draw (31) to (30);
	\end{pgfonlayer}
\end{tikzpicture}
} \quad .   
    \end{equation}
\end{example}

\begin{definition}[(Weakly) Constructible bundles]\label{ch4:: def: (weak) cbl}
    Let $f \maps \big(X \ra P \big) \longra \big(Y \ra Q \big)$ be a conically smooth stratified map. We call it a \define{weakly constructible bundle} if, for each $q\in Q$, the vertical map in the pullback
    \begin{equation}
        \begin{tikzcd}[ampersand replacement=\&]
            {X_q} \&\& X \\
            \\
            {Y_q} \&\& Y
            \arrow[hook, from=1-1, to=1-3]
            \arrow[from=1-1, to=3-1]
            \arrow["\lrcorner"{anchor=center, pos=0.125}, draw=none, from=1-1, to=3-3]
            \arrow["f", from=1-3, to=3-3]
            \arrow[hook, from=3-1, to=3-3]
        \end{tikzcd}
    \end{equation} 
    is a conically smooth fiber bundle (\Cref{ch4:: def: con sm fiber bundles}). The map $f$ is said to be a \define{constructible bundle} if:
    \begin{itemize}
        \item For $\depth(X) = 0$, the map $f$ is a smooth fiber bundle of smooth manifolds \footnote{If $\depth(X)=0$ even if $\depth(Y)>0$ the map must factor through a single stratum of $Y$.}
        \item For $\depth(X)>0$, then:
        \begin{itemize}
            \item $f$ is a weakly constructible bundle.
            \item For each $q\in Q$ the map 
            \begin{equation}
                \link_{X_{q}}(X) \longra X_q \underset{Y_{q}}{\times} \link_{Y_{q}}(Y)
            \end{equation}
            is a constructible bundle.
        \end{itemize}
    \end{itemize}
\end{definition}

\begin{observation}
Note that for any given sub-stratified space $Y \hookra X$, the conically smooth stratified space $\link_{Y}(X)$ has strictly less depth than $X$, which validates the inductive step.
\end{observation}

\begin{example}
    Let $f\maps X \longra Y$ be the projection 
    \begin{equation}
    \sC(\scalebox{0.3}{\tikzfig{figures/Strat/STR28}{\begin{tikzpicture}
	\begin{pgfonlayer}{nodelayer}
		\node [style=none] (3) at (0, -0.5) {};
		\node [style={wide_dot}] (4) at (-2, 0.5) {};
		\node [style={wide_dot}] (5) at (0, 0.5) {};
		\node [style={wide_dot}] (6) at (2, 0.5) {};
	\end{pgfonlayer}
	\begin{pgfonlayer}{edgelayer}
		\draw (4) to (5);
		\draw (5) to (6);
	\end{pgfonlayer}
\end{tikzpicture}
}}) \longra \RR_{\geq 0}
    \end{equation} 
    which we diagramatically depict as 
    \begin{equation}
        \tikzfig{figures/Strat/STR29}{\begin{tikzpicture}
	\begin{pgfonlayer}{nodelayer}
		\node [style=none] (6) at (2.5, -3) {};
		\node [style=none] (4) at (-2.5, 0) {};
		\node [style=none] (5) at (2.5, 3) {};
		\node [style=0 strat] (0) at (-2.5, 0) {};
		\node [style=none] (1) at (2.5, 3) {};
		\node [style=none] (2) at (2.5, -3) {};
		\node [style=none] (3) at (2.5, 0) {};
		\node [style=none] (7) at (0, -4) {};
		\node [style=none] (8) at (0, -6) {};
		\node [style=0 strat] (9) at (-2.5, -7) {};
		\node [style=none] (10) at (2.5, -7) {};
		\node [style=none] (11) at (-5, -0.25) {$X$};
		\node [style=none] (12) at (-5, -7) {$Y$};
	\end{pgfonlayer}
	\begin{pgfonlayer}{edgelayer}
		\draw [style=defect patch b] (6.center)
			 to (4.center)
			 to (5.center)
			 to cycle;
		\draw (0) to (1.center);
		\draw (0) to (3.center);
		\draw (0) to (2.center);
		\draw [style=ts-dashed] (1.center) to (2.center);
		\draw [style=cd-ar] (7.center) to (8.center);
		\draw (9) to (10.center);
	\end{pgfonlayer}
\end{tikzpicture}
}
    \end{equation}
    \medskip

    \nid This is a weakly constructible bundle since for every stratum of the base, the pullback is a conically smooth fiber bundle. The fiber over $\{0\}$ is the cone point and the fiber over any point in the interior stratum $\RR_{>0}$ is an interval stratified by the boundary and a single $0$-stratum in the interior.
\end{example}

A key fact, and a hard result to prove, is that weakly constructible bundles, unlike constructible bundles, do not in general pull back along arbitrary conically smooth stratified maps.

\begin{theorem}[Pullback stability]\label{ch4:: thm: cbl pullback}\cite[Lemma~6.1.11]{AFR1}
    Let $f \maps X \longra Y$ be a constructible bundle and $g\maps Y' \longra Y$ an arbitrary conically smooth stratified map, the pullback 
    \begin{equation}
        \begin{tikzcd}[ampersand replacement=\&]
            {X \underset{Y}{\times} Y'} \&\& X \\
            \\
            {Y'} \&\& Y
            \arrow[from=1-1, to=1-3]
            \arrow["{f'}"', from=1-1, to=3-1]
            \arrow["\lrcorner"{anchor=center, pos=0.125}, draw=none, from=1-1, to=3-3]
            \arrow["f", from=1-3, to=3-3]
            \arrow["g"', from=3-1, to=3-3]
        \end{tikzcd}
    \end{equation}
    exists in the category $\strat$. The map $f'\maps X\underset{Y}{\times} Y' \longra Y'$ is again a constructible bundle. It is also proper if $f$ is proper.
\end{theorem}


\begin{definition}[Refinement]\label{ch4:: def: refinement}
    A map $f\maps X \longra Y$ is a \define{refinement} if it is a homeomorphism of underlying topological spaces, and for each $p\in P$
    \begin{equation}
        f_{|X_p}\maps X_p \longra Y_{f(p)}
    \end{equation} 
    is an isomorphism onto its image as smooth manifolds.
\end{definition}

\begin{example}
    This diagram depicts a refinement:
    \begin{equation}
        \tikzfig{figures/Strat/STR25}{\begin{tikzpicture}
	\begin{pgfonlayer}{nodelayer}
		\node [style=none] (27) at (-1.25, -2) {};
		\node [style=none] (28) at (-1.25, 2) {};
		\node [style=none] (29) at (0, -2) {};
		\node [style=0 strat] (30) at (0, -2) {};
		\node [style=0 strat] (31) at (0, -1) {};
		\node [style=0 strat] (32) at (0, 1) {};
		\node [style=0 strat] (33) at (0, 2) {};
		\node [style=none] (34) at (-1.25, 2) {};
		\node [style=none] (35) at (1.25, 2) {};
		\node [style=none] (36) at (-1.25, -2) {};
		\node [style=none] (37) at (1.25, -2) {};
		\node [style=none] (38) at (-1.25, 0.75) {};
		\node [style=none] (39) at (1.25, -0.75) {};
	\end{pgfonlayer}
	\begin{pgfonlayer}{edgelayer}
		\draw [style=defect patch b] (36.center)
			 to (34.center)
			 to (35.center)
			 to (37.center)
			 to cycle;
		\draw (34.center) to (33);
		\draw (33) to (35.center);
		\draw (36.center) to (30);
		\draw (30) to (37.center);
		\draw (33) to (32);
		\draw (32) to (31);
		\draw (31) to (30);
		\draw [in=-180, out=0] (38.center) to (32);
		\draw [in=180, out=0] (31) to (39.center);
		\draw [style=ts-dashed] (34.center) to (36.center);
		\draw [style=ts-dashed] (35.center) to (37.center);
	\end{pgfonlayer}
\end{tikzpicture}
} \;\;\xra{\,\,\, f \,\,\,}\;\; \tikzfig{figures/Strat/STR26}{\begin{tikzpicture}
	\begin{pgfonlayer}{nodelayer}
		\node [style=none] (27) at (-1.25, -2) {};
		\node [style=none] (28) at (-1.25, 2) {};
		\node [style=none] (29) at (0, -2) {};
		\node [style=0 strat] (30) at (0, -2) {};
		\node [style=0 strat] (32) at (0, 1) {};
		\node [style=0 strat] (33) at (0, 2) {};
		\node [style=none] (34) at (-1.25, 2) {};
		\node [style=none] (35) at (1.25, 2) {};
		\node [style=none] (36) at (-1.25, -2) {};
		\node [style=none] (37) at (1.25, -2) {};
		\node [style=none] (38) at (-1.25, 0.75) {};
		\node [style=none] (39) at (1.25, -0.75) {};
	\end{pgfonlayer}
	\begin{pgfonlayer}{edgelayer}
		\draw [style=defect patch b] (36.center)
			 to (34.center)
			 to (35.center)
			 to (37.center)
			 to cycle;
		\draw (34.center) to (33);
		\draw (33) to (35.center);
		\draw (36.center) to (30);
		\draw (30) to (37.center);
		\draw (33) to (32);
		\draw [style=ts-dashed] (34.center) to (36.center);
		\draw [style=ts-dashed] (35.center) to (37.center);
		\draw (32) to (30);
	\end{pgfonlayer}
\end{tikzpicture}
}
    \end{equation}
    \medskip

    \nid which is the identity on the underlying topological space. Since it is stratified, it implies that $f_{|X_{p}}$ is an isomorphism onto its image. 
\end{example}



\subsection{Differential topology of conically smooth stratified spaces}
In this subsection we extend classical results from differential topology from smooth manifolds to conically smooth stratified spaces. The key constructions here will be: 
\begin{enumerate}
    \item The $(\oo,1)$-category $\Strat$, a suitable  $\mathsf{Kan}$-enrichment of $\strat$ .
    \item Categories of basics.
    \item Conically smooth stratified spaces with corners.
    \item Unzipping: a suitably functorial way of resolving singularities, akin to \emph{blow-up} in algebraic geometry.
    \item Tubular neighborhoods around strata and constructible sub-stratifications.
\end{enumerate}


One of the features of smooth manifolds which sets them apart from algebraic varieties, for example, is the ability to patch together local data and constructions into global ones using partitions of unity. These are likewise available for conically smooth stratified spaces.

\begin{lemma}[Conically smooth partitions of unity]\label{lem:partitions of unity}\cite[Lemma 7.1.1]{AFT1}
    Every conically smooth stratified space admits a conically smooth partition of unity.
\end{lemma}

\subsubsection{\texorpdfstring{Topologizing $\strat$}{Topologizing strat}}

For finer details about this construction, see Sections~2.1--2.4 in \cite{AFR1} — we will give only a schematic overview. 

We observe that the $\oo$-category of topological spaces can be defined as the Dwyer--Kan localization of $\Top$ along the set of homotopy equivalences. This is a strategy we can aptly reproduce in the case of conically smooth stratified spaces if we first define what is meant by \define{stratified homotopy equivalence}. 

\begin{definition}[Stratified homotopy equivalence]\label{ch4:: def: stratified homotopy equivalence}\cite[Definition~2.1.1]{AFR1}
    Given a pair of conically smooth stratified maps $f,g\maps X \longra Y$, a \define{stratified homotopy} 
    \begin{equation}
        H \maps X \times \RR \longra Y \times \RR
    \end{equation}
    is a conically smooth stratified map $H$ for which the following commutative diagram commutes:
    \begin{equation}
        \begin{tikzcd}[ampersand replacement=\&]
            X \&\& Y \\
            {X\times \RR} \&\& {Y \times \RR} \\
            X \&\& Y
            \arrow["f", from=1-1, to=1-3]
            \arrow["{\{0\}}"', from=1-1, to=2-1]
            \arrow["{\{0\}}", from=1-3, to=2-3]
            \arrow["H"{description}, from=2-1, to=2-3]
            \arrow["{\{1\}}", from=3-1, to=2-1]
            \arrow["g"', from=3-1, to=3-3]
            \arrow["{\{1\}}"', from=3-3, to=2-3]
        \end{tikzcd}
    \end{equation}
    The map $f$ is said to be a \define{stratified homotopy equivalence} if there exists a conically smooth stratified map $g$ together with stratified homotopies from $gf$ and $fg$ to $\id_{X}$ and $\id_{Y}$, respectively. Let us denote with $\cJ$ the collection of all stratified homotopy equivalences. 
\end{definition}

Our first pass at a definition is the following.

\begin{definition}
    Let 
    \begin{equation}
    \strat[\cJ^{-1}]
    \end{equation} 
    be the Dwyer--Kan localization of $\strat$ along the stratified homotopy equivalences. We denote the associated $(\oo,1)$-category by $\Strat$.
\end{definition}

\begin{observation}
This definition captures the correct homotopy theory of conically smooth stratified spaces, but as with most naive localizations, it does not offer a concrete model — an equivalent $(\oo,1)$-category modeled after a simplicially enriched category, for example.   
\end{observation}

A first simplification is to localize along a narrower class of morphisms.

\begin{definition}
    Let $X \times \RR \xra{\; \mathsf{pr}\;} X$ be the canonical projection. Define $\cR$ to be the collection of all projections for all $X\in \strat$. This collection is manifestly contained in the stratified homotopy equivalences, i.e.
        \begin{equation}
        \cR \subset \cJ.
        \end{equation}   
\end{definition}

\begin{lemma}\cite[Lemma~2.1.7]{AFR1}
    The canonical map between Dwyer--Kan localizations 
    \begin{equation}
        \strat[\cR^{-1}] \longra \strat[\cJ^{-1}]
    \end{equation}
    is an equivalence of $(\oo,1)$-categories.
\end{lemma}
\begin{proof}[Proof sketch]
    Let $F \maps \strat \longra \cC$ be a functor which sends every projection in $\cR$ to an equivalence in $\cC$. For each $X\in \strat$ and each endpoint inclusion $\{i\}\maps X \hookra X\times \RR$ with $i\in \{0,1\}$, the composite
    \begin{equation}
        X \xra{\;\{i\}\;} X\times \RR \xra{\;\mathsf{pr}\;} X
    \end{equation}
    is the identity. Since $F(\mathsf{pr})$ is an equivalence, the 2-out-of-3 property implies that $F(\{i\})$ is also an equivalence.

    Now let $H\maps X\times \RR \longra Y\times \RR$ be a stratified homotopy between $f,g\maps X\longra Y$. By definition of stratified homotopy, we have
    \begin{equation}
        f = \mathsf{pr}_Y \circ H \circ \{0\}, \qquad g = \mathsf{pr}_Y \circ H \circ \{1\}.
    \end{equation}
    After applying $F$, the morphisms $F(f)$ and $F(g)$ differ only by replacing the endpoint inclusion $F(\{0\})$ with $F(\{1\})$, and both of these are equivalences. Hence $F(f)$ is an equivalence if and only if $F(g)$ is.

    Finally, if $f\maps X\longra Y$ is a stratified homotopy equivalence with homotopy inverse $g\maps Y\longra X$, then the homotopies $g\circ f \simeq \id_X$ and $f\circ g \simeq \id_Y$ imply that $F(g)\circ F(f)$ and $F(f)\circ F(g)$ are equivalent to identities. Therefore $F(f)$ is itself an equivalence. This shows that every functor inverting $\cR$ also inverts $\cJ$, which is exactly the required universal property.
\end{proof}

We regard $\strat$ as a site by declaring a family of morphisms $\{U_i\hookra X\}_{i\in I}$ to be a covering family if each $U_i\hookra X$ is a conically smooth open embedding and the opens $\{U_i\}_{i\in I}$ cover the underlying topological space of $X$.
Write
\[
  \psh(\strat)\;:=\;\Fun(\strat^{\op},\Spaces)
\]
for the $(\oo,1)$-category of presheaves, and
\[
  \shv(\strat)\;\subset\;\psh(\strat)
\]
for the full subcategory of (hypercomplete) $\infty$-sheaves for this Grothendieck topology.


\begin{definition}
    The $\oo$-category of \define{constructible sheaves} on $\strat$ is the full $(\oo,1)$-subcategory 
    \begin{equation}
        \shv^{\cbl}(\strat) \; \subset \; \shv(\strat).
    \end{equation}
\end{definition}

In particular, if we restrict a constructible sheaf $\cF$ to the subsite $\man \hookra \strat$ of smooth manifolds we obtain a concordance-invariant sheaf. Contrary to the case of smooth manifolds, the concordance resolution 
\begin{equation}
    \psh(\strat) \longra \psh(\strat), \quad \cF \longmapsto \left| \cF(-\times \AA^{\bullet}) \right|
\end{equation}
does not, in general, report back a constructible sheaf. In Section 2.3 of \cite{AFR1} the authors provide a class of sheaves whose concordance resolution (alias \emph{topologizing diagram} in \cite{AFR1}) is constructible — see \cite[Theorem~2.3.14]{AFR1}. These are called isotopy sheaves, and as demonstrated in the following definition, have a very well-defined list of checkable criteria.
\begin{definition}\cite[Definition~2.3.8]{AFR1}
    The category of \define{isotopy sheaves} is the full subcategory 
    \begin{equation}
        \mathsf{Isot} \subset \Cat_{/\strat}
    \end{equation}
    consisting of right fibrations $\cI \longra \strat$ whose straightening $\cI \maps \strat^{\op} \longra \mathsf{Gpd}$ satisfies the following conditions:
    \begin{enumerate}
        \item \textbf{Sheaf:} The straightening $\cI\maps \strat^{\op} \longra \mathsf{Gpd}$ is a sheaf. 
        \item \textbf{Isotopy extension:} For each stratified space $K$, the following conditions are satisfied:
        \begin{enumerate}
            \item For each $q\geq 0$, the functor between groupoids 
            \begin{equation}
                \cI(K \times \AA^q) \longra \cI(K \times \partial \AA^q)
            \end{equation}
            is an isofibration.
            \item For each $0\leq i \leq q$, the functor between groupoids
            \begin{equation}
                \cI(K \times \AA^q) \longra \cI(K \times (\Lambda^q)_i)
            \end{equation}
            is surjective on objects and on Hom-sets.
        \end{enumerate}
    \end{enumerate}
\end{definition}

There is a limit-preserving natural inclusion $\shv^{\cbl}(\strat) \xlonghookra{\;\;} \shv(\strat)$ which admits a left adjoint:
\begin{equation}
    L\maps \shv(\strat) \longra \shv^{\cbl}(\strat).
\end{equation}
When restricted to the subcategory $\mathsf{Isot}\subset \shv(\strat)$, the value of this functor is not only a constructible sheaf, it is its associated concordance resolution.
\begin{proposition}\cite[Lemma~2.3.16]{AFR1}\label{ch4:: def: isotopy sheaf realizes to cbl}
    If $\cI$ is an isotopy sheaf, then there is a canonical natural transformation 
    \begin{equation}
        L\cI \longra \left| \cI(-\times \AA^\bullet) \right|,
    \end{equation}
    and it is an equivalence.
\end{proposition} 
We can use this result to provide a $\mathsf{Kan}$-enrichment of $\strat$ as follows:
\begin{enumerate}
    \item For every $K\in \strat$, the forgetful functor 
    \begin{equation}
        \strat_{/K} \longra \strat
    \end{equation} 
    which straightens to the representable $\strat(-,K)$, is an isotopy sheaf. The verification uses the existence of regular neighborhoods for constructible subspaces, a fact we establish later in the chapter.
    \item  The functor
    \begin{equation}
        \strat \longra \shv(\strat), \quad K \longmapsto \strat(-,K)
    \end{equation}
    lands in $\mathsf{Isot}$. Applying $L$, produces a functor 
    \begin{equation}
        \strat \longra \shv^{\cbl}(\strat), \quad K \longmapsto L\strat(-,K).
    \end{equation}
\end{enumerate}
These observations lead us to following definition.
\begin{definition}\cite[Definition~2.4.2]{AFR1}
    The $(\oo,1)$-category $\Strat$ of conically smooth stratified spaces is the essential image of 
    \begin{equation}
    \strat \longra \shv^{\cbl}(\strat), \qquad K \longmapsto L\strat(-,K).
    \end{equation}
\end{definition}
The $\mathsf{Kan}$-enrichment is then derived from \Cref{ch4:: def: isotopy sheaf realizes to cbl} as follows: given a pair of conically smooth stratified spaces $X,Y\in \strat$, the associated simplicial set of morphisms from $X$ to $Y$ is computed by 
\begin{equation}
    \Hom_{\shv^{\cbl}}\left(L\strat(-,X),L\strat(-,Y)\right).
\end{equation} 
Using the adjunction $L \dashv i$, we have that 
\begin{align}
    \Hom_{\shv^{\cbl}}\left(L\strat(-,X),L\strat(-,Y)\right) &\simeq \Hom_{\shv}\left(\strat(-,X),iL\strat(-,Y)\right) \\
    &\simeq iL\strat(-,Y)(X) \\
    &= L\strat(-,Y)(X).
\end{align}
From \Cref{ch4:: def: isotopy sheaf realizes to cbl} we have that 
\begin{equation}
    L\strat(-,Y) \simeq \left|\strat(-\times \AA^\bullet,Y)\right|,
\end{equation}
which implies that 
\begin{equation}
    \Hom_{\shv^{\cbl}}(X,Y) \,\, \simeq \,\, L\strat(-,Y)(X) \,\, \simeq \,\, \left|\strat(X\times \AA^{\bullet},Y)\right|,
\end{equation} where we are interpreting $X$ and $Y$ as their images under the localization.

\begin{theorem}\cite[Theorem 2.4.5]{AFR1}\label{ch4:: thm: Strat enrichment}
    The $(\oo,1)$-category obtained through the localization $\strat[\cJ^{-1}]$ is equivalent to the $(\oo,1)$-category $\Strat$ obtained as the $\mathsf{Kan}$-enrichment of $\strat$ via:
    \begin{equation}\label{eq: st enrich}
        \Hom_{\Strat}(X,Y) \,\, \overset{\sf def}{=} \,\, \Strat(X\times \AA^{\bullet},Y).
    \end{equation}
\end{theorem}

\begin{observation}
    We could equivalently describe the $\mathsf{Kan}$-enrichment~(\ref{eq: st enrich}) of $\strat$ by 
\begin{equation}
    \Strat(X,Y) \,\, \overset{\sf def}{=}\,\, \strat_{/\AA^{\bullet}}(X\times \AA^{\bullet},Y\times \AA^{\bullet}).
\end{equation}
\end{observation}

\begin{notation}
    We denote by $\Bsc$ and $\Snglr$ as the $\mathsf{Kan}$-enriched subcategories of $\Strat$ whose objects are those of $\bsc$ and $\snglr$, respectively. In particular, we will refer to these not as $\mathsf{Kan}$-enrichments but as $(\oo,1)$-categories.
\end{notation}

\medskip

\medskip

\subsubsection{Categories of basics}

In \Cref{ch3:: sec: tangential structures} we discussed the notion of tangential structures for smooth manifolds. In this section we aim to introduce an analog of $\BO(d)$ and the tangent classifier $M \longra \BO(d)$, i.e.\ the map which classifies the tangent bundle of the $d$-dimensional smooth manifold $M$. The goal being to introduce tangential structures and singularity types.

As is observed in the introduction of \cite{AFT1}, a space over $\BO(d)$, which we can assume to be a Kan fibration, is equivalent data to that of a functor $\BO(d)^{\op} \longra \Top$. Through the marvelous observation that 
\begin{equation}
    \sO(d) \longra \Emb_{\BO(d)}(\RR^{d},\RR^{d})
\end{equation} 
is a weak homotopy equivalence — a combination of the Gram--Schmidt orthogonalization and taking derivatives at the origin — we have that 
\begin{equation}
    \psh(\sD_{d}) \, \simeq \, \Top_{/\BO(d)}
\end{equation}
between presheaves on $\sD_d$ and the spaces mapping into $\BO(d)$. Here we are defining $\sD_d$ as $\man^{\emb}$ restricted to the $d$-dimensional Euclidean spaces. 

Recall that every $p$-stratum of a conically smooth stratified space $(X \ra P)$ is a $d_{p}$-dimensional smooth manifold. Therefore, a suitable proposed generalization of $\sD_d$ is $\Bsc$, the $(\oo,1)$-category of all basics and conically smooth open embeddings thereamong.

\begin{definition}\cite[Definition~1.1.3]{AFT1}
    The \define{tangent classifier} is the composite functor
    \begin{equation}
        \tau \maps \Snglr \xra{\,\,\cY\,\,} \psh(\Snglr) \xra{\,\,\iota^*\,\,} \psh(\Bsc).
    \end{equation}
\end{definition}

\begin{remark}
To a conically smooth stratified space $M$, the functor $\tau$ associates to it the presheaf
\begin{equation}
    \Snglr(-,M)\maps \Bsc^{\op} \longra \Top
\end{equation}
\end{remark}

\begin{observation}
Restricting the tangent classifier $\tau$ to $\Man_{d}^{\emb}$ it is proved in \cite[Corollary~4.4.9]{AFT1} that its image lands in $\psh(\sD_d) \subset \psh(\Bsc)$. More concretely, it is the usual tangent classifier for smooth manifolds which classifies the tangent bundle. 
\end{observation}

\begin{remark}
    Let $\entr(M \ra P)$ denote the $(\oo,1)$-category which is the unstraightening of $\tau(M)$. Unpacking this we note that since 
    \begin{equation}
        \tau(M\ra P) = \Snglr(\iota(-),M)\maps \Bsc^{\op} \longra \Top,
    \end{equation}
    then its unstraightening is the $(\oo,1)$-overcategory 
    \begin{equation}
        \Bsc_{/M}.
    \end{equation}
    We denote this $(\oo,1)$-category as the \define{enter-path category} of $M$. The enter-path category comes with a natural projection functor 
    \begin{equation}
        \Bsc_{/M} \longra \Bsc
    \end{equation}
    which is a right fibration.
\end{remark}

\begin{example}
    Let us consider a simple example, let $M= \sC(* \sqcup *)$
    \begin{equation}
       M \; = \; \tikzfig{figures/Strat/STR30}{\begin{tikzpicture}
	\begin{pgfonlayer}{nodelayer}
		\node [style=none] (0) at (-3, 0) {};
		\node [style=none] (1) at (3, 0) {};
		\node [style=0 strat] (2) at (0, 0) {};
	\end{pgfonlayer}
	\begin{pgfonlayer}{edgelayer}
		\draw (0.center) to (1.center);
	\end{pgfonlayer}
\end{tikzpicture}
}\, .
    \end{equation}
    Objects in $\entr(M)$ are of essentially three types:
    \begin{itemize}
        \item Open embeddings of open intervals into the left of the origin
        \begin{equation}
            \tikzfig{figures/Strat/STR31}{\begin{tikzpicture}
	\begin{pgfonlayer}{nodelayer}
		\node [style=none] (0) at (-3, 0) {};
		\node [style=none] (1) at (3, 0) {};
		\node [style=0 strat] (2) at (0, 0) {};
		\node [style=none] (3) at (-2, 0) {};
		\node [style=none] (4) at (-1, 0) {};
	\end{pgfonlayer}
	\begin{pgfonlayer}{edgelayer}
		\draw (0.center) to (1.center);
		\draw [style={open_int}] (3.center) to (4.center);
	\end{pgfonlayer}
\end{tikzpicture}
}.
        \end{equation} 
        \medskip
        \item Open embeddings of open intervals into the right of the origin
        \begin{equation}
            \tikzfig{figures/Strat/STR32}{\begin{tikzpicture}
	\begin{pgfonlayer}{nodelayer}
		\node [style=none] (0) at (-3, 0) {};
		\node [style=none] (1) at (3, 0) {};
		\node [style=0 strat] (2) at (0, 0) {};
		\node [style=none] (3) at (1, 0) {};
		\node [style=none] (4) at (2, 0) {};
	\end{pgfonlayer}
	\begin{pgfonlayer}{edgelayer}
		\draw (0.center) to (1.center);
		\draw [style={open_int}] (3.center) to (4.center);
	\end{pgfonlayer}
\end{tikzpicture}
}.
        \end{equation} 
        \medskip
        \item Open embeddings of $M$ into itself
        \begin{equation}
            \tikzfig{figures/Strat/STR33}{\begin{tikzpicture}
	\begin{pgfonlayer}{nodelayer}
		\node [style=none] (0) at (-3, 0) {};
		\node [style=none] (1) at (3, 0) {};
		\node [style=0 strat] (2) at (0, 0) {};
		\node [style=none] (3) at (1, 0) {};
		\node [style=none] (4) at (-1, 0) {};
	\end{pgfonlayer}
	\begin{pgfonlayer}{edgelayer}
		\draw (0.center) to (1.center);
		\draw [style={open_int}] (4.center) to (3.center);
	\end{pgfonlayer}
\end{tikzpicture}
}.
        \end{equation} 
        \medskip
    \end{itemize}
    Morphisms are also simple. There are no morphisms from objects of the first type into the second type, and vice-versa. There is essentially one morphism from an object of the first or second type into one of the third type, depicted below as follows:
\end{example}

This lead us naturally to the following generalization of $B$-framed manifolds.

\begin{definition} \cite[Definition~5.0.2]{AFT1} \label{ch4:: def: categories of basics}
    An $(\oo,1)$\define{-category of basics} is a right fibration $\cB \longra \Bsc$. Given such an $(\oo,1)$-categoy of basics $\cB$, the $(\oo,1)$-category of $\cB$-manifolds is the pullback $(\oo,1)$-category 
    \begin{equation}
        \begin{tikzcd}[ampersand replacement=\&]
            {\Mfld(\cB)} \&\& {\left( \rfib_{\Bsc} \right)_{/\cB}} \\
            \\
            \Snglr \&\& {\rfib_{\Bsc}}
            \arrow[from=1-1, to=1-3]
            \arrow[from=1-1, to=3-1]
            \arrow["\lrcorner"{anchor=center, pos=0.125}, draw=none, from=1-1, to=3-3]
            \arrow[from=1-3, to=3-3]
            \arrow[from=3-1, to=3-3]
        \end{tikzcd}
    \end{equation} 
\end{definition}


\begin{remark}
    Unpacking definition \Cref{ch4:: def: categories of basics} we note the analogies with the smooth case. Instead of a space $B$ and a continuous map $B \longra \BO(d)$ we now have an $(\oo,1)$-category $\cB$ and a right fibration $\Bsc$. If $\cB = B$ is an $\oo$-groupoid, then a right fibration $B \longra \Bsc$ must factor through an $\oo$-subgroupoid of $\Bsc$. Inspecting such a subgroupoid, we observe that it must be a copy of $\BO(d)$ for some $d$, thereby demonstrating that it is a stratified augmentation of the smooth unstratified case.

    A conically smooth stratified space $M$ with a $\cB$-manifold structure amounts to a morphism — a natural transformation to be more precise — in $\rfib_{\Bsc}$ with source and target $\entr(M) \longra \Bsc$ and $\cB \longra \Bsc$, respectively
    \begin{equation}
        \begin{tikzcd}[ampersand replacement=\&]
            \&\& \cB \\
            \\
            {\entr(M)} \& {} \& \Bsc
            \arrow[between={0.5}{0.9}, Rightarrow, from=1-3, to=3-2]
            \arrow[from=1-3, to=3-3]
            \arrow["g", dashed, from=3-1, to=1-3]
            \arrow[from=3-1, to=3-3]
        \end{tikzcd}
    \end{equation}
\end{remark}

\begin{observation}
As we will see, the notion of $\cB$-manifold incubates not just the more typical notion of tangential structure, like orientations and spin structures, but also \emph{singularity types} — the latter being a specification of the kinds of links allowed for any given stratum.
\end{observation}

Specifying an $(\oo,1)$-category of basics can be split into two sequential steps:

\begin{enumerate}
    \item Specifying \emph{singularity types} which amounts to specifying a fully-faithful functor 
    \begin{equation}
        \cL \xlonghookra{\, \,\,\,} \Bsc.
    \end{equation}
    \item Specifying an essentially sujective functor 
    \begin{equation}
        \cK \longra \cL.
    \end{equation}
\end{enumerate}
This sequence is sustained by the fact a right fibration (any functor between $(\oo,1)$-categories, in fact) (see \cite[Lemma~2.8.3]{Run25}) admits a factorization into a fully-faithful functor composed with an essentially surjective functor.

Let us begin with singularity types.
\begin{definition}
    Let $\cC$ be an $(\oo,1)$-category and $\cL \subseteq \cC$ an $(\oo,1)$-subcategory. $\cL$ is a \define{sieve} if $\cL \hookra \cC$ is a right fibration. This is equivalent to the following checkable condition:
    \begin{itemize}
        \item \textbf{(Sieve condition)}] \label{con: sieve condition} For each $c \in \cL$, the space 
        \begin{equation}
            \Hom_{\cC}(c',c) \neq \emptyset
        \end{equation}
        implies that $c' \in \cL$.
    \end{itemize}
\end{definition}

\begin{remark}
    Checking a $\mathsf{Kan}$-enriched full subcategory is a right fibration can be done by checking the sieve condition.  
\end{remark}

\begin{example}
    If $\cB = \sD_d$, then $\cD_d \longra \Bsc$ is simply a sieve, hence a right fibration. A $\sD_d$-manifold is a smooth $d$-dimensional manifold — we can interpret this as saying that only a single ``singularity type'' is available, the smooth $d$-dimensional one, or the absence of a singularity.
\end{example}

\begin{example}
    The full subcategory $\Bsc_{\leq k,n} \subset \Bsc$ is a sieve. More concretely, given a basic $U$ of depth $l \leq k$ and pure dimension $n$, then any basic $V$ admitting an open embedding $V \hookra U$ must be of depth $l'$ and pure dimension and $l'\leq l \leq k$. 
\end{example}

\begin{example}
    We now present an example of a subcategory which is not a sieve. Let $\Bsc_1$ be the subcategory $\Bsc_1 \subset \Bsc$  dimension 1 basics. Considering the full subacategory spanned by $\cL = \{\sC(\{n\})\}_{n \geq 1}$, we fail to construct a sieve since 
    \begin{equation}
        \Hom_{\Bsc_1}(\RR, \sC(\{n\})) \neq \emptyset
    \end{equation}
    and $\RR \notin \cL$, as an example
    \begin{equation}
        \tikzfig{figures/Strat/STR38}{\begin{tikzpicture}
	\begin{pgfonlayer}{nodelayer}
		\node [style=0 strat] (0) at (0, 0) {};
		\node [style=none] (1) at (0, 2.5) {};
		\node [style=none] (2) at (-2, -2) {};
		\node [style=none] (3) at (2, -2) {};
		\node [style=none] (4) at (0, 1.75) {};
		\node [style=none] (5) at (0, 0.75) {};
	\end{pgfonlayer}
	\begin{pgfonlayer}{edgelayer}
		\draw (0) to (1.center);
		\draw (2.center) to (0);
		\draw (3.center) to (0);
		\draw [style={open_int_blue}] (4.center) to (5.center);
	\end{pgfonlayer}
\end{tikzpicture}
}
    \end{equation}
    we indicate in blue the image of a conically smooth open embedding $\RR \hookra \sC({1,2,3})$.
\end{example}

\begin{example}
    A typical example considered in mathematical physics and knot theory is that of a single $d$-dimensional proper submanifold $L\subset M$ of an $n$-dimensional smooth manifold $M$. We would like to consider $L$ as a singularity locus, so we would like it to form a class of $\cB$-manifolds for some $\cB$. To do so, we focus on the basic neighborhoods covering points in $L\subset M$ and $M$, respectively. 
    
    Let us define $\sD_{d\subset n}$ by choosing only two objects of $\Bsc$, namely:
    \begin{equation}
        \mathsf{Ob}(\sD_{d\subset n}) = \{\RR^{n}, \RR^{n-d}\times \sC(S^{n-d-1})\}.
    \end{equation}
    The morphism Kan complexes are as follows:
    \begin{itemize}
        \item $\sD_{d\subset n}(\RR^{n},\RR^{n}) = \mathsf{Emb}(\RR^{n},\RR^{n})$.
        \item $\sD_{d\subset n}(\RR^{n},\RR^{n-d}\times \sC(S^{n-d-1})) = \mathsf{Emb}(\RR^{n},\RR^{n}\setminus \RR^{d}) \simeq \sO(n) \times S^{n-d-1}$.
        \item $\sD_{d\subset n}(\RR^{n-d}\times \sC(S^{n-d-1}),\RR^{n}) = \emptyset.$
        \item $\sD_{d\subset n}(\RR^{n-d}\times \sC(S^{n-d-1}),\RR^{n-d}\times \sC(S^{n-d-1}) \, \simeq \, \sO(d) \times \Aut(S^{n-d-1})$.
    \end{itemize}
    A $\sD_{d\subset n}$-manifold is thus the data of an smooth $n$-manifold together with an embedded proper $d$-submanifold $L \subset M$. 
\end{example}

Finally, we look into the specifiation of tangential structures. 

\begin{example}
Let us define $\sD_{n}^{\fr} = (\sD_{n})_{/\RR^{n}} \longra \sD_{n}$. This functor unstraightens to $\sD_{n}^{\op} \longra \mathsf{Kan}$ given by $\sD_{n}(-,\RR^{n})$ — a fiber over $\RR^{n}$ is therefore equivalent to a point of $\mathsf{Emb}(\RR^n, \RR^n)$ or a framing on $\RR^{n}$. A $\sD_n^{\sf fr}$-manifold is therefore a smooth $n$-manifold together with a framing or a trivialization of its tangent bundle. 
\end{example}

\begin{example}
    We define $\sD_{d\subset n}^{\fr}$ as the pullback $(\oo,1)$-category 
    \begin{equation}
        \begin{tikzcd}[ampersand replacement=\&]
            {\sD_{d\subset n}^{\fr}} \&\& {\sD_{ n}^{\fr}} \\
            \\
            {\sD_{d\subset n}} \&\& {\sD_{n}}
            \arrow[from=1-1, to=1-3]
            \arrow[from=1-1, to=3-1]
            \arrow["\lrcorner"{anchor=center, pos=0.125}, draw=none, from=1-1, to=3-3]
            \arrow[from=1-3, to=3-3]
            \arrow[from=3-1, to=3-3]
        \end{tikzcd}
    \end{equation}
    For further unpacking of this construction see \cite[Example~5.2.12]{AFT1}. A $\sD_{d\subset n}^{\fr}$-manifold is comprised of:
    \begin{enumerate}
        \item A framed $n$-manifold $M$.
        \item A properly embedded smooth $d$-submanifold $L$ equipped with a null-homotopy of the map
        \begin{equation}
            L \longra \mathsf{Gr}_{d}(\RR^n)
        \end{equation}
        classifying the tangent bundle as its embedded in $\sT M$ — also called its \emph{Gauss map}. To specify such a null-homotopy one can stipulate a trivialization of both the tangent and normal bundle of $L$ in a way compatible with the underlying trivialization of the tangent bundle of $M$. 
    \end{enumerate}

    As a simple example of a $\sD_{1\subset 2}^{\fr}$ one might consider $\RR^{2}$ with the standard framing and the $x$-axis $L = \{y = 0\}$ together with its standard framing (orientation, equivalently). A non-example in a similar vein is given by $L = S^1 \hookra \RR^2$. Note that the Grassmannian of $1$-planes in $\RR^{2}$ is homeomorphic to $S^1$. Therefore, the Gauss map $L = S^1 \xra{\,\,l\,\,} \mathsf{Gr}_{1}(\RR^2)\simeq S^1$ has a degree and it is $+1$ — it winds once around the origin. As such, it is a standard algebraic topology fact that $l$ does not admit a null-homotopy.  
\end{example}



\subsubsection{Conically smooth stratified spaces with corners}

Now we turn to the goal of incorporating boundaries and corners into the class of conically smooth stratified spaces. To do so, we begin by introducing a special poset, designed with the purpose of capturing the local models of manifolds with $\lag n \rag$-corners \emph{à la} Jänich--Joyce--Laures--Melrose (see \cite{Jan68}\cite{Joy10}\cite{Lau00}\cite{Mel96} for definitions and key results and details). We will follow \cite[Section~7.2]{AFT1}. 

\begin{notation}
    Let $I$ be a finite set. We denote $\cP(I)$ as the poset of subsets of $I$, with the caveat that the ordering is given by reverse inclusion, i.e.\ given a pair of subsets $U, U' \subseteq I$, we have $U\leq U'$ if $U' \subset U$.    
\end{notation}

\begin{example}
    Let $I= \{1,\ldots,n\}$, then we can define a stratification 
    \begin{equation}
        (\RR_{\geq 0})^{n} \xra{\;\;q\;\;} \cP(I), \quad (x_{i})_{i \in I} \longmapsto \{j\,:\, x_{j} = 0\} \subseteq I. 
    \end{equation}
    For $n=1$ and $n=2$ we have the following depictions of this stratification:
    \smallskip
    \begin{equation}
        \tikzfig{figures/Strat/STR20}{\begin{tikzpicture}
	\begin{pgfonlayer}{nodelayer}
		\node [style=none] (0) at (0, 0) {};
		\node [style=none] (2) at (5, 0) {};
		\node [style=none] (6) at (5, 0) {};
		\node [style=none] (5) at (0, 0) {};
		\node [style=0 strat] (7) at (0, 0) {};
		\node [style=none] (9) at (-1, -0.75) {$\{1\}$};
		\node [style=none] (11) at (2.75, -0.75) {$\emptyset$};
	\end{pgfonlayer}
	\begin{pgfonlayer}{edgelayer}
		\draw [style=defect patch b] (5.center) to (6.center);
		\draw [style=cd-ar] (0.center) to (2.center);
	\end{pgfonlayer}
\end{tikzpicture}
} \qquad \qquad \tikzfig{figures/Strat/STR21}{\begin{tikzpicture}
	\begin{pgfonlayer}{nodelayer}
		\node [style=none] (0) at (0, 0) {};
		\node [style=none] (3) at (5, 5) {};
		\node [style=none] (4) at (0, 5) {};
		\node [style=none] (1) at (0, 5) {};
		\node [style=none] (2) at (5, 0) {};
		\node [style=none] (6) at (5, 0) {};
		\node [style=none] (5) at (0, 0) {};
		\node [style=0 strat] (7) at (0, 0) {};
		\node [style=none] (9) at (-1, -0.75) {$\{1,2\}$};
		\node [style=none] (10) at (-1, 2.5) {$\{1\}$};
		\node [style=none] (11) at (2.75, -0.75) {$\{2\}$};
		\node [style=none] (12) at (2.75, 2.5) {$\emptyset$};
	\end{pgfonlayer}
	\begin{pgfonlayer}{edgelayer}
		\draw [style=defect patch b] (6.center)
			 to (3.center)
			 to (4.center)
			 to (5.center)
			 to cycle;
		\draw [style=cd-ar] (0.center) to (1.center);
		\draw [style=cd-ar] (0.center) to (2.center);
	\end{pgfonlayer}
\end{tikzpicture}
}
    \end{equation}
    \medskip

    \nid Note that in the $n=2$ case, the integers which label the stratum describe which coordinates are vanishing, so they change contravariantly to their standard ordering — that is why $\cP(I)$ is ordered with the ``reverse'' inclusion. 
\end{example}

\begin{observation}
    If $I = \{1\}$, then $\cP(I) = \{\emptyset, \{1\}\}$.
\end{observation}

\begin{definition}[Topologically coCartesian]\cite[Definition 7.2.2]{AFT1}\label{ch4:: def: top coCartesian}
    Let $X = (X \ra P)$ be a conically smooth stratified space. We define a continuous map 
    \begin{equation}
        X \xra{c} \cP(I)
    \end{equation}
    to be \define{topologically coCartesian} if it factors as a map of posets as follows:
    \begin{equation}
        \begin{tikzcd}[ampersand replacement=\&]
            X \&\& P \\
            \& {\cP(I)}
            \arrow[from=1-1, to=1-3]
            \arrow["c"', from=1-1, to=2-2]
            \arrow["{\overline{c}}", dashed, from=1-3, to=2-2]
        \end{tikzcd}
    \end{equation} 
    and if, for every $T \subset I$ i.e.\ for every $T \in \cP(I)$, the inclusion of the $T$-stratum $ X_T \xlonghookra X$ into $X$ extends to a conically smooth open embedding over the poset inclusion:
    \begin{equation}
        \begin{tikzcd}[ampersand replacement=\&]
            {(\RR_{\geq 0})^{T} \times X_T} \&\& X \\
            \\
            {\cP(T)} \&\& {\cP(I)}
            \arrow[hook, from=1-1, to=1-3]
            \arrow["{q\times c_{|\cP(T)}}"', from=1-1, to=3-1]
            \arrow["c", from=1-3, to=3-3]
            \arrow[hook, from=3-1, to=3-3]
        \end{tikzcd}
    \end{equation}
\end{definition} 

\begin{example}
    If $M \longra *$ is a smooth manifold and $I = \{*\}$, then a topologically coCartesian map
    \begin{equation}
        M \xra{\;\; c \;\;} \{\emptyset, \{*\}\}
    \end{equation}  
    determines a pair of subspaces $M_{\emptyset}$ and $M_{*}$ together with a (conically) smooth map 
    \begin{equation}
        (\RR_{\geq 0}) \times M_{*} \xra{\;\;\;} M
    \end{equation}
    such that $\{0\}\times M_\ast \xra{\;\;\;} M$ is the stratum inclusion. 
\end{example}

\begin{example}
    Let the following diagram
    \smallskip 
    \begin{equation}
        M \;\; = \;\; \tikzfig{figures/Strat/STR22}{\begin{tikzpicture}
	\begin{pgfonlayer}{nodelayer}
		\node [style=none] (0) at (-5, 5) {};
		\node [style=none] (1) at (5, 5) {};
		\node [style=none] (2) at (-5, -5) {};
		\node [style=none] (3) at (5, -5) {};
		\node [style=dot] (4) at (-5, -5) {};
		\node [style=dot] (5) at (-5, 5) {};
		\node [style=dot] (6) at (5, 5) {};
		\node [style=dot] (7) at (5, -5) {};
		\node [style=none] (9) at (-1.5, 1) {};
		\node [style=none] (10) at (1.5, 1) {};
		\node [style=none] (11) at (-1.5, 5) {};
		\node [style=none] (12) at (1.5, 5) {};
		\node [style=0 strat] (13) at (-1.5, 5) {};
		\node [style=0 strat] (14) at (1.5, 5) {};
	\end{pgfonlayer}
	\begin{pgfonlayer}{edgelayer}
		\draw [style=defect patch b] (0.center)
			 to (2.center)
			 to (3.center)
			 to (1.center)
			 to cycle;
		\draw [style={open_int_blue}] (0.center) to (2.center);
		\draw [style={open_int_blue}] (2.center) to (3.center);
		\draw [style={open_int_blue}] (3.center) to (1.center);
		\draw [style={open_int_blue}] (1.center) to (0.center);
		\draw [style={open_int}, in=-90, out=-90, looseness=1.75] (9.center) to (10.center);
		\draw [style={open_int}] (11.center) to (9.center);
		\draw [style={open_int}] (12.center) to (10.center);
	\end{pgfonlayer}
\end{tikzpicture}
}
    \end{equation}
    \medskip

    \nid denote a stratification $M$ of the unit square $[0,1]^2$. This decomposes as:
    \smallskip
    \begin{equation}
        M_{\emptyset}\;\; = \;\; \tikzfig{figures/Strat/STR23}{\begin{tikzpicture}
	\begin{pgfonlayer}{nodelayer}
		\node [style=none] (0) at (-4, 4) {};
		\node [style=none] (1) at (4, 4) {};
		\node [style=none] (2) at (-4, -4) {};
		\node [style=none] (3) at (4, -4) {};
		\node [style=none] (9) at (-1.5, 0) {};
		\node [style=none] (10) at (1.5, 0) {};
		\node [style=0 strat white] (12) at (-1.5, 4) {};
		\node [style=0 strat white] (13) at (1.5, 4) {};
		\node [style=none] (14) at (1.5, 0) {};
	\end{pgfonlayer}
	\begin{pgfonlayer}{edgelayer}
		\draw [style=defect patch b] (0.center)
			 to (2.center)
			 to (3.center)
			 to (1.center)
			 to cycle;
		\draw [style={open_int}, in=-90, out=-90, looseness=1.75] (9.center) to (10.center);
		\draw [style={open_int_blue_2_dash}] (0.center) to (2.center);
		\draw [style={open_int_blue_2_dash}] (2.center) to (3.center);
		\draw [style={open_int_blue_2_dash}] (1.center) to (3.center);
		\draw [style={open_int_blue_2_dash}] (1.center) to (0.center);
		\draw [style={open_int}] (12) to (9.center);
		\draw [style={open_int}] (13) to (14.center);
	\end{pgfonlayer}
\end{tikzpicture}
}, \qquad \qquad  M \setminus M_{\emptyset} \;\; = \;\; \tikzfig{figures/Strat/STR24}{\begin{tikzpicture}
	\begin{pgfonlayer}{nodelayer}
		\node [style=dot] (4) at (-4, -4) {};
		\node [style=dot] (5) at (-4, 4) {};
		\node [style=dot] (6) at (4, 4) {};
		\node [style=dot] (7) at (4, -4) {};
		\node [style=dot] (8) at (-1.5, 4) {};
		\node [style=dot] (9) at (1.5, 4) {};
	\end{pgfonlayer}
	\begin{pgfonlayer}{edgelayer}
		\draw [style={open_int_blue}] (5) to (4);
		\draw [style={open_int_blue}] (4) to (7);
		\draw [style={open_int_blue}] (7) to (6);
		\draw [style={open_int_blue}] (6) to (5);
	\end{pgfonlayer}
\end{tikzpicture}
}
    \end{equation}
    \medskip
    with $M\setminus M_{\emptyset}$ further decomposing in the obvious way. 
\end{example}


\begin{definition}
    Fix a finite set $I$. The category $\Bsc^{\lag I \rag}$ is defined as follows:
    \begin{enumerate}
        \item \textbf{Objects:} A basic $U\in \Bsc$ endowed with a topologically coCartesian map $U \xra{\; c \;} \cP(I)$. 
        \item \textbf{Morphisms:} Given a pair of objects $\left(c \maps U \ra P \xra{\overline{c}} \cP(I)\right)$ and $\left(d \maps V \ra Q \xra{\overline{d}} \cP(I)\right)$, the set of morphisms is defined as the pullback:
        \begin{equation}
             \begin{tikzcd}[ampersand replacement=\&]
                {\Bsc^{\lag I \rag}\left(\left(U,c\right),\left(V,d\right)\right)} \&\& {\Bsc(U,V)} \\
                \\
                {\Hom_{\Top_{/\cP(I)}}(P,Q)} \&\& {\Hom_{\Top}(P,Q).}
                \arrow[from=1-1, to=1-3]
                \arrow[from=1-1, to=3-1]
                \arrow["\lrcorner"{anchor=center, pos=0.125}, draw=none, from=1-1, to=3-3]
                \arrow[from=1-3, to=3-3]
                \arrow[from=3-1, to=3-3]
            \end{tikzcd}
        \end{equation}
        Since the lower horizontal map is injective so then the top horizontal must be an inclusion on $\pi_0$. In particular, this  implies that $\Bsc^{\lag I \rag}$ is a sieve. We can thus define \begin{equation}
            \Snglr^{\lag I \rag} \, \overset{\sf def}{=} \, \Mfld(\Bsc^{\lag I \rag})
        \end{equation}
        which we will call the $(\oo,1)$\define{-category of conically smooth stratified spaces with} $\lag I \rag$\define{-corners}.   
    \end{enumerate}
\end{definition}


\subsubsection{Unzipping and vector fields}

We now turn to a key construction, \emph{unzipping}, which offers a bridge between the world of conically smooth stratified spaces and that of smooth manifolds with corners, thereby opening up a fruitful avenue for translating statements, definitions, and theorems in one into definitions, statements, and theorems in the other.   

In particular, unzipping will lead us to the definition of parallel and inward-pointing vector fields. These are the correct notions of conically smooth infinitesimal automorphisms for conically smooth stratified spaces with $\lag I \rag$-corners. In turn, these lead to vector field flows and ultimately to a generalization of the Ehresmann fibration theorem (see \Cref{ap:: prop: Ehresmann fibration no boundary}) for conically smooth stratified spaces.

In a slogan: unzipping associates to every conically smooth stratified spaces $(M \ra P)$ of interior depth $k$ a smooth manifold $\unzip(M)$ with $\lag k \rag$-corners (no singularities) and it does so functorially. This procedure works by morally mirroring the notion of algebraic blow-up for algebraic varieties: it removes the singularity locus and glues in its associated spherical bundle. A simple example is shown below:
\begin{equation}
    \tikzfig{figures/Strat/STR39}{\begin{tikzpicture}
	\begin{pgfonlayer}{nodelayer}
		\node [style=none] (0) at (-3.5, 0) {};
		\node [style=none] (1) at (3.5, 0) {};
		\node [style=0 strat] (2) at (0, 0) {};
	\end{pgfonlayer}
	\begin{pgfonlayer}{edgelayer}
		\draw [style=defect patch b] (1.center)
			 to [in=270, out=-90, looseness=1.75] (0.center)
			 to [in=90, out=90, looseness=1.75] cycle;
		\draw [style=ts-dashed-grey, in=90, out=90, looseness=1.75] (0.center) to (1.center);
		\draw [style=ts-dashed-grey, in=270, out=270, looseness=1.75] (1.center) to (0.center);
	\end{pgfonlayer}
\end{tikzpicture}
} \;\;\; \xra{\,\, \unzip \,\,} \;\;\; \tikzfig{figures/Strat/STR40}{\begin{tikzpicture}
	\begin{pgfonlayer}{nodelayer}
		\node [style=none] (0) at (-3.5, 0) {};
		\node [style=none] (1) at (3.5, 0) {};
		\node [style=none] (2) at (-1, 0) {};
		\node [style=none] (3) at (1, 0) {};
	\end{pgfonlayer}
	\begin{pgfonlayer}{edgelayer}
		\draw [style=defect patch b] (1.center)
			 to [in=270, out=-90, looseness=1.75] (0.center)
			 to [in=90, out=90, looseness=1.75] cycle;
		\draw [style=ts-dashed-grey, in=90, out=90, looseness=1.75] (0.center) to (1.center);
		\draw [style=ts-dashed-grey, in=270, out=270, looseness=1.75] (1.center) to (0.center);
		\draw [style=white patch] (2.center)
			 to [in=90, out=90, looseness=1.75] (3.center)
			 to [in=270, out=-90, looseness=1.75] cycle;
		\draw [style=defect line 3, in=450, out=90, looseness=1.75] (3.center) to (2.center);
		\draw [style=defect line 3, in=270, out=270, looseness=1.75] (3.center) to (2.center);
	\end{pgfonlayer}
\end{tikzpicture}
}
\end{equation}
We remove the only singular locus, the 0-dimensional stratum at the origin, and we replace it using the spherical bundle associated to its normal bundle.

\begin{definition}[Unzip at depth $k$]\cite[Definition~7.3.2]{AFT1}
    Fix $k \geq 0$ and let $\overline{X} = (X,c_X)$ be a conically smooth stratified space with $\lag I \rag$-corners, whose interior depth is $k$. We define the following category $\mathsf{Res}(\overline{X},k)$
    \begin{itemize}
        \item \textbf{Objects:} A conically smooth stratified space with $\lag * \sqcup I \rag$-corners 
        \begin{equation}
            \overline{Y} \xra{c_{Y}} \cP(*\sqcup I)
        \end{equation} 
        together with a conically smooth map 
        \begin{equation}
            p \maps \overline{Y} \longra X
        \end{equation}
        such that:
        \begin{itemize}
            \item The following diagram commutes:
            \begin{equation}
                \begin{tikzcd}[ampersand replacement=\&]
                    {\overline{Y}} \&\& X \\
                    \\
                    {\cP(*)\times \cP(I)} \&\& {\cP(I)}
                    \arrow[from=1-1, to=1-3]
                    \arrow[from=1-1, to=3-1]
                    \arrow[from=1-3, to=3-3]
                    \arrow[from=3-1, to=3-3]
                \end{tikzcd}
            \end{equation}
            where we have used the identification $\cP(*\sqcup I) \simeq \cP(*) \times \cP(I)$.
            \item The subspaces of $Y$ \[\overline{\partial}_* \overline{Y} = p^{-1}\overline{X}_{k}\]
            agree.
        \end{itemize}
        \item A morphism from $(Y,c_{Y},p)$ to $(Y',c_{Y'},p')$ is a conically smooth map 
        \begin{equation}
            \begin{tikzcd}[ampersand replacement=\&]
                {\overline{Y}} \&\& {\overline{Y}'} \\
                \\
                \& {\cP(*\sqcup I)\times X}
                \arrow[from=1-1, to=1-3]
                \arrow[from=1-1, to=3-2]
                \arrow[from=1-3, to=3-2]
            \end{tikzcd}
        \end{equation} 
    \end{itemize}
    We define a $k$-unzip of $\overline{X}$ to be 
    \begin{equation}
        \pi_k \maps \unzip_{k}(X) \longra X 
    \end{equation}
    a \emph{final} object of $\mathsf{Res}(\overline{X},k)$. If we restrict $\unzip_k(X)$ to the boundary $\overline{\partial}_{*}\unzip_{k}(X)$ we get the \define{link}:
    \begin{equation}
        \link_{k}(X) \xra{\, \pi_k \,} X_{k}.
    \end{equation} 
\end{definition}

\begin{example}
    Let us explore unzipping through a simple example. Let 
    \begin{equation}
        M \;\; = \;\; \tikzfig{figures/Strat/STR41}{\begin{tikzpicture}
	\begin{pgfonlayer}{nodelayer}
		\node [style=none] (0) at (-3, 3) {};
		\node [style=none] (1) at (-3, -3) {};
		\node [style=none] (2) at (3, 3) {};
		\node [style=none] (3) at (3, -3) {};
		\node [style=0 strat] (4) at (0, 0) {};
		\node [style=none] (5) at (0, 3) {};
		\node [style=none] (6) at (0, -3) {};
	\end{pgfonlayer}
	\begin{pgfonlayer}{edgelayer}
		\draw [style=defect patch b] (1.center)
			 to (3.center)
			 to (2.center)
			 to (0.center);
		\draw [style=ts-dashed-grey] (0.center) to (1.center);
		\draw [style=ts-dashed-grey] (1.center) to (3.center);
		\draw [style=ts-dashed-grey] (3.center) to (2.center);
		\draw [style=ts-dashed-grey] (2.center) to (0.center);
		\draw [style={open_int}] (5.center) to (6.center);
	\end{pgfonlayer}
\end{tikzpicture}
}
    \end{equation}
    \medskip
    
    \nid be $\RR^2$ stratified by the origin and the $y$-axis minus the origin. The 0-dimensional stratum has depth 2, therefore we have that 
    \begin{equation}
        \unzip_2(M) \;\; = \;\; \tikzfig{figures/Strat/STR42}{\begin{tikzpicture}
	\begin{pgfonlayer}{nodelayer}
		\node [style=none] (0) at (-3, 3) {};
		\node [style=none] (1) at (-3, -3) {};
		\node [style=none] (2) at (3, 3) {};
		\node [style=none] (3) at (3, -3) {};
		\node [style=none] (5) at (0, 3) {};
		\node [style=none] (6) at (0, -3) {};
		\node [style=none] (7) at (-0.75, 0) {};
		\node [style=none] (8) at (0.75, 0) {};
		\node [style=0 strat] (9) at (0, 0.75) {};
		\node [style=0 strat] (10) at (0, -0.75) {};
	\end{pgfonlayer}
	\begin{pgfonlayer}{edgelayer}
		\draw [style=defect patch b] (1.center)
			 to (3.center)
			 to (2.center)
			 to (0.center);
		\draw [style=ts-dashed-grey] (0.center) to (1.center);
		\draw [style=ts-dashed-grey] (1.center) to (3.center);
		\draw [style=ts-dashed-grey] (3.center) to (2.center);
		\draw [style=ts-dashed-grey] (2.center) to (0.center);
		\draw [style=white patch] (8.center)
			 to [in=-90, out=-90, looseness=1.75] (7.center)
			 to [in=90, out=90, looseness=1.75] cycle;
		\draw [style={open_int}] (5.center) to (9);
		\draw [style={open_int}] (10) to (6.center);
		\draw [style={open_int_blue}, in=180, out=-180, looseness=1.50] (9) to (10);
		\draw [style={open_int_blue}, in=0, out=0, looseness=1.50] (10) to (9);
	\end{pgfonlayer}
\end{tikzpicture}
}
    \end{equation}
    \medskip

    \nid The resulting conically smooth stratified space $\unzip_2(M)$ is a conically smooth stratified space with $\lag * \rag$-corners, i.e.\ a boundary — the link is the blue boundary circle stratified with the 0-dimensional dots. The interior depth of $\unzip_2(M)$ is 1, so we can apply $\unzip_1$, and this would remove the red stratum and replace it by the sphere bundle associated to its normal bundle 
    \begin{equation}
        \unzip_{1}\left(\unzip_{2}(M)\right) \;\; = \;\; \tikzfig{figures/Strat/STR43}{\begin{tikzpicture}
	\begin{pgfonlayer}{nodelayer}
		\node [style=none] (0) at (-3, 3) {};
		\node [style=none] (1) at (-3, -3) {};
		\node [style=none] (2) at (-0.5, -3) {};
		\node [style=none] (3) at (-0.5, -0.75) {};
		\node [style=none] (4) at (-0.5, 0.75) {};
		\node [style=none] (5) at (-0.5, 3) {};
		\node [style=none] (6) at (3, -3) {};
		\node [style=none] (7) at (3, 3) {};
		\node [style=none] (8) at (0.5, 3) {};
		\node [style=none] (9) at (0.5, 0.75) {};
		\node [style=none] (10) at (0.5, -0.75) {};
		\node [style=none] (11) at (0.5, -3) {};
		\node [style=0 strat] (12) at (-0.5, -0.75) {};
		\node [style=0 strat] (13) at (-0.5, 0.75) {};
		\node [style=0 strat] (14) at (0.5, 0.75) {};
		\node [style=0 strat] (15) at (0.5, -0.75) {};
		\node [style=none] (16) at (0.5, 3) {};
	\end{pgfonlayer}
	\begin{pgfonlayer}{edgelayer}
		\draw [style=defect patch b] (5.center)
			 to (4.center)
			 to [bend left=270, looseness=1.75] (3.center)
			 to (2.center)
			 to (1.center)
			 to (0.center)
			 to cycle;
		\draw [style=defect patch b] (11.center)
			 to (10.center)
			 to [bend left=270, looseness=1.75] (9.center)
			 to (8.center)
			 to (7.center)
			 to (6.center)
			 to cycle;
		\draw [style=ts-dashed-grey] (2.center) to (1.center);
		\draw [style=ts-dashed-grey] (1.center) to (0.center);
		\draw [style=ts-dashed-grey] (0.center) to (5.center);
		\draw [style=ts-dashed-grey] (16.center) to (7.center);
		\draw [style=ts-dashed-grey] (7.center) to (6.center);
		\draw [style=ts-dashed-grey] (6.center) to (11.center);
		\draw [style={open_int_blue}] (5.center) to (13);
		\draw [style={open_int_blue}] (16.center) to (14);
		\draw [style={open_int_blue}] (15) to (11.center);
		\draw [style={open_int_blue}] (12) to (2.center);
		\draw [style={open_int_blue}, bend left=90, looseness=1.50] (12) to (13);
		\draw [style={open_int_blue}, bend left=90, looseness=1.50] (14) to (15);
	\end{pgfonlayer}
\end{tikzpicture}
}
    \end{equation} 
\end{example}

\begin{remark}\cite[Example~7.3.4]{AFT1}
    Let us consider $\unzip_2(M)$ from the example above. A possible object of $\mathsf{Res}\left(\unzip_2(M),1\right)$ is a curve $\gamma \maps (0,1] \longra \unzip_2(M)$ which must be transverse to $W$ (the depth $1$ stratum) at $1$ and only at $1$ — this ensures it is a manifold with boundary in $\unzip_2(M)$ and its boundary is contained in $W$. For example, let the yellow curve $\gamma$ be such an object
    \begin{equation}
        \tikzfig{figures/Strat/STR44}{\begin{tikzpicture}
	\begin{pgfonlayer}{nodelayer}
		\node [style=none] (0) at (-3, 3) {};
		\node [style=none] (1) at (-3, -3) {};
		\node [style=none] (2) at (3, 3) {};
		\node [style=none] (3) at (3, -3) {};
		\node [style=none] (5) at (0, 3) {};
		\node [style=none] (6) at (0, -3) {};
		\node [style=none] (7) at (-0.75, 0) {};
		\node [style=none] (8) at (0.75, 0) {};
		\node [style=0 strat] (9) at (0, 0.75) {};
		\node [style=0 strat] (10) at (0, -0.75) {};
		\node [style=none] (11) at (2, 2) {};
		\node [style=none] (13) at (2, -0.75) {};
		\node [style=none] (14) at (0, -1.5) {};
	\end{pgfonlayer}
	\begin{pgfonlayer}{edgelayer}
		\draw [style=defect patch b] (1.center)
			 to (3.center)
			 to (2.center)
			 to (0.center);
		\draw [style=ts-dashed-grey] (0.center) to (1.center);
		\draw [style=ts-dashed-grey] (1.center) to (3.center);
		\draw [style=ts-dashed-grey] (3.center) to (2.center);
		\draw [style=ts-dashed-grey] (2.center) to (0.center);
		\draw [style=white patch] (8.center)
			 to [in=-90, out=-90, looseness=1.75] (7.center)
			 to [in=90, out=90, looseness=1.75] cycle;
		\draw [style={open_int}] (5.center) to (9);
		\draw [style={open_int}] (10) to (6.center);
		\draw [style={open_int_blue}, in=180, out=-180, looseness=1.50] (9) to (10);
		\draw [style={open_int_blue}, in=0, out=0, looseness=1.50] (10) to (9);
		\draw [style={open_int_green}, in=-15, out=-105, looseness=0.75] (11.center) to (13.center);
		\draw [style={open_int_green}, in=0, out=165] (13.center) to (14.center);
	\end{pgfonlayer}
\end{tikzpicture}
}
    \end{equation}
    \medskip

    \nid Moreover, there is a clear morphism in $\mathsf{Res}\left(\unzip_2(M),1\right)$ from $\gamma$ to $\unzip_2(M)$, described below 
    \begin{equation}
        \tikzfig{figures/Strat/STR44}{\begin{tikzpicture}
	\begin{pgfonlayer}{nodelayer}
		\node [style=none] (0) at (-3, 3) {};
		\node [style=none] (1) at (-3, -3) {};
		\node [style=none] (2) at (3, 3) {};
		\node [style=none] (3) at (3, -3) {};
		\node [style=none] (5) at (0, 3) {};
		\node [style=none] (6) at (0, -3) {};
		\node [style=none] (7) at (-0.75, 0) {};
		\node [style=none] (8) at (0.75, 0) {};
		\node [style=0 strat] (9) at (0, 0.75) {};
		\node [style=0 strat] (10) at (0, -0.75) {};
		\node [style=none] (11) at (2, 2) {};
		\node [style=none] (13) at (2, -0.75) {};
		\node [style=none] (14) at (0, -1.5) {};
	\end{pgfonlayer}
	\begin{pgfonlayer}{edgelayer}
		\draw [style=defect patch b] (1.center)
			 to (3.center)
			 to (2.center)
			 to (0.center);
		\draw [style=ts-dashed-grey] (0.center) to (1.center);
		\draw [style=ts-dashed-grey] (1.center) to (3.center);
		\draw [style=ts-dashed-grey] (3.center) to (2.center);
		\draw [style=ts-dashed-grey] (2.center) to (0.center);
		\draw [style=white patch] (8.center)
			 to [in=-90, out=-90, looseness=1.75] (7.center)
			 to [in=90, out=90, looseness=1.75] cycle;
		\draw [style={open_int}] (5.center) to (9);
		\draw [style={open_int}] (10) to (6.center);
		\draw [style={open_int_blue}, in=180, out=-180, looseness=1.50] (9) to (10);
		\draw [style={open_int_blue}, in=0, out=0, looseness=1.50] (10) to (9);
		\draw [style={open_int_green}, in=-15, out=-105, looseness=0.75] (11.center) to (13.center);
		\draw [style={open_int_green}, in=0, out=165] (13.center) to (14.center);
	\end{pgfonlayer}
\end{tikzpicture}
} \;\; \longra \;\; \tikzfig{figures/Strat/STR45}{\begin{tikzpicture}
	\begin{pgfonlayer}{nodelayer}
		\node [style=none] (0) at (-3.25, 3) {};
		\node [style=none] (1) at (-3.25, -3) {};
		\node [style=none] (2) at (-0.5, -3) {};
		\node [style=none] (3) at (-0.5, -0.75) {};
		\node [style=none] (4) at (-0.5, 0.75) {};
		\node [style=none] (5) at (-0.5, 3) {};
		\node [style=none] (6) at (3.25, -3) {};
		\node [style=none] (7) at (3.25, 3) {};
		\node [style=none] (8) at (0.5, 3) {};
		\node [style=none] (9) at (0.5, 0.75) {};
		\node [style=none] (10) at (0.5, -0.75) {};
		\node [style=none] (11) at (0.5, -3) {};
		\node [style=0 strat] (12) at (-0.5, -0.75) {};
		\node [style=0 strat] (13) at (-0.5, 0.75) {};
		\node [style=0 strat] (14) at (0.5, 0.75) {};
		\node [style=0 strat] (15) at (0.5, -0.75) {};
		\node [style=none] (16) at (0.5, 3) {};
		\node [style=none] (17) at (2.5, 1.75) {};
		\node [style=none] (18) at (2.5, -1) {};
		\node [style=none] (19) at (0.5, -1.75) {};
	\end{pgfonlayer}
	\begin{pgfonlayer}{edgelayer}
		\draw [style=defect patch b] (5.center)
			 to (4.center)
			 to [bend left=270, looseness=1.75] (3.center)
			 to (2.center)
			 to (1.center)
			 to (0.center)
			 to cycle;
		\draw [style=defect patch b] (11.center)
			 to (10.center)
			 to [bend left=270, looseness=1.75] (9.center)
			 to (8.center)
			 to (7.center)
			 to (6.center)
			 to cycle;
		\draw [style=ts-dashed-grey] (2.center) to (1.center);
		\draw [style=ts-dashed-grey] (1.center) to (0.center);
		\draw [style=ts-dashed-grey] (0.center) to (5.center);
		\draw [style=ts-dashed-grey] (16.center) to (7.center);
		\draw [style=ts-dashed-grey] (7.center) to (6.center);
		\draw [style=ts-dashed-grey] (6.center) to (11.center);
		\draw [style={open_int_blue}] (5.center) to (13);
		\draw [style={open_int_blue}] (16.center) to (14);
		\draw [style={open_int_blue}] (15) to (11.center);
		\draw [style={open_int_blue}] (12) to (2.center);
		\draw [style={open_int_blue}, bend left=90, looseness=1.50] (12) to (13);
		\draw [style={open_int_blue}, bend left=90, looseness=1.50] (14) to (15);
		\draw [style={open_int_green}, in=-15, out=-105, looseness=0.75] (17.center) to (18.center);
		\draw [style={open_int_green}, in=0, out=165] (18.center) to (19.center);
	\end{pgfonlayer}
\end{tikzpicture}
}
    \end{equation}
\end{remark}
\medskip

Unzipping allows us to discuss the ``right'' notion of vector fields for conically smooth stratified spaces. Before doing so, we introduce a few key components and useful notation.

Given a smooth manifold $M$, let us denote $\vartheta(M)$ for the vector space $\Gamma(\sT M \ra M)$ of smooth sections or smooth vector fields on $M$. Fixing a smooth map $f\maps M \longra N$, we have the following linear maps
\begin{equation}
    \vartheta(N) \longra \Gamma(f^* \sT N \ra M), \qquad (N \xra{X} \sT N) \longmapsto (M \xra{f} N \xra{X} \sT N)
\end{equation}
which we will call restriction, and 
\begin{equation}
    \vartheta(M) \xra{\sD f} \Gamma(f^* \sT N \ra M), \qquad (M \xra{X} \sT M) \longmapsto (M \xra{X} \sT M \xra{\sD f} \sT N)
\end{equation}
which we will call the derivative map. We will be considering smooth vector fields on smooth manifolds with corners, so let us provide a gentle reminder of how these are defined. Smooth vector fields are defined locally, so it suffices to determine what these mean on open subsets of the form 
\begin{equation}
    \RR^{i} \times \RR^{j}_{\geq 0}, \qquad \vartheta(\RR^{i} \times \RR^{j}_{\geq 0}) = \Gamma\left( \sT \RR^{i+j}_{|\RR^{i}\times \RR^{j}_{\geq 0}} \longra \RR^{i}\times \RR^{j}_{\geq 0} \right).
\end{equation}
In particular, if $M$ is a smooth manifold with boundary, say $\RR_{\geq 0}$, then a possible vector field might point outward. We are interested in defining sub-sheaves of vector fields which point-inward or tangently to the faces of the boundary.

\begin{definition}
    Let $\overline{M}$ be a smooth manifold with corners. We define the sub-sheaf of \define{inward-pointing vector fields} as
    \begin{equation}
        \vartheta^{\sf in}(\overline{M}) \, = \, \{X \maps \RR^{i}\times \RR^{j}_{\geq 0} \ra \RR^{i+j}\,:\,  \mathsf{pr}_{i}\left(X_{|\RR^{i}\times \RR^{S} \times \{0\}}\right) > 0, \,\,\,\, \forall\, i \in \{1,\ldots,j\}\setminus S \,\,\,\, \forall\, S \in \cP(\{1,\ldots,j\})\}
    \end{equation}
    and that of \define{parallel vector fields} as
    \begin{equation}
        \vartheta^{\sf ||}(M) \, = \, \{X \maps \RR^{i}\times \RR^{j}_{\geq 0} \ra \RR^{i+j}\,:\,  X_{|\RR^{i}\times \RR^{S} \times \{0\}} \subset \RR^{i}\times \RR^{S} \times \{0\} \,\,\,\, \forall\, S \in \cP(\{1,\ldots,j\})\}.
    \end{equation} 
\end{definition}

\begin{example}
    Consider the case $i=1$ and $j=1$, i.e.,\ $\RR \times \RR_{\geq 0}$.
    \begin{equation}
        \tikzfig{figures/Strat/STR46}{\begin{tikzpicture}
	\begin{pgfonlayer}{nodelayer}
		\node [style=none] (4) at (-4.5, 5) {};
		\node [style=none] (5) at (-4.5, 0) {};
		\node [style=none] (6) at (4.5, 5) {};
		\node [style=none] (7) at (4.5, 0) {};
		\node [style=none] (0) at (-4.5, 0) {};
		\node [style=none] (1) at (4.5, 0) {};
		\node [style=none] (8) at (0, 0) {};
		\node [style=none] (9) at (0, 6) {};
		\node [style=none] (10) at (-4.5, 0) {};
		\node [style=none] (11) at (1, 6.5) {$\RR^{\{2\}}$};
		\node [style=none] (12) at (5.5, 0) {$\RR^{\{1\}}$};
		\node [style=none] (13) at (-2, 0) {};
		\node [style=none] (14) at (-0.75, 0.5) {};
		\node [style=none] (15) at (1, 0) {};
		\node [style=none] (16) at (2, 0) {};
		\node [style=none] (17) at (2.5, -0.75) {};
	\end{pgfonlayer}
	\begin{pgfonlayer}{edgelayer}
		\draw [style=defect patch b] (5.center)
			 to (4.center)
			 to (6.center)
			 to (7.center)
			 to cycle;
		\draw [style={open_int_blue}] (0.center) to (1.center);
		\draw [style=defect patch b] (0.center) to (1.center);
		\draw [style=dir thick] (8.center) to (9.center);
		\draw [style=cd-ar red] (13.center) to (14.center);
		\draw [style=cd-ar red] (8.center) to (15.center);
		\draw [style=cd-ar red] (16.center) to (17.center);
	\end{pgfonlayer}
\end{tikzpicture}
} 
    \end{equation} 
    \medskip

    \nid The vectors are inward-pointing, parallel, and non-inward-pointing from left to right, respectively.
\end{example}

Lastly, we note that given a conically smooth stratified space $M\ra P$ and a linear subposet $S \subset P$ we can articulate the corresponding face of its unzip:
\begin{equation}
    \partial_{S} \unzip(M) \, = \, \cap_{s \in S} \partial_{s} \unzip(M).
\end{equation}
It comes with a constructible bundle (smooth fiber bundle) $\pi_S \maps \partial_{S} \unzip(M) \longra M_{\mathsf{Min}(S)}$.

\begin{definition}[Parallel vector fields]\cite[Definition~8.1.2]{AFT1}
    Let $M = (M\ra P)$ be a conically smooth stratified space. We define the vector space $\vartheta(M)$ of \define{parallel vector fields} on $M$ as the limit 
    \begin{equation}
        \begin{tikzcd}[ampersand replacement=\&]
            {\vartheta(M)} \&\&\& {\vartheta^{||}\left(\unzip(M)\right)} \\
            \\
            \&\&\& {\prod_{S\underset{\text{linear}}{\subset}P}\vartheta\left(\partial_{S}\unzip(M)\right)} \\
            \\
            {\prod_{p\in P} \vartheta(M_p)} \&\&\& {\prod_{S\underset{\text{linear}}{\subset}P} \Gamma\left(\pi_{S}^*\sT M_{\mathsf{Min}(S)}\ra \partial_{S}\unzip(M)\right)}
            \arrow["{\widetilde{(-)}}", from=1-1, to=1-4]
            \arrow["{(-)_{|M_{-}}}"', from=1-1, to=5-1]
            \arrow["{(-)_{|\partial_{-}}}", from=1-4, to=3-4]
            \arrow["{\sD\pi_{-}}", from=3-4, to=5-4]
            \arrow["{(\pi_{-})^{*}}"', from=5-1, to=5-4]
        \end{tikzcd}
    \end{equation}
\end{definition}


\begin{example}
    The notion of parallel vector field is best understood through an example. Let us consider the following stratification of the open disk
    \begin{equation}
        M \;\; = \;\; \tikzfig{figures/Strat/STR47}{\begin{tikzpicture}
	\begin{pgfonlayer}{nodelayer}
		\node [style=none] (0) at (-3.5, 0) {};
		\node [style=none] (1) at (3.5, 0) {};
		\node [style=0 strat] (2) at (0, 0) {};
		\node [style=none] (3) at (0, 3.675) {};
		\node [style=none] (4) at (-2.25, -2.925) {};
		\node [style=none] (5) at (2.425, -2.725) {};
	\end{pgfonlayer}
	\begin{pgfonlayer}{edgelayer}
		\draw [style=defect patch b] (1.center)
			 to [in=270, out=-90, looseness=1.75] (0.center)
			 to [in=90, out=90, looseness=1.75] cycle;
		\draw [style=ts-dashed-grey, in=90, out=90, looseness=1.75] (0.center) to (1.center);
		\draw [style=ts-dashed-grey, in=270, out=270, looseness=1.75] (1.center) to (0.center);
		\draw [style={open_int}] (3.center) to (2);
		\draw [style={open_int}] (2) to (4.center);
		\draw [style={open_int}] (2) to (5.center);
	\end{pgfonlayer}
\end{tikzpicture}
}.
    \end{equation}
    \medskip
    
    \nid We might change it by both breaking into strata or taking its total unzip:
    \begin{equation}
        \tikzfig{figures/Strat/STR48}{\begin{tikzpicture}
	\begin{pgfonlayer}{nodelayer}
		\node [style=none] (6) at (-0.25, 3.675) {};
		\node [style=none] (7) at (-2.5, -2.825) {};
		\node [style=none] (8) at (-0.25, 0.175) {};
		\node [style=none] (11) at (0.925, 3.675) {};
		\node [style=none] (12) at (3.175, -2.825) {};
		\node [style=none] (13) at (0.925, 0.175) {};
		\node [style=none] (15) at (-1.925, -3.575) {};
		\node [style=none] (16) at (2.325, -3.575) {};
		\node [style=none] (17) at (0.325, -0.575) {};
		\node [style=0 strat] (18) at (0.325, 0) {};
		\node [style=none] (19) at (0.325, 0.5) {};
		\node [style=none] (20) at (0.325, 3.675) {};
		\node [style=none] (21) at (-2.25, -3.25) {};
		\node [style=none] (22) at (0, -0.25) {};
		\node [style=none] (23) at (0.675, -0.275) {};
		\node [style=none] (24) at (2.7, -3.2) {};
	\end{pgfonlayer}
	\begin{pgfonlayer}{edgelayer}
		\draw [style=defect patch b] (8.center)
			 to (6.center)
			 to [bend right=75, looseness=1.25] (7.center)
			 to cycle;
		\draw [style=defect patch b] (13.center)
			 to (11.center)
			 to [bend left=75, looseness=1.25] (12.center)
			 to cycle;
		\draw [style=defect patch b] (17.center)
			 to (15.center)
			 to [bend right=45] (16.center)
			 to cycle;
		\draw [style={open_int}] (20.center) to (19.center);
		\draw [style={open_int}] (21.center) to (22.center);
		\draw [style={open_int}] (23.center) to (24.center);
		\draw [style=ts-dashed-grey] (6.center) to (8.center);
		\draw [style=ts-dashed-grey] (8.center) to (7.center);
		\draw [style=ts-dashed-grey, bend left=75, looseness=1.25] (7.center) to (6.center);
		\draw [style=ts-dashed-grey] (17.center) to (15.center);
		\draw [style=ts-dashed-grey, bend right=45] (15.center) to (16.center);
		\draw [style=ts-dashed-grey] (16.center) to (17.center);
		\draw [style=ts-dashed-grey] (13.center) to (12.center);
		\draw [style=ts-dashed-grey] (13.center) to (11.center);
		\draw [style=ts-dashed-grey, bend left=75, looseness=1.25] (11.center) to (12.center);
	\end{pgfonlayer}
\end{tikzpicture}
} \qquad , \qquad \tikzfig{figures/Strat/STR49}{\begin{tikzpicture}
	\begin{pgfonlayer}{nodelayer}
		\node [style=none] (6) at (-0.25, 3.675) {};
		\node [style=none] (7) at (-2.5, -2.825) {};
		\node [style=none] (8) at (-0.25, 0.175) {};
		\node [style=none] (11) at (0.925, 3.675) {};
		\node [style=none] (12) at (3.175, -2.825) {};
		\node [style=none] (13) at (0.925, 0.175) {};
		\node [style=none] (15) at (-1.925, -3.575) {};
		\node [style=none] (16) at (2.325, -3.575) {};
		\node [style=none] (17) at (0.325, -0.575) {};
		\node [style=none] (21) at (-2.25, -3.25) {};
		\node [style=none] (23) at (0.75, -0.5) {};
		\node [style=none] (24) at (2.7, -3.2) {};
		\node [style=none] (25) at (-1.5, 0) {};
		\node [style=none] (26) at (2, 0) {};
		\node [style=0 strat] (27) at (-1.175, -1.05) {};
		\node [style=0 strat] (28) at (-0.25, 1.5) {};
		\node [style=0 strat] (29) at (0.925, 1.5) {};
		\node [style=0 strat] (30) at (1.725, -0.925) {};
		\node [style=0 strat] (31) at (1.025, -1.675) {};
		\node [style=0 strat] (32) at (-0.5, -1.7) {};
	\end{pgfonlayer}
	\begin{pgfonlayer}{edgelayer}
		\draw [style=defect patch b] (8.center)
			 to (6.center)
			 to [bend right=75, looseness=1.25] (7.center)
			 to cycle;
		\draw [style=defect patch b] (13.center)
			 to (11.center)
			 to [bend left=75, looseness=1.25] (12.center)
			 to cycle;
		\draw [style=defect patch b] (17.center)
			 to (15.center)
			 to [bend right=45] (16.center)
			 to cycle;
		\draw [style=white patch] (26.center)
			 to [in=270, out=-90, looseness=1.75] (25.center)
			 to [in=90, out=90, looseness=1.50] cycle;
		\draw [style={open_int_blue}] (6.center) to (28);
		\draw [style={open_int_blue}] (7.center) to (27);
		\draw [style={open_int_blue}] (15.center) to (32);
		\draw [style={open_int_blue}] (16.center) to (31);
		\draw [style={open_int_blue}] (12.center) to (30);
		\draw [style={open_int_blue}] (11.center) to (29);
		\draw [style={open_int_blue}, in=-165, out=120] (27) to (28);
		\draw [style={open_int_blue}, in=-150, out=-30, looseness=0.50] (32) to (31);
		\draw [style={open_int_blue}, in=-30, out=60] (30) to (29);
	\end{pgfonlayer}
\end{tikzpicture}
}
    \end{equation}
    \medskip 

    \nid Any parallel vector field has to be such that it restricts to a vector field on every stratum; moreover, they must satisfy a compatibility condition stipulated by their lift to $\partial_S \unzip(M)$ for varying $S$.     
\end{example}

\begin{remark}
    Fixing a conically smooth stratified space $M$, the collection of parallel vector fields organizes into a sheaf of vector spaces:
    \begin{equation}
        \vartheta_{M} \maps \mathsf{Opens}(M)^{\op} \longra \mathsf{Vect}_{\RR}, \quad \vartheta_M(U) \, = \, \vartheta(U).
    \end{equation} 
\end{remark}


\begin{proposition}[Flows]\cite[Lemma~8.1.3]{AFT1}\label{ch4:: def:flows}
For a conically smooth stratified space $M = (M \longra P)$, and for $V$ a parallel vector field on $M$, there exists a conically smooth map 
\begin{equation}\label{eq:def:flows}
    \gamma^V \maps M \times \RR \dashrightarrow M
\end{equation}
defined on an open neighborhood of $M\times \{0\}$, satisfying the following conditions:
\begin{enumerate}
    \item For each conically smooth map $\epsilon \maps M \longra \RR$, the composition 
    \begin{equation}\label{eq:def:flows:iso}
        M \xra{~ \id_M \times \epsilon ~} M \times \RR \dashrightarrow{\gamma^V} M
    \end{equation}
    is a conically smooth isomorphism, provided it is defined. 
    \item \label{additivity for flow} There is an equality 
    \begin{equation}\label{eq:def:flows:action}
        \gamma^V_{t}\big(\gamma^V_{s}(x)\big) = \gamma^V_{s+t}(x)
    \end{equation}
    provided the maps are defined.
    \item For each point $x\in M$ belonging to a stratum $x\in M_p \subset M$,
    \begin{equation}\label{eq:def:flows:tangent}
        \frac{\sf d}{{\sf d}t}\gamma_t^{V}(x)|_{t=0} = V_{|M_p}(x).
    \end{equation}
\end{enumerate}
We call $\gamma^V$ the \define{conically smooth flow of }$V$.
\end{proposition}

\begin{definition}[Inward-pointing vector fields]\cite[Definition~8.1.2]{AFT1}
    Let $\overline{M} = (\overline{M}\ra P,c)$ be a conically smooth stratified space with corners. We define the vector space $\vartheta^{\sf in}(M)$ of \define{inward vector fields} on $\overline{M}$ as the limit 
    \begin{equation}
        \begin{tikzcd}[ampersand replacement=\&]
            {\vartheta^{\sf in}(M)} \&\&\& {\vartheta^{||}\left(\unzip(\overline{M})\right)} \\
            \\
            \&\& {\prod_{S\underset{\text{linear}}{\subset}P}\vartheta\left(\partial_{S}\unzip(M)\right)} \& {\prod_{S\underset{\text{linear}}{\subset}P}\Gamma\left(\sT \unzip(\overline{M})_{|\partial_S}\ra \partial_S \right)} \\
            \\
            {\prod_{p\in P} \vartheta^{\sf in}(M_p)} \&\& {\prod_{S\underset{\text{linear}}{\subset}P} \Gamma\left(\pi_{S}^*\sT \overline{M}_{\mathsf{Min}(S)}\ra \partial_{S}\unzip(\overline{M})\right)}
            \arrow["{\widetilde{(-)}}", from=1-1, to=1-4]
            \arrow[from=1-1, to=3-3]
            \arrow["{(-)_{|\overline{M_{-}}}}"', from=1-1, to=5-1]
            \arrow["{(-)_{|\partial_{-}}}", from=1-4, to=3-4]
            \arrow[from=3-3, to=3-4]
            \arrow["{\sD\pi_{-}}", from=3-3, to=5-3]
            \arrow["{(\pi_{-})^{*}}"', from=5-1, to=5-3]
        \end{tikzcd}
    \end{equation}
\end{definition}


\begin{proposition}[Inward flows]\cite[Lemma~8.1.3]{AFT1}\label{ch4:: def: in flows}
For a conically smooth stratified space with corners $\overline{M} = (M \longra P,c)$, and for $V$ an inward vector field on $\overline{M}$, there exists a conically smooth map 
\begin{equation}\label{eq:def:in flows}
    \gamma^V \maps M \times \RR \dashrightarrow M
\end{equation}
defined on an open neighborhood of $M\times \{0\}$, satisfying the following conditions:
\begin{enumerate}
    \item For each conically smooth map $\epsilon \maps \overline{M} \longra \RR$, and for each linear subposet $S\subset P$, the composition 
    \begin{equation}\label{eq:def:in flows:iso}
        \partial_{S}\unzip(\overline{M}) \xra{\id \times \epsilon_{|\partial_{S}}} \partial_{S}\unzip(\overline{M})\times \RR_{\geq 0} \dashrightarrow \unzip(\overline{M})
    \end{equation}
    is a refinement onto its image, provided it is defined. 
    \item There is an equality 
    \begin{equation}\label{eq:def:in flows:action}
        \gamma^V_{t}\big(\gamma^V_{s}(x)\big) = \gamma^V_{s+t}(x)
    \end{equation}
    provided the maps are defined.
    \item For each point $x\in M$ belonging to a stratum $x\in M_p \subset M$,
    \begin{equation}\label{eq:def:in flows:tangent}
        \frac{\sf d}{{\sf d}t}\gamma_t^{V}(x)|_{t=0} = V_{|M_p}(x).
    \end{equation}
    
\end{enumerate}
\end{proposition}

From the existence of flows we can establish the existence of regular neighborhoods for arbitrary proper constructible embedding $Y\hookra X$.

\begin{proposition}\cite[Proposition~8.2.5]{AFT1}
    Let $Y\hookra X$ be a proper constructible embedding. Then there is a conically smooth map 
    \begin{equation}
        \sC(\pi_Y) \longra X
    \end{equation}
    under $Y\hookra X$ such that it is a refinement onto its image, and its image is open. The conically smooth stratified space $\sC(\pi_Y)$ is built as the fiberwise open cone of $\link_{Y}(X) \xra{\pi_Y} Y$.
\end{proposition}

The existence of regular neighborhoods allows us to establish the following important structural fact about $\Strat$.

\begin{proposition}\cite[Lemma~2.5.2]{AFR1}
    Given a conically smooth stratified space $M$ together with its deepest stratum $M_k \subset M$, the following is a pushout square in the $(\oo,1)$-category $\Strat$:
    \begin{equation}
        \begin{tikzcd}[ampersand replacement=\&]
            {\link_{M_{k}}(M)} \&\& {\unzip(M)} \\
            \\
            {M_k} \&\& M
            \arrow[hook, from=1-1, to=1-3]
            \arrow[from=1-1, to=3-1]
            \arrow[from=1-3, to=3-3]
            \arrow[hook, from=3-1, to=3-3]
        \end{tikzcd}
    \end{equation} 
\end{proposition}

Conically smooth stratified spaces can thus be obtained from their unzip, by iteratively gluing strata to the shallower faces. More concretely, assuming $k$ is the interior depth of $M$:
\begin{equation}
    \unzip_{k}(M) = X_{k-1} \coprod_{\link_{X_{k-1}}(\unzip_{k-1}(M))} \unzip_{k-1}(M)
\end{equation}

\subsubsection{Ehresmann fibration theorem}
A key result in building the infrastructure around higher categories of cobordisms is the Ehresmann fibration theorem, in its versions with boundary and without boundary (see \Cref{ap:: prop: Ehresmann fibration no boundary}). We now turn to proving its conically smooth analogue — in essence, a conically smooth version of \emph{Thom's first isotopy lemma} (see Mather's seminal notes \cite[Proposition~11.1]{Mat70}). 

We begin by establishing the fact that, akin to smooth submersions, conically smooth submersions admit lifts of parallel vector fields.

\begin{proposition}\label{prop: lifts weak}
    Let $\pi\maps X \longra K$ be a weakly constructible bundle. Then every parallel vector field $V \in \vartheta(K)$ lifts to a parallel vector field $\widetilde{V}\in \vartheta(X)$, i.e., $\sD \pi (\widetilde{V}) = V$.
\end{proposition}
\begin{proof}
    The map $\pi\maps X \longra K$ determines a morphism between sheaves of vector spaces on $\unzip(X)$:
\begin{equation}
    \sD \tilde{\pi} \colon \vartheta_{\unzip(X)} \ra \tilde{\pi}^{\ast}\vartheta_{\unzip(K)}
\end{equation}
because the differential of a face-preserving map again preserves the corresponding boundary faces. Locally on $\unzip(K)$, the map $\tilde{\pi}$ is isomorphic to a projection $U\times F \longra U$, and for a projection every parallel vector field on $U$ lifts to $(V,0)$ on $U\times F$. Hence $\sD\tilde{\pi}$ is surjective as a map of sheaves. We therefore have the commutative diagram
\begin{equation}
\begin{tikzcd}[ampersand replacement=\&]
	{(q_X)_\ast\vartheta_{\unzip(X)}} \&\& {(q_X)_\ast (\tilde{\pi}^\ast\vartheta_{\unzip(K)})} \\
	\&\& {\pi^\ast (q_K)_{\ast}(\vartheta_{\unzip(K)})} \\
	{\vartheta_X} \&\& {\pi^\ast\vartheta_K}
	\arrow["{(q_X)_\ast(\sD \tilde{\pi})}", from=1-1, to=1-3]
	\arrow[from=1-1, to=3-1]
	\arrow[from=1-3, to=2-3]
	\arrow[from=2-3, to=3-3]
	\arrow[dashed, from=3-1, to=3-3]
\end{tikzcd}
\end{equation}
The vertical maps identify parallel vector fields on $X$ and $K$ with the corresponding pushforwards from their unzips, so the bottom horizontal map is induced from the top one. Since the top map is surjective on sections over every open set, the same is true for $\sD\pi$. Thus every parallel vector field on $K$ lifts locally to $X$, and because $\vartheta_X$ is a sheaf of vector spaces these local lifts glue to a global lift.
\end{proof}

\begin{proposition}\label{prop:lifts of vector fields}
    Let $f\maps X \longra Y$ be a conically smooth submersion, and let $V$ be a parallel vector field on $Y$. There exists a parallel vector field $W$, defined on $X$, such that $\sD f(W) = V.$  
\end{proposition}
\begin{proof}
     Because $f$ is a submersion, we can find a basis $\cU_{Y}= \{\varphi_\alpha(O_\alpha)\}_\alpha$ for the topology of $Y$ together with conically smooth open embeddings $\psi_\alpha \maps U_\alpha \times O_\alpha \hookrightarrow X$ for which 
    \begin{equation}
        \big(f\circ \psi_\alpha \big)(u,o) = \varphi_\alpha(o).
    \end{equation} 
    On the local model $U_\alpha \times O_\alpha \longra O_\alpha$, the restriction $V_\alpha := V_{|O_\alpha}$ lifts to the parallel vector field $(0,V_\alpha)$. Transporting this field across $\psi_\alpha$ gives a parallel lift $W_\alpha$ on the open subset $\psi_\alpha(U_\alpha\times O_\alpha)\subset X$.

    Choose a conically smooth partition of unity $\{\eta_\alpha\}$ on $Y$ subordinate to the cover $\cU_Y$. Then
    \begin{equation}
        W:=\sum_\alpha (f^\ast\eta_\alpha)\,W_\alpha
    \end{equation}
    is a well-defined parallel vector field on $X$. Since each $W_\alpha$ lifts $V_{|O_\alpha}$ and the coefficients come from the base, the same partition-of-unity argument as in the smooth case shows that $W$ lifts $V$. 
\end{proof}

These results lead us closer to the proof of the Ehresmann fibration theorem. We are only missing a standard fact about flow extensions over compact subsets. 

\begin{lemma}\label{lem:compact flows}
    Let $K\subset X$ be a compact subset of a conically smooth stratified space $Z =(X\longra P)$, and let $V$ be a parallel vector field on $X$ with maximal flow given by $\gamma^V \maps A \longra X$, with $A\subset X\times \RR$ an open neighborhood of $X\times \{0\}$. Let $x\in X_p \subset X$, and let 
    \begin{equation}\label{eq:lem:compact flows}
        \{x\}\times J_x  = A \cap (\{x\}\times \RR)
    \end{equation}
    such that $J_x = (a_x,b_x)$ and contains $0$. If $b_x<\infty$, then there exists $\epsilon > 0$ such that $\gamma^V_t(x)\notin K$ for $b_x - \epsilon < t < b_x$. Likewise, if $a_x>-\infty$, then there exists $\epsilon > 0$ such that $\gamma^V_t(x)\notin K$ for $a_x < t < a_x+\epsilon$.
\end{lemma}
\begin{proof}
    Since $A$ is an open neighborhood of $K\times \{0\}$, compactness gives $\epsilon>0$ such that
    \begin{equation}\label{eq:lem:prf:compact flows cover}
        K\times [-\epsilon,\epsilon] \subset A.
    \end{equation}
    In other words, every point of $K$ can be flowed for at least time $\epsilon$ in both directions.

    Suppose $b_x<\infty$ and that there exists a sequence $t_n \uparrow b_x$ with $\gamma^V_{t_n}(x)\in K$. For each $n$, the point $\gamma^V_{t_n}(x)$ lies in $K$, so by \Cref{eq:lem:prf:compact flows cover} the flow is defined for an additional time interval $[-\epsilon,\epsilon]$ starting at $\gamma^V_{t_n}(x)$. Therefore $\gamma^V_{t_n+\epsilon}(x)$ is defined. For $n$ large enough, however, $t_n+\epsilon>b_x$, contradicting the maximality of $J_x$. Hence $\gamma^V_t(x)\notin K$ for all $t$ sufficiently close to $b_x$ from the left.

    The argument for the left endpoint is identical, applied to a sequence $t_n \downarrow a_x$.
\end{proof}


\begin{remark}\label{rmk:flow out of compact}
\Cref{lem:compact flows} shows that a maximal trajectory cannot terminate while remaining inside a fixed compact subset. In particular, if the image of a trajectory is contained in a compact subset for every bounded time interval, then the corresponding flow is defined for all time.
\end{remark}

We can now finally state and prove the main theorem of this section — the conically smooth stratified analogue of the Ehresmann fibration theorem. 

\begin{theorem}[Thom/Ehresmann/Morse for conically smooth stratified spaces]\label{ch4:: thm : stratified ehresmann}
    Let $\pi\maps X \longra Y$ be a conically smooth proper submersion. Then $\pi$ is a conically smooth fiber bundle.
\end{theorem}
\begin{proof}
Recall that $\pi$ is a conically smooth fiber bundle if for every $y\in Y$, there exists a neighborhood $U\subset Y$, which we can assume centered at $y$, and a conically smooth isomorphism
\begin{equation}
    h_U\maps \pi^{-1}(U) \longra \pi^{-1}(\{y\})\times U
\end{equation}
participating in the following commutative diagram: 
\begin{equation}\label{eq:thm:Ehresmann}
\begin{tikzcd}[ampersand replacement=\&]
	{\pi^{-1}(U)} \&\& {\pi^{-1}(\{y\})\times U} \\
	\\
	\& U \& {.}
	\arrow["{h_U}", from=1-1, to=1-3]
	\arrow["{\pi_{|\pi^{-1}(U)}}"', from=1-1, to=3-2]
	\arrow["{{\sf pr}}", from=1-3, to=3-2]
\end{tikzcd}
\end{equation}
Because the definition is local on the base, after shrinking around $y$ we may suppose that $Y=\RR^k$ and that $y=0\in \RR^k$. Let $\partial_1,\ldots, \partial_k$ be the coordinate vector fields on $\RR^k$. By \Cref{prop:lifts of vector fields}, there exists a parallel vector field $\partial_{i}^X$ on $X$ for each $i\in \{1,\cdots, k\}$ such that 
\begin{equation}\label{eq:thm:str_ehrs:prf:lift}
    (\sD_{x}\pi)\,\big(\partial^X_{i}(x)\big)= \partial_{i}(\pi(x)) \quad \forall x\in X. 
\end{equation}
By \Cref{eq:def:flows}, each of the lifted parallel vector fields $\partial^X_{i}$ on $X$ generates a local flow 
\begin{equation}\label{eq:thm:prf:flow}
\gamma^{\partial^X_{i}}\maps X \times \RR \dashrightarrow X,
\end{equation} defined on an open neighborhood $A_i \subset X\times \RR$ of $X\times \{0\}.$ Projecting this flow down to $\RR^k$, we note that it is given by the simple explicit formula
\begin{equation}
    \pi\big(\gamma^{\partial^{X}_{i}}_{t}(x)\big) = \pi(x) + (0,\ldots,0,t,0,\ldots,0),
\end{equation} where the only non-vanishing coordinate is $i$. 

Fix $x\in X$ and $i\in \{1,\ldots,k\}$, and write
\begin{equation}
    \{x\}\times J_{x,i}=A_i\cap (\{x\}\times \RR)
\end{equation}
for the maximal interval of definition of the flow through $x$. Suppose $b_{x,i}<\infty$. Choose numbers $a<0<b_{x,i}<b$. The formula above shows that
\begin{equation}
    \pi\big(\gamma^{\partial^{X}_i}_{t}(x)\big)\in K:=\pi(x)+\{(0,\ldots,0,t,0,\ldots,0)\mid t\in [a,b]\}\subset \RR^k
\end{equation}
for every $t\in J_{x,i}\cap (a,b)$. The set $K$ is compact, so properness of $\pi$ implies that $\pi^{-1}(K)$ is compact. Hence the image of the trajectory on $J_{x,i}\cap (a,b)$ lies in a fixed compact subset of $X$. By \Cref{lem:compact flows}, this is incompatible with a finite right endpoint. Therefore $b_{x,i}=+\infty$. The same argument applied to the left endpoint shows $a_{x,i}=-\infty$. Thus each $\partial_i^X$ is complete.

We introduce the notation $t_{i}(x):=\pi_{i}\big(\pi(x)\big)$, where $\pi_i$ is the projection onto the $i$-th factor. We define
\begin{equation}
    h\maps X \longrightarrow \pi^{-1}(\{0\})\times \RR^k
\end{equation}
as follows:
\begin{align}
    h(x) &= \bigg(\Big(\gamma^{\partial^{X}_{1}}_{-t_1(x)}\circ \cdots \circ \gamma^{\partial^{X}_{k}}_{-t_k(x)}\Big)(x), \pi(x)\bigg).
\end{align}
Projecting the first component to $\RR^k$ and using the displayed formula for the base flow, each negative-time flow subtracts the corresponding coordinate of $\pi(x)$. Hence the first component indeed lands in $\pi^{-1}(\{0\})$, so $h$ is well-defined. The map is conically smooth because it is built from conically smooth flows and projections. To show it is an isomorphism, define
\begin{equation}
    \bar{h}\big(x_0,(t_1,\cdots,t_k)\big) = \Big(\gamma^{\partial^{X}_{k}}_{t_k}\circ \cdots \circ \gamma^{\partial^{X}_{1}}_{t_1}\Big)(x_0).
\end{equation} 
This map is again conically smooth. The identities $\bar{h}\circ h = \id_X$ and $h\circ \bar{h}=\id_{\pi^{-1}(0)\times \RR^k}$ follow by repeatedly applying the additivity relation from \Cref{eq:def:flows}; no commutativity between the different flows is needed because the two compositions are written in opposite orders. Therefore $\pi$ is locally isomorphic to the projection $\pi^{-1}(\{0\})\times \RR^k \longra \RR^k$, as required.
\end{proof}

\begin{remark}
    We will need a variant of the stratified Ehresmann lemma for conically smooth stratified spaces \emph{with corners}, for our applications to stratified bordisms. This is the content of the next theorem.
\end{remark}

\begin{cor}\label{ch4:: cor: stratified ehresmann boundary}
    Let $\big(M \ra P\big)$ and $K$ be a conically smooth stratified space and a smooth manifold, respectively. Assume $\big(M \ra P\big)$ and $K$ both carry boundary structures. If $\pi\maps M \ra K$ is a conically smooth proper submersion such that $\pi(\partial M) \subseteq \partial K$, then $\pi$ is a conically smooth fiber bundle. 
\end{cor}
\begin{proof}
    The argument of the proof follows that of the unstratified case, exhibited in \Cref{ap:: prop: Ehresmann fibration no boundary}. For boundary points, the vector field lift must necessarily be inward-pointing, and the existence of the corresponding flow is ensured by \Cref{ch4:: def: in flows}.
\end{proof}




\subsubsection{Exit-path \texorpdfstring{$(\oo,1)$}{(infty,1)}-category}

The singular complex $\sing(M)$ of a smooth manifold $M$ is a fundamental construction in homotopy theory and manifold topology. Any conically smooth stratified space $\big(M \ra P\big)$ also admits a singular complex $\sing(M)$. Yet given, for example, a refinement of $\big(M\ra P\big) \longra (M' \ra Q)$, where the underlying spaces are homeomorphic, this induces an evident equivalence between singular complexes $\sing(M)\longra \sing(M')$ — the caveat is that $M$ and $M'$ are not in general stratified homotopy equivalent.

The construction of the exit-path category (or its 2-truncation bicategory) dates back to David Treumann's work on exit-path 2-categories and constructible stacks \cite{Tre09} together with Lurie's work (see \cite{LurieHA}) on the $(\oo,1)$-category version — the latter is the one which we will be discussing here. Before proceeding with the definition, we must first introduce an important stratification of the topological standard $n$-simplex $\Delta^n$.

\begin{definition}\label{ch4:: def: standard stratification}
There is a cosimplicial object in $\strat$, hence also in $\Strat$,
\begin{align}
    \mathsf{st} \maps \bDelta &\longra \strat \\
    [p] &\longmapsto \big(\,\Delta^{p} \xra{m} [p] \,\big)
\end{align}
where $m\maps (t_0,\cdots, t_p) \longmapsto \mathsf{Max}\{i\, |\, t_i \neq 0\}$.
\end{definition}

This cosimplicial object is well-defined, since giving a morphism $f \maps [p] \ra [q]$ in $\bDelta$, the usual map of spaces $|f| \maps \bDelta^{p} \longra \bDelta^{q}$ extends to a map of stratified spaces 
\begin{equation}
    \begin{tikzcd}
        {\bDelta^{p}} && {\bDelta^{q}} \\
        \\
        {[p]} && {[q]}
        \arrow["{|f|}", from=1-1, to=1-3]
        \arrow["m"', from=1-1, to=3-1]
        \arrow["m", from=1-3, to=3-3]
        \arrow["f"', from=3-1, to=3-3]
    \end{tikzcd}
\end{equation}
since $f$ is monotonic. 

\begin{observation}
We grant that this is a very unorthodox stratification of the topological $p$-simplex. Other, more pedestrian stratifications one could conceive of include, for example, a stratification by dimension of the face of the simplex or, as a substratification, just two strata: the interior and the boundary. To help make sense of this stratification we draw the stratifications until $p=3$.

\smallskip
\begin{equation}
    \tikzfig{figures/Strat/STR13}{\begin{tikzpicture}
	\begin{pgfonlayer}{nodelayer}
		\node [style=none] (14) at (11.5, 5) {};
		\node [style=none] (15) at (14, 0) {};
		\node [style=none] (16) at (14.75, 2) {};
		\node [style=none] (2) at (-2, 0) {};
		\node [style={endpoint_yellow_small}] (3) at (-7, 0) {};
		\node [style=none] (4) at (6, 0) {};
		\node [style=none] (6) at (3.5, 4) {};
		\node [style=none] (7) at (1, 0) {};
		\node [style=none] (8) at (6, 0) {};
		\node [style={endpoint_yellow_small}] (5) at (1, 0) {};
		\node [style=none] (9) at (14, 0) {};
		\node [style=none] (17) at (11.5, 5) {};
		\node [style=none] (18) at (9, 0) {};
		\node [style=none] (19) at (14, 0) {};
		\node [style=none] (20) at (11.5, 5) {};
		\node [style=none] (21) at (9, 0) {};
		\node [style=none] (22) at (14.75, 2) {};
		\node [style={endpoint_yellow_small}] (13) at (9, 0) {};
		\node [style=none] (10) at (14.75, 2) {};
		\node [style=none] (11) at (9, 0) {};
		\node [style=none] (12) at (14, 0) {};
		\node [style={endpoint_yellow_small}] (23) at (-10, 0) {};
		\node [style=none] (24) at (-10, -1.75) {$\Delta^0$};
		\node [style=none] (25) at (-4.5, -1.75) {$\Delta^1$};
		\node [style=none] (26) at (3.5, -1.75) {$\Delta^2$};
		\node [style=none] (27) at (11.5, -1.75) {$\Delta^3$};
		\node [style=none] (28) at (-10, 0.75) {$0$};
		\node [style=none] (29) at (-7, -0.5) {$0$};
		\node [style=none] (30) at (-2, -0.5) {$1$};
		\node [style=none] (31) at (1, -0.5) {$0$};
		\node [style=none] (32) at (6, -0.5) {$1$};
		\node [style=none] (33) at (3.5, 4.5) {$2$};
		\node [style=none] (34) at (14, -0.5) {$1$};
		\node [style=none] (35) at (9, -0.5) {$0$};
		\node [style=none] (36) at (15.25, 2) {$2$};
		\node [style=none] (37) at (11.5, 5.5) {$3$};
	\end{pgfonlayer}
	\begin{pgfonlayer}{edgelayer}
		\draw [style=defect line 1] (20.center) to (19.center);
		\draw [style=defect patch 1 dark] (14.center)
			 to (15.center)
			 to (16.center)
			 to cycle;
		\draw [style=defect line 3] (3) to (2.center);
		\draw [style=defect patch 2] (5) to (6.center);
		\draw [style=defect patch 2] (6.center) to (4.center);
		\draw [style=defect patch 2] (7.center)
			 to (8.center)
			 to (6.center)
			 to cycle;
		\draw [style=defect line 3] (5) to (4.center);
		\draw [style=defect patch 2] (13) to (10.center);
		\draw [style=defect patch 2] (10.center) to (9.center);
		\draw [style=defect patch 1 dark] (17.center)
			 to (19.center)
			 to (18.center)
			 to cycle;
		\draw [style=defect patch 1 dark] (20.center)
			 to (22.center)
			 to (21.center)
			 to cycle;
		\draw [style=defect line 3] (13) to (9.center);
		\draw [style=defect patch 2] (11.center)
			 to (12.center)
			 to (10.center)
			 to cycle;
	\end{pgfonlayer}
\end{tikzpicture}
}
\end{equation}
\medskip

\noindent There is a more natural way of understanding this asymmetric stratification through the lens of cones. Namely, remember the closed cone construction $\overline{\sC}(-)$ which produces a stratified space with boundary. If one takes the closed cone on the empty set 
\begin{equation}
\overline{\sC}(\emptyset \ra [0])
\end{equation} which is precisely $\Delta^1 \ra [1]$. Moreover, iterating this $(p-1)$ times we get 
\begin{equation}
    \overline{\sC}^{p}(\emptyset \ra [0]) = \big(\,\Delta^{p} \ra [p]\,\big).
\end{equation} 
So that for every $p\geq 0$, we have the stratified identification 
\begin{equation}
\overline{\sC}\big(\bDelta^{p-1}\xra{m}[p-1]\big) \, \simeq \, \big(\bDelta^{p} \xra{m} [p] \big).
\end{equation}
\end{observation}

The next theorem is both a definition and a theorem.

\begin{theorem}\cite[Corollary 3.3.6]{AFR1}
    The restricted Yoneda embedding associated to $\mathsf{st}\maps \bDelta \ra \Strat$ 
    \begin{equation}
    \exit \maps \Strat \xra{\,\,\mathsf{Yoneda}\,\,} \psh(\Strat) \xra{\,\,\,\mathsf{st}^{*}\,\,\,} \psh(\bDelta)
    \end{equation}
    lands in complete Segal spaces. We call the value of $\exit$ on a conically smooth stratified space its \textbf{exit-path $(\oo,1)$-category}. 
\end{theorem}

\begin{observation}
Unpacking this definition we see that given a conically smooth stratified space $\big(M \ra P\big)$, the space of objects of its exit-path $(\oo,1)$-category is:
\begin{equation}
    \exit(M)_{0} = \Strat(*,M)
\end{equation}
which as a Kan complex has as $q$-simplices the set $\Strat(*\times \AA^q, M)$, an element of which, as a map out of a smooth manifold, must factor through a single smooth stratum $M_p$ for some $p\in P$. The space of objects can thus be identified with
\begin{equation}
    \exit(M)_{0} \,\, \simeq \,\, \coprod_{p \in P} M_p
\end{equation} 
the space which is the disjoint union of the individual strata.
\end{observation}

\begin{observation} Observe that exit-paths do not compose strictly. Exit-paths are the $0$-simplices $\big(\gamma\maps \bDelta^1 \ra \big(M \ra P\big)\big)\in \Strat(\bDelta^{1}\times \AA^0, M)$ — these are stratified maps with an important feature: as any map out of $[0,1]$, it traces a path in $M$, but as a stratified map it forces the path to start at a deeper stratum than the one it ends. Moreover, once the path leaves its starting stratum — if it does — it never returns to it. In a very precise sense, exit-paths are extremely restricted. 


Their restricted nature invites deeper reflection on how they compose — given the fact that they form an $(\oo,1)$-category. At a distance, they do not compose strictly, as is exemplified by the next figure
\begin{equation}
    \tikzfig{figures/Strat/STR14}{\begin{tikzpicture}
	\begin{pgfonlayer}{nodelayer}
		\node [style=none] (0) at (-3.5, 3.5) {};
		\node [style=none] (1) at (3.5, 3.5) {};
		\node [style=none] (2) at (-3.5, -3.5) {};
		\node [style=none] (3) at (3.5, -3.5) {};
		\node [style=none] (4) at (0, -3.5) {};
		\node [style=none] (5) at (0, 3.5) {};
		\node [style=dot] (6) at (0, 0) {};
		\node [style=none] (7) at (0, 1.75) {};
		\node [style=none] (8) at (-2, 1.75) {};
		\node [style=none] (9) at (0.5, 0.75) {$\gamma_1$};
		\node [style=none] (10) at (-1, 2.5) {$\gamma_2$};
	\end{pgfonlayer}
	\begin{pgfonlayer}{edgelayer}
		\draw [style=defect patch 2] (1.center)
			 to (0.center)
			 to (2.center)
			 to (3.center)
			 to cycle;
		\draw [style=defect line 3] (5.center) to (4.center);
		\draw [style=dir thick] (6) to (7.center);
		\draw [style=dir thick, in=30, out=-165] (7.center) to (8.center);
	\end{pgfonlayer}
\end{tikzpicture}
}
\end{equation}
where the underlying stratified space is $\underline{\RR}^2 =\big(\RR^2 \ra [2]\big)$. The two exit-paths $\gamma_1,\gamma_2\maps \bDelta^{1} \ra \RR^2$ are such that the source of $\gamma_2$ is the target of $\gamma_1$, which would allow for a composite continuous map $|\bDelta^1| \ra \RR^2$ -- yet, as stratified maps i.e., exit-paths, they cannot compose strictly. Of concern is whether the composition is well-defined up to homotopy. To see why this is true we make use of the homotopy colimit associated to the \emph{double closed cones} in $\Strat$

\begin{equation}
    \begin{tikzcd}
        {\overline{\sC}(\emptyset)} && {\overline{\sC}(L)} \\
        \\
        {\overline{\sC}^2(\emptyset)} && {\overline{\sC}^2(L)}
        \arrow["{\overline{\sC}(\emptyset \hookra L)}", hook, from=1-1, to=1-3]
        \arrow["{\{1\}}"', hook, from=1-1, to=3-1]
        \arrow["{\{1\}}", hook, from=1-3, to=3-3]
        \arrow["{\overline{\sC}^2(\emptyset \hookra L)}"', hook, from=3-1, to=3-3]
        \arrow["\lrcorner"{anchor=center, pos=0.125, rotate=180}, draw=none, from=3-3, to=1-1]
    \end{tikzcd}
\end{equation}
where $L$ is any compact conically smooth stratified space — note that in $\strat$ this fails to be a colimit, indicating the need to work in the suitably homotopical context of $(\oo,1)$-categories. An implication of this diagram being a pushout in $\Strat$, is that for every other conically smooth stratified space $Z$, the diagram 

\begin{equation}
    \label{ch4:: eq: exit-paths compose}
    \begin{tikzcd}
        {\Strat\big(\overline{\sC}(\emptyset),Z\big)} && {\Strat\big(\overline{\sC}(L),Z\big)} \\
        \\
        {\Strat\big(\overline{\sC}^2(\emptyset),Z\big)} && {\Strat\big(\overline{\sC}^2(L),Z\big)}
        \arrow["{\overline{\sC}(\emptyset \hookra L)}"', from=1-3, to=1-1]
        \arrow["{\{1\}}", from=3-1, to=1-1]
        \arrow["{\{1\}}"', from=3-3, to=1-3]
        \arrow["{\overline{\sC}^2(\emptyset \hookra L)}", from=3-3, to=3-1]
    \end{tikzcd}
\end{equation}
is a pullback of Kan complexes in $\sset$. If $Z$ is $(M\ra P)$ and if $L=\Delta^{\{0\}}$, then we have that 
\begin{equation}
     \Strat\big(\overline{\sC}(\Delta^{\{0\}}),Z\big)\underset{\Strat\big(\overline{\sC}(\emptyset),Z\big)}{\times} \Strat\big(\overline{\sC}^2(\emptyset),Z\big)
\end{equation}
records the space of composable exit-paths. The pullback condition then says that given a pair of composable exit-paths $\gamma_1,\gamma_2$, up to homotopy, there is a unique exit 2-simplex $\sigma\maps \Delta^{\{0<1<2\}} \ra M$ whose $\Delta^{\{0<1\}}$ and $\Delta^{\{1<2\}}$ faces are $\gamma_1$ and $\gamma_2$, respectively.
\end{observation}



\subsubsection{Embedding into Euclidean space}

Studying spaces of embeddings of smooth manifolds has turned out almost as important and fruitful as studying the manifolds themselves. In this section, we initiate the exploration of the homotopical properties of embedding spaces of conically smooth stratified spaces into Euclidean space. We begin by posing a conjecture that every conically smooth stratified space $M$ of dimension $n$ admits a conically smooth embedding into $\RR^{2n + 1}$ — we prove it for the special case of depth 1, and indicate a possible path to the proof for general depth —, we prove that every conically smooth stratified space $M$ admits a conically smooth embedding into some Euclidean space, that the space of conically smooth embeddings of $M$ into $\RR^{\oo}$ is weakly contractible, and finally that every conically smooth stratified space $M$ with boundary admits a neat embedding into $\RR_{\geq 0} \times \RR^{m}$.  

One of the most important results in differential topology is the \emph{Whitney embedding theorem}. In its weakest form it asserts that every compact smooth manifold admits a proper embedding into Euclidean space.

\begin{theorem}[Whitney embedding theorem — weak version]\label{ch4:: thm: Whitney weak}
    Let $M$ be a compact smooth manifold. There exists a smooth embedding $M \hookra \RR^{N}$ for a sufficiently large $N$
\end{theorem}

\noindent The proof of this result rests almost entirely on the fact that smooth manifolds admit partitions of unity, and Whitney showed that it can be improved significantly by using transversality arguments. 

\begin{theorem}[Whitney embedding theorem — strong version]\label{ch4:: them: Whitney strong}
    \label{thm:ch7::whitney strong}
Let $M$ be a compact smooth $n$-manifold. There exists an immersion of $M$ into $\RR^{2n-1}$ and an embedding into $\RR^{2n}$.     
\end{theorem}

In working with conically smooth stratified spaces, which are defined abstractly — as opposed to Whitney stratified spaces, which are always explicitly embedded into a smooth manifold — in order to understand the topology of their embedding spaces, we must first show a result analogous to the Whitney embedding theorem, ideally in its strongest form. Before doing so, we introduce a quick technical result.

\begin{lemma}\cite[Lemma~1.3.1]{AFR1}\label{ch4:: lem: shrinking} For every open cover $\cU$ of a conically smooth stratified space $K$, there exists a refining open cover $\cV$ with the property that for each $V\in \cV$, the image of $V \hookra K$ factors as
    \begin{equation}
        \begin{tikzcd}
            V && {\RR^{i}\times \sC(L)} \\
            \\
            && {\DD^{i}\times \overline{\sC}(L)}
            \arrow[from=1-1, to=1-3]
            \arrow["i"', hook, from=1-1, to=3-3]
            \arrow["{\overline{\phi}}"', hook, from=3-3, to=1-3]
        \end{tikzcd}
    \end{equation}
    where the vertical map is a conically smooth embedding and the diagonal map is an open embedding i.e., an isomorphism into the interior $\RR^{i}\times \sC(L) \subset \DD^{i}\times \overline{\sC}(L)$.
\end{lemma}

Using it, we can prove the stratified analogue of Whitney's weak embedding theorem \Cref{ch4:: thm: Whitney weak}.

\begin{proposition}\label{ch4:: prop: weak Whit strat}
    Let $\big(M \ra P\big)$ be a finite conically smooth stratified space. There is a proper conically smooth embedding $\big(M \ra P\big) \hookra \RR^{N}$ for some $N$. 
\end{proposition}
\begin{proof}
    We argue by induction on depth. If $\mathsf{depth}(M)=0$, the result is the usual weak Whitney embedding theorem for smooth manifolds.

    Assume $\mathsf{depth}(M)>0$. By \Cref{ch4:: lem: shrinking} and finiteness of $M$, we can choose a finite conically smooth atlas $\cA=\{U_i\}_{i=1}^{k}$ such that each chart admits a conically smooth identification
    \begin{equation}
        U_i \simeq \RR^{l_i}\times \sC(Z_i),
    \end{equation}
    where $Z_i$ is compact and satisfies $\mathsf{depth}(Z_i)<\mathsf{depth}(M)$. Choose a conically smooth partition of unity $\{\psi_i\}_{i=1}^{k}$ subordinate to this cover.

    By the inductive hypothesis, for each $i$ there exists a proper conically smooth embedding
    \begin{equation}
        \varphi_i \maps Z_i \hookra \RR^{N_i} \hookra S^{N_i}.
    \end{equation}
    Taking open cones and then the product with the identity on $\RR^{l_i}$ gives a conically smooth embedding
    \begin{equation}
        f_i \maps U_i \hookra \RR^{l_i}\times \sC(S^{N_i}) \simeq \RR^{l_i+N_i+1}.
    \end{equation}

    Because $\psi_i$ is supported in $U_i$, both maps $x\mapsto \psi_i(x)$ and $x\mapsto \psi_i(x)\,f_i(x)$ extend by zero to global conically smooth maps on $M$. Define
    \begin{equation}
        F := \big(\psi_1,\psi_1f_1,\ldots,\psi_k,\psi_kf_k\big)\maps M \longra \RR^{N}
    \end{equation}
    for a sufficiently large $N$.

    The same argument as in the smooth case now shows that $F$ is a conically smooth embedding: the scalar components $\psi_i$ detect an index for which $\psi_i>0$, and on the open set $\{\psi_i>0\}$ the pair $\big(\psi_i,\psi_i f_i\big)$ recovers the local embedding $f_i$. Hence $F$ is injective and locally a conically smooth embedding, therefore a global conically smooth embedding.

    Finally, only finitely many compactified charts occur in the construction, so $F$ is proper. This completes the inductive step.
\end{proof}

With this result, we can now prove the following important structural proposition about the space of conically smooth embeddings into Euclidean spaces of sufficiently high dimension.

\begin{proposition} \label{ch4:: prop: contractibility of embeddings}
    Let $\big(M \ra P\big)$ be a finite conically smooth stratified space. The space 
    \begin{equation}
        \Strat^{\sf emb}\big(M, \RR^{\oo}\big) \, \overset{\sf def}{=} \, \hoColimit[\RR^{k}\hookra \RR^{k+1}] \Strat^{\sf emb}(M, \RR^{k})
    \end{equation} 
    is weakly contractible.
\end{proposition}

\begin{proof}
We first observe that the vector space $\RR^{\oo}$ admits a short exact sequence 
\begin{equation}
    0 \longra \RR^{\oo} \xlongra{S^{\sf ev}} \RR^{\oo} \xra{P^{\sf ev}} \RR^{\oo} \longra 0.
\end{equation}

\noindent The two maps are defined as follows 
\begin{equation}
    S^{\sf ev}\big(x_i\big) \, = \, (0,x_1, 0,x_2, 0, \cdots), \quad P^{\sf ev}\big(x_i\big) \, = \, (x_1, x_3, x_5, \cdots).
\end{equation}
Equivalently, there is a linear isomorphism 
\begin{equation}
    \RR^{\oo} \oplus \RR^{\oo} \xra{\,\,\,\,\,S^{\sf ev} + S^{\sf od}\,\,\,\,\,} \RR^{\oo},
\end{equation}
where $S^{\sf od}\big(x_i\big) = (x_1, 0, x_2, 0, \cdots)$.  

Let us consider the smooth homotopy 
\begin{equation}
    H \maps \RR^{\oo} \times \AA^1 \longra \RR^{\oo}\times \AA^1, \quad (x,t) \longmapsto \big((1-t)\cdot x + t \cdot S^{\mathsf{ev}}(x), t \big)
\end{equation}
Assume that $H(-,t)$ fails to be injective for some $t'\in \AA^1$. Since the case $t'=0$ is just the identity, we may assume $t'\neq 0$, in which case there exist $x-x' \neq 0$ satisfying 
\begin{equation}
    (1-t')x + t' S^{\mathsf{ev}}(x) = (1-t')x' + t' S^{\mathsf{ev}}(x').
\end{equation}
This implies that 
\begin{equation}
S^{\sf ev}(x-x') = -\frac{1-t'}{t'}(x-x')
\end{equation}
or that $(x-x')$ is an eigenvector for $S^{\sf ev}$ — a contradiction.
\smallskip
\noindent We will now use this (conically) smooth isotopy of (conically) smooth embeddings to build a simplicial homotopy  
\begin{equation}
    \Strat^{\sf emb}(M, \RR^{\oo})\times \Delta^{1} \longra \Strat^{\sf emb}(M, \RR^{\oo})
\end{equation}
from $\id_{\Strat^{\sf emb}(M,\RR^{\oo})}$ to $(S^{\sf ev})_{*}$. First, observe that such a simplicial homotopy amounts to a simplicial map 
\begin{equation}
    \Strat^{\sf emb}(M, \RR^{\oo}) \longra \Strat^{\sf emb}(M \times \AA^1, \RR^{\oo})
\end{equation}   
owing to the constructibility of $\Strat^{\sf emb}$. We define this map as follows:
\begin{align}    
    \big(\, \varphi_p \maps M \times \AA^p \ra \RR^{\oo} \times \AA^p \,\big) &\longmapsto 
    \big(\, M \times \AA^p \times \AA^{1} \xra{\varphi_p \times \id} \RR^{\oo} \times \AA^p \times \AA^1  \\
    &\xra{\,\,\id \times \mathsf{swap}\,\,} \RR^{\oo} \times \AA^1 \times \AA^p \\
    &\xra{\,\, H \times \id \,\,} \RR^{\oo} \times \AA^{1} \times \AA^p \\ 
    &\xra{\,\, \sf pr \,\,} \RR^{\oo} \times \AA^p \, \big).
\end{align}
The map is well-defined since $H\big(\varphi_p(x,s),t\big)$ is a conically smooth embedding for every $s \in \AA^p$. Moreover, we have that 
\begin{equation}
    H(\varphi_p(x,s),0) = \big(\varphi_{p}(x,s),0\big), \quad H(\varphi_{p}(x,s),1) = \big(S^{\sf ev}(\varphi_{p}(x,s)),1\big),
\end{equation}
further observing the correct start- and end-points for the conically smooth isotopy.

From \Cref{ch4:: prop: weak Whit strat} we know there exists a conically smooth embedding $e \maps M \hookra \RR^{\oo}$. Since $M$ is finite, after postcomposing with a standard embedding of $\RR^{N}$ into a sphere in $\RR^{N+1}$ and then into $\RR^{\oo}$, we may assume the image of $e$ lies in the unit sphere. We now articulate a contraction of $\Strat^{\sf emb}(M,\RR^{\oo})$ onto the point represented by $S^{\sf od}\circ e$:

\begin{align}
    \Strat^{\sf emb}(M,\RR^{\oo})_p &\longra \Strat^{\sf emb}(M\times \AA^{1},\RR^{\oo})_p \\
    \big(\, \varphi_p \maps M \times \AA^p \ra \RR^{\oo} \times \AA^p \,\big) &\longmapsto \big[(x,t,s) \mapsto \big((1-t)\, S^{\sf ev}(\varphi_p(x,s)) + t\, S^{\sf od}(e(x)), s\big)\big].
\end{align}
The map is well-defined because the even coordinates recover $(1-t)\varphi_p(x,s)$ while the odd coordinates recover $t\,e(x)$. Since $e(x)$ lies on the unit sphere, its norm recovers $t$, and then $e$ recovers $x$; together with the final $\AA^p$-coordinate this shows injectivity. This provides a conically smooth isotopy of conically smooth embeddings starting at $(S^{\sf ev})_*$ and ending at the constant point $S^{\sf od}\circ e$. Through a suitable rescaling, both isotopies can be concatenated to connect $\id_{\Strat^{\sf emb}}$ to a constant map. 
\end{proof}

\Cref{ch4:: prop: contractibility of embeddings} will pay off when we attempt to construct the $(\oo,1)$-category of stratified bordisms using conically smooth embedded stratified spaces together with carefully chosen submersivity conditions — as we did in Chapter 3.   

We now consider the strong version of the Whitney embedding theorem for conically smooth stratified spaces. We do not prove it for arbitrary depth, instead leaving it as the following conjecture.

\begin{conjecture}\label{ch4:: conj: whitney strat}
Let $(M \ra P)$ be a compact $d$-dimensional conically smooth stratified space. There exists a conically smooth proper embedding 
\begin{equation}
    (M \ra P) \hookra \RR^{2d + 1}.
\end{equation}
\end{conjecture}

We provide a proof for depth 1, but before introducing the proof, we mention the work of Hiroko Natsume proving \Cref{ch4:: conj: whitney strat} for Whitney stratified spaces \cite{N80}. In particular, we make use of the following result in \cite{N80}.

\begin{proposition}\cite[Lemma~1]{N80}\label{ch4:: prop: fiberwise emb}
    Let $\pi\maps E \ra M$ be a fiber bundle with a compact manifold $F$ as the fiber and $h$ a diffeomorphism of $M$ onto $M'$. Let $S$ be a sphere bundle over $M'$ and $p$ the projection of $S$ to $M'$ 
    \begin{equation}
        \begin{tikzcd}[ampersand replacement=\&]
            E \&\& S \\
            \\
            M \&\& {M'}
            \arrow["\pi"', two heads, from=1-1, to=3-1]
            \arrow["p", two heads, from=1-3, to=3-3]
            \arrow["\simeq", from=3-1, to=3-3]
        \end{tikzcd}
    \end{equation}
    Assuming that 
    \begin{equation}
        \mathsf{dim}(S) \geq 2\cdot \mathsf{dim}(E) + 2,
    \end{equation}
    then there is a fiberwise smooth embedding $\tilde{h}\maps E \hookra S$ 
    \begin{equation}
        \begin{tikzcd}[ampersand replacement=\&]
            E \&\& S \\
            \\
            M \&\& {M'}
            \arrow["{\tilde{h}}", dashed, hook, from=1-1, to=1-3]
            \arrow["\pi"', two heads, from=1-1, to=3-1]
            \arrow["p", two heads, from=1-3, to=3-3]
            \arrow["\simeq", from=3-1, to=3-3]
        \end{tikzcd}
    \end{equation} 
\end{proposition}

Lastly, before proving the stratified strong embedding theorem for depth 1, we must introduce yet another theorem of Whitney.

\begin{theorem}[Whitney extension theorem]\label{ch4:: thm: whitney extension theorem}
    Let $A$ be a smooth $m$-dimensional manifold, and let $B\subset A$ be a closed subset of $A$. Let $M=A\setminus B$ be a smooth manifold, and let $f\maps B \longra \RR^{n}$. Then there is a map $F$ of $A$ into $\RR^{n}$ with the following properties:
    \begin{enumerate}
        \item $F_{|B}=f$.
        \item $F$ is smooth.
        \item If $n\geq 2m +1$ and $f$ is an embedding of $B$, then $F$ is an embedding of $A$.
    \end{enumerate}
\end{theorem}

We are now ready to state and prove the strong stratified Whitney embedding theorem in depth 1.

\begin{proposition}\label{ch4:: prop: Whitney strong depth 1}
    Let $\big(  M\longra P \big)$ be a compact conically smooth stratified space of dimension $n$ and depth 1. There exists a proper conically smooth embedding $M \hookra \RR^{2n + 1}$.
\end{proposition}

\begin{proof}
    Since $M$ has depth $1$, we may assume that the poset $P$ is isomorphic to $\{d<n\}$. Write $M_d$ for the deepest stratum and $M_n$ for the top-dimensional stratum.
    
    The stratum $M_d$ is a compact smooth $d$-manifold, so \Cref{thm:ch7::whitney strong} gives a smooth embedding
    \begin{equation}
        \varphi_d \maps M_d \hookra \RR^{2d} \hookra \RR^{2n+1}.
    \end{equation}
    The image $\varphi_d(M_d)$, together with its complement, endows $\RR^{2n+1}$ with a stratification 
    \begin{equation}
        \label{eq:ch7::image of embedding}
        \big(\RR^{2n+1} \longra \{d < (2n + 1)\}\big)
    \end{equation}    
    which is conically smooth and of type $\sD_{d\subset 2n+1}$. Our goal is to extend $\varphi_d$ to a conically smooth embedding of all of $M$. We first extend it over a regular neighborhood of $M_d$, and then over the top stratum.

    Taking the closures of the regular neighborhoods provides the following identifications of conically smooth stratified spaces with boundary:
    \begin{equation}
        \label{eq:ch7::regular neighborhood}
        \overline{\sC}(\pi_d) ~\simeq~ \Big(\Link_{d}(M) \times [0,1] \Big)\underset{\Link_d(M)}{\bigcup} M_d, \qquad \overline{\sC}(\widetilde{\pi}_d) ~\simeq~ \Big(\Link_{d}\big(\RR^{2n+1}\big) \times [0,1] \Big) \underset{\Link_d(\RR^{2n+1})}{\bigcup} \varphi_d(M_d).
    \end{equation}
    These identifications provide coordinates $[z,t]\in \overline{\sC}(\pi_d)$ such that $[z,t] = (z,t)$ for $z\in \Link_d(M)$ and $t\in (0,1]$, and $[z,0] = \pi_d(z)$ for $z\in\Link_d(M)$ and $t=0$. Using these coordinates, we begin by constructing a fiberwise embedding of $\Link_d(M)$ into $\Link_d(\RR^{2n+1})$ — this is possible due to \Cref{ch4:: prop: fiberwise emb}, since 
    \begin{equation}
        \mathsf{dim}\big(\Link_{d}(M)\big) = n - 1, \quad \text{and} \quad \mathsf{dim}\big(\Link_{d}(\RR^{2n+1})\big) = 2n.
    \end{equation} and $2(n-1)+2 \leq 2n$.
    Let us denote this fiberwise embedding by $\widetilde{\varphi_d}$. We can now extend this embedding as follows:
    \begin{equation}
        \overline{\varphi}_d([z,t]) = 
        \begin{cases}
            (\widetilde{\varphi_d}(z),t) & \text{ for } t>0 \\
            \varphi_d(\pi_d(z)) & \text{ for } t = 0.
        \end{cases}
    \end{equation} 
    This map is conically smooth, since it is the dashed arrow in the following commutative diagram in $\Strat$:
    \begin{equation}
    \begin{tikzcd}[ampersand replacement=\&]
        \& {\Link_d(\RR^{2n+1})} \&\& {\Link_d(M)\times [0,1]} \\
        {\Link_d(M)} \& {} \& {\Link_d(M)\times [0,1]} \\
        \& {\varphi_d(M_d)} \&\& {\overline{\sC}(\widetilde{\pi}_d)} \\
        {M_d} \&\& {\overline{\sC}(\pi_d)}
        \arrow[""{name=0, anchor=center, inner sep=0}, "0", hook, from=1-2, to=1-4]
        \arrow[two heads, from=1-2, to=3-2]
        \arrow[from=1-4, to=3-4]
        \arrow["{\widetilde{\varphi_d}}"{pos=0.2}, from=2-1, to=1-2]
        \arrow["0"{pos=0.7}, hook, from=2-1, to=2-3]
        \arrow[two heads, from=2-1, to=4-1]
        \arrow["{\widetilde{\varphi_d}\times \id}"'{pos=0.6}, from=2-3, to=1-4]
        \arrow[from=2-3, to=4-3]
        \arrow[from=3-2, to=3-4]
        \arrow["{\varphi_d}", from=4-1, to=3-2]
        \arrow[from=4-1, to=4-3]
        \arrow["\lrcorner"{anchor=center, pos=0.125, rotate=180}, draw=none, from=4-3, to=2-1]
        \arrow[dashed, from=4-3, to=3-4]
        \arrow["\lrcorner"{anchor=center, pos=0.125, rotate=180}, draw=none, from=3-4, to=0]
    \end{tikzcd}
    \end{equation}
    Thus far, we have embedded $\overline{\sC}(\pi_d)$ into $\RR^{2n+1}$. To extend this to the entire $(M\longra P)$, we work on the top stratum $M_n$, which is a smooth $n$-manifold. The restriction of $\overline{\varphi}_d$ to the closed subset $\overline{\sC}(\pi_d)\cap M_n \subset M_n$ is a smooth embedding into $\RR^{2n+1}$, so by \Cref{ch4:: thm: whitney extension theorem} it extends to a smooth embedding of all of $M_n$ into $\RR^{2n+1}$. Because this extension agrees with $\overline{\varphi}_d$ on the overlap $\overline{\sC}(\pi_d)\cap M_n$, the two pieces glue to a globally defined conically smooth embedding $M \hookra \RR^{2n+1}$. Since $M$ is compact, this embedding is automatically proper.     
\end{proof}



A manifold with boundary can be embedded into another manifold with boundary in rather pathological ways, and the same can be said about conically smooth stratified spaces. The next definition fixes pathological behaviour by requiring additional regularity in the neighborhood of the boundary. An embedding is \emph{neat} if it admits a product-like structure near the boundary. 


\begin{proposition}
    \label{ch4:: prop: neat embedding strat}
    Let $\big(M \ra P\big)$ be a compact conically smooth stratified space with boundary $\big(M \ra P \fib \cP(*)\big)$. There exists a neat embedding 
    \begin{equation}
        \big(M \ra P\big) \hookra \RR_{\geq 0} \times \RR^{m}
    \end{equation}
    for some $m \geq 0$ --- where $\big(\RR_{\geq 0} \times \RR^{m}\big)$ admits the usual boundary structure.
\end{proposition}
\begin{proof}
    From \Cref{ch4:: prop: weak Whit strat}, we have a proper conically smooth embedding 
    \begin{equation}
        e \maps M \hookra \RR^{m}.
    \end{equation}
    Using an inward-pointing vector field and its flow, we obtain a collar
    \begin{equation}
        c\maps \partial M\times [0,\epsilon) \hookra M
    \end{equation}
    together with the associated boundary-defining function $\rho \maps M \longra \RR_{\geq 0}$. By the standard collar-straightening argument, after modifying $e$ by a conically smooth isotopy supported in the image of $c$, we may assume that on $c(\partial M\times [0,\epsilon/2))$ the embedding $e$ is independent of the collar coordinate.

    Taking the product of $e$ and $\rho$ gives a conically smooth map 
    \begin{equation}
        M \xra{\;\; \rho \times e \;\;} \RR_{\geq 0} \times \RR^{m}
    \end{equation}
    which is proper and injective. On the collar neighborhood $c(\partial M\times [0,\epsilon/2))$, it has the form
    \begin{equation}
        (p,t)\longmapsto \big(t,e_{|\partial M}(p)\big),
    \end{equation}
    so it is product-like near the boundary. This is exactly the neatness condition.
\end{proof}


\section{Constructible tangent bundle and structures}

For a conically smooth stratified space $M$, the constructible tangent bundle is a functor
\begin{equation}
    \sT_M \maps \exit(M)\longra \Vect^{\inj}.
\end{equation}
Concretely, it sends a point $x\in M_p$ to the ordinary tangent space $\sT_xM_p$ of the stratum containing it, and sends an exit-path to an injective linear map. The target $\Vect^{\inj}$ reflects the basic geometric feature that, as one exits a deeper stratum into a more generic one, tangent directions may be gained but not lost.

More precisely, $\sT_M$ is obtained as the composite
\begin{equation}
    \exit(M) \xra{\;\Delta_M\;} \Exit \xra{\;\sT^{\sf fib}\;} \Vect^{\inj}.
\end{equation}
The first functor $\Delta_M$ classifies the trivial constructible bundle
\begin{equation}
    M\times M \xra{\;\mathsf{pr}_1\;} M
\end{equation}
together with the diagonal section $\mathsf{diag}\maps M \longra M\times M$. The second functor $\sT^{\sf fib}$ should be thought of as the universal operation that linearizes such pointed constructible bundles. We now briefly introduce the $(\oo,1)$-categories $\Bun$, $\Exit$, and $\Vect^{\inj}$ that make this construction possible.

\begin{intuition}
One of the central constructions of Ayala--Francis--Rozenblyum (see \cite{AFR2}) is the $(\oo,1)$-category $\Bun$. Its basic universal property is that functors 
\begin{equation}
    \exit(K) \longra \Bun
\end{equation}
classify constructible bundles over $K$. The $(\oo,1)$-category $\Exit$ refines $\Bun$ by remembering a section: functors
\begin{equation}
    \exit(K) \longra \Exit
\end{equation}
classify constructible bundles over $K$ equipped with a section. The name of this category — the \emph{absolute exit-path $(\oo,1)$-category} — comes from the following property, recorded in \cite[Proposition~6.4.3]{AFR1}: for every constructible bundle $f\maps X \ra K$, the diagram
\begin{equation}
        \begin{tikzcd}[ampersand replacement=\&]
            \exit(X) \&\& \Exit \\
            \\
            \exit(K) \&\& {\Bun}
            \arrow["{(X \xra{f} K)}", hook, from=1-1, to=1-3]
            \arrow["\exit(f)"', two heads, from=1-1, to=3-1]
            \arrow[two heads, from=1-3, to=3-3]
            \arrow["(X \xra{f} K)", from=3-1, to=3-3]
        \end{tikzcd}
\end{equation} 
is a homotopy pullback diagram, indicating that the fibers over $\{f\maps X \ra K\}$ are given by $\exit(X)$. 
\end{intuition}

In particular, the projection $\mathsf{pr}_1\maps M\times M \longra M$ together with the diagonal section $\mathsf{diag}\maps M \longra M\times M$ determines a canonical functor
\begin{equation}
    \Delta_M \maps \exit(M)\longra \Exit.
\end{equation}

We now move on to the $(\oo,1)$-category $\Vect^{\sf inj}$.

\begin{intuition}
The $(\oo,1)$-category $\Vect^{\inj}$ is a generalization of the $\oo$-groupoid 
\begin{equation}
    \coprod_{i \geq 0} \BO(i).
\end{equation}
If $M$ is a smooth manifold, interpreted also as an $\oo$-groupoid, then functors
\begin{equation}
    M \longra \coprod_{i} \BO(i)
\end{equation}
classify vector bundles over $M$. The same role is played by $\Vect^{\inj}$ but for \emph{constructible} vector bundles \footnote{For simplicity, we avoid defining constructible vector bundles yet the careful reader might find it in \cite[Definition~2.2]{AFR2}.} over a conically smooth stratified space $M$, if we replace $M$ by its exit-path $(\oo,1)$-category $\exit(M)$. In practice, due to the $\inj$ superscript, we retain only constructible vector bundles whose transition maps are injective. Equivalently, along an exit-path the fiber dimension can stay the same or increase, but it cannot decrease; this is what lower semi-continuity means in the present context. This $(\oo,1)$-category is instantiated as a complete Segal space by defining it as follows:
\begin{equation}
    [p] \longmapsto \bigg|\bigg\{V_0 \hookra V_1 \hookra \cdots \hookra V_p \text{ over } \bDelta^{\bullet}\bigg\}\bigg|
\end{equation} 
where the right-hand-side is the space of injective composable sequences of finite rank vector bundles over $\bDelta^\bullet$, and such that for each $0\leq i \leq j \leq p$, the composition
\begin{equation}
    f_{j} \circ \cdots \circ f_{i}
\end{equation}
has locally constant rank. 
\end{intuition}

The remaining ingredient is the universal functor which extracts fiberwise linear data from pointed constructible bundles.
\begin{theorem}\label{ch4:: def: cbl tgt bundle}\cite[Proposition~2.15]{AFR2}
    There exists a unique symmetric monoidal functor under $\Vect^{\sim}$,
    \begin{equation}
        \sT^{\sf fib}\maps \Exit \longra \Vect^{\inj}
    \end{equation}
    such that:
    \begin{itemize}
        \item \label{closed-morphism} Sends closed morphisms $(\Strat)^{\cls} \hookra \Exit$ to equivalences.
        \item \label{embeddings} Sends embeddings $(\Strat)^{\emb} \hookra \Exit$ to equivalences.
    \end{itemize}
\end{theorem}

These two requirements encode the basic functorialities one expects from any object deserving the name tangent bundle. In the smooth case, if \(f\colon M \hookrightarrow N\) is an open embedding, then the derivative $\sD f$ identifies \(\sT M\) with the pullback of \(\sT N\), so tangent data is local for embeddings. In any suitably differential stratified setting, one further expects a restriction phenomenon along closed subspaces and closed strata; morally, passing to such a piece should not change the intrinsic linear data there. Requiring \(\sT^{\sf fib}\) to send these morphisms to equivalences is precisely the categorical way of encoding those expectations.

\begin{definition}
    The \define{constructible tangent bundle} of a conically smooth stratified space $M$ is the composite
    \begin{equation}
        \sT_M \maps \exit(M) \xra{\;\Delta_M\;} \Exit \xra{\;\sT^{\sf fib}\;} \Vect^{\inj}.
    \end{equation}
\end{definition}

This formal definition matches the concrete description given at the start of the section: $\sT_M$ sends a point $x\in M_p$ to the ordinary tangent space $\sT_xM_p$, and an exit-path to the corresponding injective linear map. A more conceptual and detailed discussion of how this composite is evaluated, both on objects and on exit-paths, can be found in \cite[Section~2]{AFR2}.


Having set up the constructible tangent bundle, we are naturally brought to the notion of conically smooth stratified \emph{tangential structure}.

\begin{definition}\label{ch4:: def: tangential structure}\cite[Definition~2.17]{AFR2}
    A \define{tangential structure} is an $(\oo,1)$-category $\cB$ equipped with a functor $\cB \longra \Vect^{\sf inj}$, and the $(\oo,1)$-category of tangential structures is the overcategory
    \begin{equation}
        (\Cat_{\oo})_{/\Vect^{\inj}}.
    \end{equation}
\end{definition}

This leads to the next definition.

\begin{definition}\label{ch4:: def: tangential structure 2}\cite[Definition~2.19]{AFR2}
    Fix a tangential structure $\cB \longra \Vect^{\inj}$ and a conically smooth stratified space $M$. A $\cB$\define{-framing on} $M$ is a lift $\varphi$
    \begin{equation}
    \begin{tikzcd}[ampersand replacement=\&]
        \& \cB \\
        {\exit(M)} \& {\Vect^{\inj}}.
        \arrow[from=1-2, to=2-2]
        \arrow["\varphi", dashed, from=2-1, to=1-2]
        \arrow["{\sT_M}"', from=2-1, to=2-2]
    \end{tikzcd}
    \end{equation} 
\end{definition}

\begin{remark}
    Note that, if $M$ is a smooth manifold then every functor $M \longra \cB$ factors through $\cB^{\sim} \hookra \cB$, the maximal $\oo$-subgroupoid. In particular this means that a $\cB$-framing for $M$ is a direct generalization of a $B$-framing, as discussed in Chapter 3.
\end{remark}


\begin{example}
    For each $n \geq 0$, the full subcategory $\Vect^{\sf inj}_{\leq n} \hookra \Vect^{\sf inj}$ on vector spaces of dimension at most $n$ is manifestly a tangential structure. A $\Vect^{\sf inj}_{\leq n}$-framing on a conically smooth stratified space $M$ is unique if it exists, and it exists precisely when all strata of $M$ have dimension at most $n$. This will be useful to us when discussing certain categories of cobordisms.
\end{example}

In this thesis we will focus exclusively on one especially tractable class of stratified tangential structures — \emph{solid tangential structures}.

\begin{definition}\cite[Definition~2.35]{AFR2}
    A \define{solid tangential structure} is a right fibration
    \begin{equation}
        \cB \longra \Vect^{\sf inj}.
    \end{equation}
\end{definition}

By straightening, to give a solid tangential structure is equivalent to stipulating the data of a functor
\begin{equation}
    F \maps (\Vect^{\sf inj})^{\op} \longra \Spaces.
\end{equation}
Thus, for every finite-dimensional vector space $V$, one has a space $F(V)$ of allowed structures on $V$, and for every injective linear map $V \hookra W$, one has a restriction map
\begin{equation}
    F(W) \longra F(V).
\end{equation}

In particular, a solid framing on a conically smooth stratified space $M$ is a lift of the constructible tangent bundle
\begin{equation}
    \sT_M \maps \exit(M) \longra \Vect^{\sf inj},
\end{equation}
and hence amounts to a compatible choice of such structure on each tangent space, varying contravariantly along exit-paths.

The point of this definition is that solid structures only remember how tangential data restricts along injections; they do not impose the stronger coherence data encoded by more general tangential structures, such as vari-framings. In this sense they are less sensitive to the underlying stratification. Indeed, it is proven by Ayala--Francis--Rozenblyum in \cite[Proposition~2.61]{AFR2} that solid tangential structures are invariant under refinement of the stratification.

We are especially interested in a subclass of solid framings, namely the \emph{solid $B$-framings}, for $B$ a space.

\begin{definition}
    Let $B$ be a space equipped with a map $B \longra \Vect^{\inj}$. A \define{solid $B$-framing} on a conically smooth stratified space $M$ is a lift
    \begin{equation}
    \begin{tikzcd}[ampersand replacement=\&]
        \& B \\
        {\exit(M)} \& {\Vect^{\inj}}.
        \arrow[from=1-2, to=2-2]
        \arrow["\varphi", dashed, from=2-1, to=1-2]
        \arrow["{\sT_M}"', from=2-1, to=2-2]
    \end{tikzcd}
    \end{equation}
\end{definition}

\begin{remark}
    Since $B$ is a space, a solid $B$-framing remembers only a homotopy-coherent choice of $B$-structure on each constructible tangent space, compatible under restriction along exit-paths. In particular, it is even less sensitive to the incidence data of the strata than a general tangential structure, and by the result quoted above it is unchanged under refinement of the stratification.
\end{remark}

\begin{example}
    The simplest case is $B=*$. Choosing the object $\RR^n\in \Vect^{\inj}$ determines a map $* \longra \Vect^{\inj}$, and for a smooth $n$-manifold $M$ a solid $*$-framing is exactly an ordinary framing $\sT M \cong M\times \RR^n$.
\end{example}

To treat families of conically smooth stratified framed manifolds, we need the fiberwise version of the constructible tangent bundle. Let $\pi\maps M \longra K$ be a conically smooth fiber bundle. The pointed constructible bundle
\begin{equation}
    M\times_{K} M \xra{\;\mathsf{pr}_1\;} M
\end{equation}
together with the diagonal section $\mathsf{diag}\maps M \longra M\times_{K} M$ determines a canonical functor
\begin{equation}
    \Delta_{\pi}\maps \exit(M)\longra \Exit.
\end{equation}

\begin{definition}
    The \define{fiberwise constructible tangent bundle} of $\pi$ is the composite
    \begin{equation}
        \sT^{\mathrm{fib}}_{\pi}\maps \exit(M) \xra{\;\Delta_{\pi}\;} \Exit \xra{\;\sT^{\sf fib}\;} \Vect^{\inj}.
    \end{equation}
\end{definition}

This is the stratified analogue of the ordinary vertical tangent bundle of a smooth fiber bundle. Indeed, if $\pi\maps M \longra K$ is a smooth fiber bundle, then the usual vertical tangent bundle is the kernel of $\sD\pi$,
\begin{equation}
    \sT^{\mathrm{vert}}_{\pi}M\coloneqq \ker(\sD\pi)\longra M,
\end{equation}
and $\sT^{\mathrm{fib}}_{\pi}$ plays exactly this role in the constructible setting. In the special case $K=*$, the fiber product $M\times_{*}M$ is just $M\times M$, so $\sT^{\mathrm{fib}}_{\pi}$ reduces to the constructible tangent bundle $\sT_M$.

\begin{definition}
    Let $\cB \longra \Vect^{\inj}$ be a tangential structure and let $\pi\maps M \longra K$ be a conically smooth fiber bundle. A \define{vertical $\cB$-framing} on $\pi$ is a lift
    \begin{equation}
    \begin{tikzcd}[ampersand replacement=\&]
        \& \cB \\
        {\exit(M)} \& {\Vect^{\inj}}.
        \arrow[from=1-2, to=2-2]
        \arrow["\varphi", dashed, from=2-1, to=1-2]
        \arrow["{\sT^{\mathrm{fib}}_{\pi}}"', from=2-1, to=2-2]
    \end{tikzcd}
    \end{equation}
\end{definition}

If $B$ is a space and $\pi\maps X \longra *$ is the terminal map, then $\sT^{\mathrm{fib}}_{\pi}=\sT_X$, so a vertical solid $B$-framing on $\pi$ is exactly the usual solid $B$-framing on $X$; when $X$ is a smooth manifold, this is the ordinary $B$-framing of Chapter 3. This is the notion we will need later when discussing families of conically smooth stratified framed manifolds.

